\documentclass[a4paper,11pt]{article}
% %%%%%%%%%%%%%%%%%%%%%%%%%%%%%%%%%%%%%%%%%%%%%%%%%%%%%%%%%%%%%%%%%%%%%
\usepackage[a4paper, total={6.5in, 9.8in}]{geometry}
% %%%%%%%%%%%%%%%%%%%%%%%%%%%%%%%%%%%%%%%%%%%%%%%%%%%%%%%%%%%%%%%%%%%%%
\usepackage{amssymb}
\usepackage{amsfonts}
\usepackage{amsmath}
\usepackage{makeidx}
\usepackage{fontspec}
\usepackage{xunicode}
\usepackage{xltxtra}
\usepackage{listings}
\usepackage{xcolor}
\usepackage{graphicx}
\usepackage{morefloats}
\usepackage[yyyymmdd]{datetime}
\usepackage{appendix}
\usepackage{tikz}
\usepackage{enumitem}
\usepackage[most]{tcolorbox}
\usepackage{hyperref}
\hypersetup{
  pdftitle={Array Module Documentation (array\textunderscore mod)},
  pdfauthor={Konstantinos N. Anagnostopoulos},
  colorlinks=true,
  linkcolor=blue!70!black,
  filecolor=magenta,
  citecolor=blue!70!black,
  anchorcolor=blue!70!black,
  urlcolor=blue!70!black,
  bookmarksopen=true,
  bookmarksdepth=3
}
\urlstyle{same}
\setlength{\marginparwidth}{2cm}
\emergencystretch=1.5em

\setmainfont[Mapping=tex-text]{GFS Didot}

% Load user's LaTeX macros and listings configuration from Documentation/array_mod_doc/mymacros.tex
%\newcommand{\part}{\partial}
%\renewcommand{\Huge}{\Large} % προβλημα με ελληνικες γραμματοσειρες
%\renewcommand{\huge}{\Large} % προβλημα με ελληνικες γραμματοσειρες
%%\renewcommand\contentsname{ΠΕΡΙΕΧΟΜΕΝΑ}
%%\renewcommand\listfigurename{ΚΑΤΑΛΟΓΟΣ ΤΩΝ ΕΙΚΟΝΩΝ}
%%\renewcommand\listtablename{ΚΑΤΑΛΟΓΟΣ ΠΙΝΑΚΩΝ}
%%%\newcommand\refname{ΑΝΑΦΟΡΕΣ}
%%\renewcommand\refname{ΑΝΑΦΟΡΕΣ}
%%\renewcommand\indexname{Index}
%%\renewcommand\figurename{Σχήμα}
%%\renewcommand\tablename{Πίνακας}
%%\renewcommand\partname{Μέρος}
%%\renewcommand\chaptername{ΚΕΦΑΛΑΙΟ}
%%\renewcommand\appendixname{ΠΑΡΑΡΤΗΜΑ}
%%%\newcommand\abstractname{ΠΕΡΙΛΗΨΗ}
%%\renewcommand\abstractname{ΠΕΡΙΛΗΨΗ}
%%\renewcommand\bibname{ΒΙΒΛΙΟΓΡΑΦΙΑ}
%%\def\today{\number\day \space\ifcase\month\or
%%  Ιανουαρίου \or Φεβρουαρίου\or Μαρτίου\or Απριλίου\or Μαΐου\or Ιουνίου\or
%%  Ιουλίου\or Αυγούστου\or Σεπτεμβρίου\or Οκτωβρίου\or Νοεμβρίου\or
%%  Δεκεμβρίου\fi 
%%   \space\number\year}

% %%%%%%%%%%%%%%%%%%%%%%%%%%%%%%%%%%%%%%%%%%%%%%%%%%%%%%%%%%
% %%%%%%%%%%%%%%%%%%%%%%%%%%%%%%%%%%%%%%%%%%%%%%%%%%%%%%%%%%
%                     math definitions
% %%%%%%%%%%%%%%%%%%%%%%%%%%%%%%%%%%%%%%%%%%%%%%%%%%%%%%%%%%
\newcommand  {\e}       [1]{{\rm e}^{#1}}
\newcommand  {\vev}     [1]{\ensuremath{\langle #1 \rangle}}
\newcommand  {\veveq}   [1]{\ensuremath{{\langle #1 \rangle}_{\mbox{\scriptsize eq}}}}
\newcommand  {\norm}    [1]{\ensuremath{\left\lVert #1 \right\rVert}}  % double vertical line norm like ||x||
\newcommand  {\rf}      [1]{(\ref{#1})}
\newcommand  {\eq}      [1]{Eq.(\protect\ref{#1})}
\newcommand  {\eqs}     [2]{Eqs.(\protect\ref{#1})--(\protect\ref{#2})}
\newcommand  {\fig}     [1]{Figure \protect\ref{#1}}
\newcommand  {\figs}    [2]{Figures \protect\ref{#1} -- \protect\ref{#2}}
\newcommand  {\tbl}     [1]{Table  \protect\ref{#1}}
\newcommand  {\code}    [1]{\colorbox{listinggray}{{\tt #1}}}                           % \definecolor{listinggray}{gray}{0.9}  -- defined below
\newcommand  {\lcode}   [1]{\colorbox{listinggray}{\lstinline[language=bash       ]{#1}}}   
\newcommand  {\lcodeb}  [1]{\colorbox{listinggray}{\lstinline[language=bash       ]{#1}}}   
\newcommand  {\lcodef}  [1]{\colorbox{listinggray}{\lstinline[language=Fortran    ]{#1}}}    
\newcommand  {\lcodec}  [1]{\colorbox{listinggray}{\lstinline[language=C          ]{#1}}}   
\newcommand  {\lcodea}  [1]{\colorbox{listinggray}{\lstinline[language=Awk        ]{#1}}}   
\newcommand  {\lcodeg}  [1]{\colorbox{listinggray}{\lstinline[language=Gnuplot    ]{#1}}}   
\newcommand  {\lcodem}  [1]{\colorbox{listinggray}{\lstinline[language=Mathematica]{#1}}}   
\newcommand  {\cbox}    [1]{\colorbox{listinggray}{{\tt #1}}}                           
\renewcommand{\implies}    {\ensuremath{\Rightarrow}}
\newcommand  {\nn}         {\nonumber\\}
\newcommand  {\ket}     [1]{\ensuremath{\,| #1 \rangle}}
\newcommand  {\bra}     [1]{\ensuremath{\langle #1 |\,}}
\newcommand  {\braket}  [2]{\ensuremath{\langle #1 | #2 \rangle}}
\newcommand  {\calO}       {\ensuremath{\mathcal{O}}}
\newcommand  {\calA}       {\ensuremath{\mathcal{A}}}
\newcommand  {\calM}       {\ensuremath{\mathcal{M}}}
\newcommand  {\calD}       {\ensuremath{\mathcal{D}}}
\newcommand  {\calN}       {\ensuremath{\mathcal{N}}}
\newcommand  {\calQ}       {\ensuremath{\mathcal{Q}}}
\newcommand  {\smbox}   [1]{\ensuremath{\mbox{\eng{\scriptsize{#1}}}}}
\newcommand  {\embox}   [1]{\ensuremath{\mbox{\eng{{#1}}}}}
\newcommand  {\padd}       {\ensuremath{P_{\smbox{add}}}}
\newcommand  {\abs}     [1]{\ensuremath{\left| #1 \right|}}
\newcommand  {\pd}      [2]{\ensuremath{\frac{\partial   #1}{\partial #2}}}
\newcommand  {\pdd}     [2]{\ensuremath{\frac{\partial^2 #1}{\partial {#2}^2}}}
\newcommand  {\tr}         {\ensuremath{{\rm tr\,}}}
\newcommand  {\Tr}         {\ensuremath{{\rm Tr\,}}}
\newcommand  {\Reals}      {\ensuremath{\mathbb{R}}}
\newcommand  {\Complexes}  {\ensuremath{\mathbb{C}}}
\newcommand  {\Naturals}   {\ensuremath{\mathbb{N}}}
\newcommand  {\Rationals}  {\ensuremath{\mathbb{Q}}}
\newcommand  {\Integers}   {\ensuremath{\mathbb{Z}}}
\newcommand  {\UnitMatrix} {\ensuremath{\mathbb{I}}}
\newcommand  {\el}         {\ensuremath{\in}}        % Mathematica notation
\renewcommand{\Re}         {\ensuremath{\mbox{Re}}}
\renewcommand{\Im}         {\ensuremath{\mbox{Im}}}
%\newcommand  {\Seff}       {\ensuremath{S_{\text{\scriptsize{eff}}}}}
%\newcommand  {\tSeff}      {\ensuremath{\tilde{S}_{\text{\scriptsize{eff}}}}}
\newcommand  {\Seff}       {\ensuremath{S_{\text{eff}}}}
\newcommand  {\tSeff}      {\ensuremath{\tilde{S}_{\text{eff}}}}
\newcommand  {\dtau}       {\ensuremath{d\tau}}
\newcommand  {\taut}       {\ensuremath{\tau_{\mbox{\scriptsize aut}}}}
\newcommand  {\tint}       {\ensuremath{\tau_{\mbox{\scriptsize int}}}}
\newcommand  {\MyPhi}      {\ensuremath{P}}                             % notation of pdf of stochastic variables
\newcommand  {\MyPhieq}    {\ensuremath{\MyPhi_{\mbox{\scriptsize eq}}}}
\newcommand  {\mscript} [2]{\ensuremath{{#1}_{\mbox{\scriptsize #2}}}}
\newcommand  {\mscr}    [2]{\ensuremath{{#1}_{\mbox{\scriptsize #2}}}}
\newcommand  {\Pf}         {\ensuremath{\text{Pf}\,}}
\newcommand  {\PfM}         {\ensuremath{\text{Pf}\,\mathcal{M}}}
\newcommand  {\detM}        {\ensuremath{\det\mathcal{M}}}
\newcommand  {\detD}        {\ensuremath{\det\mathcal{D}}}
\newcommand  {\mf}          {\ensuremath{m_{\text{f}}}}
\newcommand  {\ncg}         {\ensuremath{N_{\text{cg}}}}
\newcommand  {\NS}          {\ensuremath{\mathscr{D}}}  % symbol for spinor dimensionality 2^{D/2-1}. Choices: N_S,N_s, \mathscr{D}, \mathfrak{D}, d_s
% %%%%%%%%%%%%%%%%%%%%%%%%%%%%%%%%%%%%%%%%%%%%%%%%%%%%%%%%%%
% tofo notes: \usepackage[colorinlistoftodos,shadow]{todonotes}, \listoftodos, \todo[optios]{text}, \missingfigure[figwidth=0.5\textwidth]{text}, \todo[inline,size=\Large,color=red!40]{#1} (option disable: \usepackage[disable,...]{todonotes})
\usepackage[colorinlistoftodos,shadow]{todonotes}
\newcommand {\XXX}     [1]{\todo[size=\tiny,color=green!40]{#1}}
\newcommand {\XXXS}    [1]{\todo[size=\small,color=green!40]{#1}}
\newcommand {\XXXs}    [1]{\todo[size=\tiny ,color=green!40]{#1}}
\newcommand {\XXXi}    [1]{\todo[inline,color=green!40]{#1}}
\newcommand {\XXXX}    [1]{\todo[inline,color=green!40]{#1}}
\newcommand {\XXXtodo} [1]{\todo{#1}}
\newcommand {\XXXred}  [1]{\todo[size=\small,color=red!40]{#1}}
\newcommand {\XXXXXX}  [1]{\todo[inline,size=\Large,color=red!40]{#1}}
\newcommand {\XXXref}  [1]{\todo[size=\small,linecolor=green!70!white,backgroundcolor=blue!20!white,bordercolor=red]{Add Ref.: #1}}
\newcommand {\XXXfig}  [1]{\todo[inline,color=green!40]{#1}} %insert in captions
\newcommand {\XXXtab}  [1]{\todo[inline,color=green!40]{#1}}
\newcommand {\XXXaddfig}[1]{\missingfigure[figwidth=0.5\textwidth]{#1}}
% %%%%%%%%%%%%%%%%%%%%%%%%%%%%%%%%%%%%%%%%%%%%%%%%%%%%%%%%%%
% index alias:
\newcommand {\ii}      [1]{\index{#1}}
% %%%%%%%%%%%%%%%%%%%%%%%%%%%%%%%%%%%%%%%%%%%%%%%%%%%%%%%%%%
% Other packages:
\usepackage{booktabs} % For tables in chapter 06_Parallel
%\usepackage{afterpage} % For big floats placement. Use \afterpage{ \begin{figure}[!htb] ...  \end{figure} }
\usepackage{float}    % For big floats. use  \begin{figure}[H] to actually place it here!
\usepackage{mathrsfs} % \mathscr{D}
% %%%%%%%%%%%%%%%%%%%%%%%%%%%%%%%%%%%%%%%%%%%%%%%%%%%%%%%%%%
%                     hyphenation
% %%%%%%%%%%%%%%%%%%%%%%%%%%%%%%%%%%%%%%%%%%%%%%%%%%%%%%%%%%
%\hyphenation{every-where}
\hyphenation{Lan-ge-vin}
% %%%%%%%%%%%%%%%%%%%%%%%%%%%%%%%%%%%%%%%%%%%%%%%%%%%%%%%%%%
%                     include subsubsections
% %%%%%%%%%%%%%%%%%%%%%%%%%%%%%%%%%%%%%%%%%%%%%%%%%%%%%%%%%%
\setcounter{secnumdepth}{3}
\setcounter{tocdepth}{3}


% %%%%%%%%%%%%%%%%%%%%%%%%%%%%%%%%%%%%%%%%%%%%%%%%%%%%%%%%%
%   Code listing: 
% %%%%%%%%%%%%%%%%%%%%%%%%%%%%%%%%%%%%%%%%%%%%%%%%%%%%%%%%%
%\usepackage{listings}
% /usr/share/doc/texlive-doc/latex/listings/listings.pdf
% Languages: C, Fortran, [77]Fortran,  Gnuplot, Awk, csh, TeX,
% LaTeX,make, Perl, Lisp,HTML....
%\usepackage{color}
%%\lstset{language=Gnuplot,alsolanguage=csh,basicstyle=\scriptsize,deletekeywords={gawk,emacs},morekeywords={print,printf}}
%%\begin{lstlisting}[caption={Hello there},label=lbl]
%%\begin{lstlisting}[title={My title}]
%% ..................
%%\end{lstlisting}
\definecolor{listinggray}{gray}{0.9}
\definecolor{lbcolor}{rgb}{0.9,0.9,0.9}
\lstset{
        backgroundcolor=\color{lbcolor},
        tabsize=4,
        rulecolor=,
        language=Fortran,
        basicstyle=\footnotesize,%\scriptsize,
        upquote=true,
        aboveskip={1.5\baselineskip},
        columns=fixed,
        showstringspaces=false,
        extendedchars=true,
        inputencoding=utf8x,
        breaklines=true,
        prebreak = \raisebox{0ex}[0ex][0ex]{\ensuremath{\hookleftarrow}},
        frame=single,
        showtabs=false,
        showspaces=false,
        showstringspaces=false,
        identifierstyle=\ttfamily,
        keywordstyle=\color[rgb]{0,0,1},
        commentstyle=\color[rgb]{0.133,0.545,0.133},
        stringstyle=\color[rgb]{0.627,0.126,0.941},
}
% from /usr/share/texlive/texmf-dist/tex/latex/listings/lstlang1.sty
\lstdefinelanguage[95]{Fortran}[77]{Fortran}%
  {deletekeywords=SAVE,%
   morekeywords={subroutine,ACTION,ADVANCE,ALLOCATE,ALLOCATABLE,ASSIGNMENT,CASE,%
      CONTAINS,CYCLE,DEALLOCATE,DEFAULT,DELIM,EXIT,INCLUDE,IN,NONE,IN,%
      OUT,INTENT,INTERFACE,IOLENGTH,KIND,LEN,MODULE,NAME,NAMELIST,NMT,%
      NULLIFY,ONLY,OPERATOR,OPTIONAL,OUT,PAD,POINTER,POSITION,PRIVATE,%
      PUBLIC,READWRITE,RECURSIVE,RESULT,SELECT,SEQUENCE,SIZE,STAT,%
      TARGET,USE,WHERE,WHILE,BLOCKDATA,DOUBLEPRECISION,%
      ENDBLOCKDATA,ENDFILE,ENDFUNCTION,ENDINTERFACE,%
      ENDMODULE,ENDPROGRAM,ENDSELECT,ENDSUBROUTINE,ENDTYPE,ENDWHERE,%
      INOUT,SELECTCASE,FORALL,RESHAPE,SUM,PRODUCT,ELSEWHERE,%
      sin,cos,tan,atan2,exp,cexp,dexp,cdexp,log,abs,conjg,dconjg,%
      dble,dimag,cdabs,imag,tiny,huge,epsilon,range,trim,atan,pure,%
      ifdef,ifndef,define,endif,elif,%
      MATMUL,TRANSPOSE,UBOUND,LBOUND,SAVE,%
      MPI_LOGICAL,MPI_INTEGER,MPI_INTEGER8,MPI_REAL,MPI_REAL8,MPI_COMPLEX,MPI_COMPEX16,MPI_CHARACTER,MPI_STATUS_IGNORE,%
      MPI_SUM,MPI_PROD,MPI_MAX,MPI_MIN,%
      concurrent,CONCURRENT,random_number,%
      co_broadcast,co_sum,co_max,com_min,co_reduce,sync,all,this_image,image_index,num_images,error,lcobound,ucobound,codimension,%
      mpi_init,mpi_comm_size,mpi_comm_rank,mpi_finalize,mpi_gather,mpi_allgather,mpi_send,mpi_recv,%
      mpi_reduce,mpi_allreduce,MPI_COMM_WORLD,mpi_bcast,mpi_wtime,mpi_get_count,mpi_scatter,mpi_barrier,mpi_sendrecv,MPI_TAG,MPI_SOURCE,%
      mpi_comm_group,mpi_group_incl,mpi_comm_create,mpi_comm_split,mpi_comm_null,mpi_comm_free,mpi_sendrecv_replace,%
      blacs_get,blacs_gridinit,blacs_gridinfo,blacs_exit,blacs_pinfo,blacs_barrier,blacs_exit,blacs_pnum,blacs_gridmap,%
      pdelget,pdelset,pdtran,dgemm,pdgemm,numroc,descinit,zheev,pzheev,dsyev,pdsyev,%
      mpiu_init,mpiu_end,mpiu_sync,mpiu_abort,mpiu_world_set,mpiu_bcast,mpiu_scatter,mpiu_gather,mpiu_sndrcv,mpiu_sum,mpiu_prod,%
      indxl2g,indxg2l,indxg2p},%
   deletecomment=[f],% no fixed comment line: 1998 Magne Rudshaug
   morecomment=[l]!%
  }%
%% csh definition (c) 1998 Kai Below
%%
\lstdefinelanguage{csh}
  {morekeywords={alias,awk,cat,echo,else,end,endif,endsw,exec,exit,%
      foreach,glob,goto,history,if,logout,nice,nohup,onintr,repeat,sed,%
      set,setenv,shift,source,switch,then,time,while,umask,unalias,%
      unset,wait,while,@,env,argv,child,home,ignoreeof,noclobber,%
      noglob,nomatch,path,prompt,shell,status,verbose,print,printf,%
      sqrt,BEGIN,END,gfortran,ls,rm,less,chmod,perl,cd,pwd,mkdir,rmdir,%
      sed,cp,mv,which,where,whatis,cmd,man,info,locate,head,tail,sort,grep,%
      gawk,emacs,du,df,top,ps,sudo,apt-get},%
   morecomment=[l]\#,%
   morestring=[d]"%
  }[keywords,comments,strings]%

% %%%%%%%%%%%%%%%%%%%%%%%%%%%%%%%%%%%%%%%%%%%%%%%%%%%%%%%%%
%   Making emacs refcard table in 01computer.tex
% %%%%%%%%%%%%%%%%%%%%%%%%%%%%%%%%%%%%%%%%%%%%%%%%%%%%%%%%%
% We assume 3 columns in the table. Change multicolumn commands and &
% appropriately if changed
\usepackage{longtable}
\newcommand{\emtentry}   [1]{      \multicolumn{3}{l}{\parbox[t]{10cm}{\scriptsize{\eng{#1}}}}\\}
\newcommand{\emtsection} [1]{\hline\multicolumn{3}{l}{\eng{\bfseries{{#1}}}}\\\hline}
\newcommand{\emtkey}     [2]{\eng{\scriptsize{#1}} & \engt{\scriptsize{#2}} & \\}
\newcommand{\kbd}        [1]{{\tt\scriptsize {#1}}\null}
\newcommand{\emtthreecol}[3]{\eng{\scriptsize{#1}}&\engt{\scriptsize{#2}}&\engt{\scriptsize{#3}}\\}
%\newcommand{\emtkey}[2]{\eng{#1} & \multicolumn{2}{l}{\engt{#2}}\\}
% refer to problems: askhseis:
\newcommand{\pref}[1]{\thechapter.\ref{#1}}

\newenvironment{latintext}{}{}
%\newcommand{\latintext}{}
\newcommand{\greektext}{}

\usepackage{rotating}
% rotate tables:
%%\begin{sidewaystable}
%%    \centering
%%    \caption{Your caption here}
%%   \begin{tabular}{ll}
%%    First First & First Second\\
%%    Second First & Second Second
%%    \end{tabular}
%%\end{sidewaystable}


% Load custom documentation box and man-page styling macros
% ===================================================================
% Custom Macros and Styling for array_mod Extended Documentation
% ===================================================================

% Color Palette
\definecolor{modPrimary}{RGB}{22, 64, 122}      % Deep Navy Blue
\definecolor{modSecondary}{RGB}{31, 111, 139}   % Teal / Slate Blue
\definecolor{modAccent}{RGB}{168, 42, 42}       % Crimson / Burgundy for Procedure Names
\definecolor{modArg}{RGB}{0, 102, 51}           % Forest Green for Dummy Arguments
\definecolor{modType}{RGB}{102, 51, 153}        % Royal Purple for Types / Intents
\definecolor{modBoxBg}{RGB}{244, 247, 251}      % Very Light Ice Blue background
\definecolor{modBoxFrame}{RGB}{44, 88, 150}     % Medium Navy frame
\definecolor{modWarnBg}{RGB}{254, 249, 238}     % Warm Amber background
\definecolor{modWarnFrame}{RGB}{204, 122, 0}    % Amber frame
\definecolor{modCodeBg}{RGB}{242, 244, 247}     % Clean light slate gray for listings
\definecolor{modSynBg}{RGB}{236, 241, 248}      % Soft blue-gray for Synopsis blocks

% Enhance default listings appearance while preserving mymacros.tex settings
\lstset{
    backgroundcolor=\color{modCodeBg},
    rulecolor=\color{modBoxFrame!50},
    frame=single,
    basicstyle=\footnotesize\ttfamily,
    keywordstyle=\color[rgb]{0,0,0.75}\bfseries,
    commentstyle=\color[rgb]{0.13,0.50,0.13}\itshape,
    stringstyle=\color[rgb]{0.58,0.10,0.75},
    upquote=true,
    breaklines=true,
    columns=fixed,
    showstringspaces=false
}

% Typography Macros for Procedures, Arguments, Types, and Files
\newcommand{\procname}[1]{\textcolor{modAccent}{\texttt{\textbf{#1}}}}
\newcommand{\argname}[1]{\textcolor{modArg}{\texttt{\textbf{#1}}}}
\newcommand{\typename}[1]{\textcolor{modType}{\texttt{#1}}}
\newcommand{\intentname}[1]{\textcolor{modSecondary}{\texttt{#1}}}
\newcommand{\modname}[1]{\textcolor{modPrimary}{\texttt{\textbf{#1}}}}
\newcommand{\filename}[1]{\textcolor{modSecondary}{\texttt{#1}}}

% Man-Page Section Label Macro
\newcommand{\manheading}[1]{%
  \par\vspace{0.35em}\noindent%
  {\sffamily\bfseries\footnotesize\textcolor{modPrimary}{\MakeUppercase{#1}}}\par\vspace{0.15em}\noindent%
}

% Colored Informational & Warning Boxes
\newtcolorbox{notebox}[1][Note]{
  enhanced,
  breakable,
  colback=modBoxBg,
  colframe=modBoxFrame,
  coltitle=white,
  fonttitle=\sffamily\bfseries\small,
  title={#1},
  arc=1.5mm,
  boxrule=0.6pt,
  left=8pt, right=8pt, top=6pt, bottom=6pt
}

\newtcolorbox{warnbox}[1][Important]{
  enhanced,
  breakable,
  colback=modWarnBg,
  colframe=modWarnFrame,
  coltitle=white,
  fonttitle=\sffamily\bfseries\small,
  title={#1},
  arc=1.5mm,
  boxrule=0.6pt,
  left=8pt, right=8pt, top=6pt, bottom=6pt
}

% Man-Page Procedure Entry Header Box & TOC Registration
% Usage: \manentry{kind}{proc_display_name}{toc_name}{brief_summary}{label_id}
\newcommand{\manentry}[5]{%
  \par\vspace{1.2em}%
  \phantomsection%
  \label{#5}%
  \addcontentsline{toc}{subsubsection}{\texorpdfstring{\texttt{\textbf{#3}} \textnormal{\footnotesize--- #4}}{#3}}%
  \begin{tcolorbox}[
    enhanced,
    colback=modBoxBg,
    colframe=modPrimary,
    boxrule=0.8pt,
    leftrule=4.5pt,
    arc=1.5mm,
    left=8pt, right=8pt, top=6pt, bottom=6pt,
    before skip=0.5em,
    after skip=0.4em
  ]
  \noindent{\sffamily\bfseries\scriptsize\textcolor{modSecondary}{\MakeUppercase{#1}}}\hfill{\sffamily\scriptsize\textcolor{gray}{array\_mod(3f)}}\\[2pt]
  {\large\procname{#2}} \quad {\normalsize --- \textbf{#4}}
  \end{tcolorbox}%
  \noindent\textbf{\sffamily\footnotesize\textcolor{modPrimary}{NAME}}\par\vspace{0.1em}%
  \noindent\hspace{1.2em}\procname{#2} --- #4\par%
}

% Argument Item Macro inside Dummy Arguments list
% Usage: \argitem{arg_name}{type_and_intent_spec}{description}
\newcommand{\argitem}[3]{%
  \item[\argname{#1}] (\typename{#2}): #3%
}


\setcounter{secnumdepth}{3}
\setcounter{tocdepth}{3}

\title{%
  \vspace{-1.2cm}%
  \begin{tcolorbox}[
    enhanced,
    colback=modPrimary,
    colframe=modPrimary!80!black,
    coltext=white,
    arc=2mm,
    boxrule=1pt,
    top=12pt, bottom=12pt, left=12pt, right=12pt
  ]
  \centering
  {\LARGE\bfseries Array Module Documentation}\\[6pt]
  {\large Extended Reference Manual for the Fortran \texttt{array\_mod}, \texttt{tensorprod\_mod}, and \texttt{timer\_mod} Modules}
  \end{tcolorbox}%
  \vspace{-0.4cm}%
}
\author{\textbf{Konstantinos N. Anagnostopoulos}\\
  {\small Physics Department, National Technical University of Athens}\\
  {\small \texttt{konstant@mail.ntua.gr} $\cdot$ \url{https://www.physics.ntua.gr/konstant}}}
\date{\small \today}

\begin{document}

\maketitle

\begin{abstract}
\noindent
This document provides the complete reference manual for the public procedures and interfaces of the Fortran \modname{array\_mod} module (and its companion modules \modname{tensorprod\_mod} and \modname{timer\_mod}) in the \texttt{Matrix} library (\filename{Software/Tools/Matrix/}). Built on top of \texttt{BLAS}, \texttt{LAPACK}, and \texttt{PFAPACK}, \modname{array\_mod} offers high-performance, ergonomic Fortran interfaces for double-precision real (\typename{real(8)}) and complex (\typename{complex(8)}) linear algebra, including traces, matrix symmetries, sorting, Fisher--Yates shuffling, matrix--matrix and matrix--vector multiplication, matrix inversion, determinants, Pfaffians, eigenvalues and right eigenvectors, Kronecker tensor products, Pauli and Dirac $\Gamma$ matrices, uniform and Gaussian random number generation, formatted I/O, Mathematica export, IEEE \texttt{NaN} handling, and execution timing. Every public procedure is documented in Unix man-page style with its name, synopsis, brief summary, detailed mathematical and algorithmic description, dummy arguments, and usage examples.
\end{abstract}

\vspace{0.3cm}
\begin{tcolorbox}[
  enhanced,
  breakable,
  colback=modBoxBg!60,
  colframe=modPrimary!60,
  arc=1.5mm,
  boxrule=0.6pt,
  left=10pt, right=10pt, top=6pt, bottom=8pt
]
\tableofcontents
\end{tcolorbox}
\newpage

% Section 1: Usage, Compiling and Linking, Download, Hello World!
% ===================================================================
% Section 1: Introduction and Usage
% ===================================================================
\section{Usage}
\label{sec:usage}

The \modname{array\_mod} module (part of the \texttt{Matrix} library) provides high-performance, idiomatic Fortran interfaces for linear algebra calculations and array manipulations on double-precision real (\typename{real(8)} / \typename{real(real64)}) and double-precision complex (\typename{complex(8)} / \typename{complex(real64)}) arrays.

Designed for scientific computing and matrix model simulations, \modname{array\_mod} wraps performance-critical Fortran~77 \texttt{BLAS} and \texttt{LAPACK} routines, as well as the \texttt{PFAPACK} library, behind clean modern Fortran generic interfaces. Key capabilities include:
\begin{itemize}
  \item \textbf{Basic Matrix Operations}: Matrix trace, quadratic and connected traces ($\Tr A$, $\Tr A^2$, $\Tr(AB)$), traceless projection, Hermitian/symmetric/antisymmetric transformations and norm checks, diagonal extraction and construction, identity and Pauli matrices ($\sigma_0,\sigma_1,\sigma_2,\sigma_3$), $L_1$-averaged array norms, Quicksort with permutation tracking, and Fisher--Yates random shuffling.
  \item \textbf{Matrix Multiplication}: High-performance matrix--matrix and matrix--vector products via overloaded operator \procname{.mm.}, functional interface \procname{lmatmul} (supporting general, Hermitian, symmetric, and rank-$k$ products $A A^\dagger, A^\dagger A, A A^T, A^T A$), and zero-overhead subroutine interfaces (\procname{mmmult}, \procname{mvmult}, \procname{vmmult}).
  \item \textbf{Linear Algebra}: Matrix inversion (\procname{inverse}), determinants (\procname{determinant}) and overflow-safe logarithmic determinants (\procname{lndet}), Pfaffians (\procname{pfaffian}) and logarithmic Pfaffians (\procname{lnPfaffian}) for complex skew-symmetric matrices, and full spectral decompositions (\procname{eigenvalues} and \procname{eigenvectors}) with automatic sorting options.
  \item \textbf{Tensor Products}: Kronecker products (\procname{tensorprod}) of up to six matrices or vectors, along with vector reshaping (\procname{vec2col}, \procname{vec2row}) and column-major matrix vectorization (\procname{vectorize}), exported from \modname{tensorprod\_mod}.
  \item \textbf{Random Numbers}: Overloaded \procname{random\_number} for uniform and Box--Muller Gaussian distributions on real and complex scalars and arrays of rank 1, 2, and 3, paired with entropy seeding from \filename{/dev/urandom} (\procname{matrix\_random\_init}).
  \item \textbf{Utilities \& Timing}: Formatted I/O (\procname{save}, \procname{read}), rich array inspection and Mathematica/Gamma-matrix export (\procname{print}, \procname{printna}), IEEE quiet \procname{NaN} generation and testing (\procname{isNaN}), and CPU/wall-clock execution profiling exported from \modname{timer\_mod} (\procname{timer\_mod\_init}, \procname{timer\_mod\_now}, \procname{timer\_mod\_print}).
\end{itemize}

A minimal program illustrating the use of \modname{array\_mod} is shown below:
\lstset{language=[95]Fortran}
\begin{lstlisting}
program hello_world
  use :: array_mod
  complex(8), allocatable :: m(:,:)
  complex(8)              :: z

  n = 8
  allocate(m(n,n))

  call matrix_random_init           ! seed random number from /dev/urandom
  call random_number(m, sigma=1._8) ! fill with Gaussian complex random numbers

  z = determinant(m)
  print *, 'det(m)= ', z
end program hello_world
\end{lstlisting}

% -------------------------------------------------------------------
\subsection{Compiling and Linking}
\label{subsec:compiling}

To use \modname{array\_mod}, the \texttt{Matrix} library archive (\filename{libMatrix.a}) and its Fortran module description files (\filename{*.mod}) must first be built, and your program must be linked against \texttt{libMatrix}, \texttt{LAPACK}, and \texttt{BLAS}.

If your program is stored in the file \filename{hello.f90} inside the \filename{Software/Tools/Matrix/} directory, build the library and compile with \texttt{gfortran} as follows:
\lstset{language=bash}
\begin{lstlisting}
make gnu
gfortran hello.f90 -L. -lMatrix -llapack -lblas -o hello
\end{lstlisting}
Notice that \lcodeb{make gnu} builds the \texttt{PFAPACK} library first, compiles all module object and \filename{.mod} files, archives them (including the \texttt{PFAPACK} objects) into \filename{libMatrix.a}, and copies the artifacts into \filename{lib/gnu/}. If \texttt{PFAPACK} has already been compiled once, you can also rebuild the module objects and archive directly via \lcodeb{make lib}.

If your source file \filename{hello.f90} is located in a different working directory and the compiled module files and \filename{libMatrix.a} reside in \filename{/path/to/Matrix/} (or \filename{/path/to/Matrix/lib/gnu/}), use the \texttt{-I} flag to specify the location of the \filename{.mod} files and the \texttt{-L} flag for the library archive path:
\lstset{language=bash}
\begin{lstlisting}
gfortran -I/path/to/Matrix/ hello.f90 -L/path/to/Matrix/ -lMatrix -llapack -lblas -o hello
\end{lstlisting}

\begin{notebox}[Compiler Flags and Multi-Architecture Targets in the Makefile]
The included \filename{Makefile} provides preconfigured targets for several compilers and high-performance computing architectures:
\begin{itemize}
  \item \textbf{GNU Fortran (\texttt{gfortran})}: Run \lcodeb{make gnu}. Uses default flags:
\lstset{language=bash}
\begin{lstlisting}
FC     = gfortran
FFLAGS = -O3 -mtune=native -funroll-loops -fopenmp -pthread
LIBS   = -llapack -lblas
\end{lstlisting}
  Note that when compiling with \texttt{-fopenmp}, you must also pass \texttt{-fopenmp} when linking your main program, and control thread concurrency at runtime via \lcodeb{export OMP_NUM_THREADS=8}.
  \item \textbf{Intel Fortran (\texttt{ifx} / \texttt{ifort} with MKL)}: Run \lcodeb{make intel}. Uses \texttt{-O3 -xHost -qopenmp -qmkl=parallel} and copies artifacts to \filename{lib/intel/}.
  \item \textbf{HPC Cluster Targets}: Specialized targets are provided in the \filename{Makefile} for AMD AOCC/\texttt{flang} with AOCL (\lcodeb{make aris}), GNU with Intel MKL (\lcodeb{make arismkl}), and NVIDIA HPC SDK \texttt{nvfortran} (\lcodeb{make nvaris}, \lcodeb{make icarus}, \lcodeb{make icarusnogpu}).
  \item \textbf{Debugging}: Run \lcodeb{make debug} (or \lcodeb{make d}) to recompile with \texttt{FFLAGS=-g}.
\end{itemize}
\end{notebox}

% -------------------------------------------------------------------
\subsection{Download}
\label{subsec:download}

The complete source code archive of the \texttt{Matrix} library can be downloaded from:
\begin{center}
  \url{https://physics.ntua.gr/konstant/PUB/Matrix.tgz}
\end{center}
The official GitHub repository is maintained at:
\begin{center}
  \url{https://github.com/KNAnagnostop/FortranMatrix}
\end{center}

% -------------------------------------------------------------------
\subsection{Hello World!}
\label{subsec:hello_world}

The complete \filename{hello.f90} program below (included in the root directory of the library) tours the essential capabilities of \modname{array\_mod}, including Gaussian random matrix generation, determinants, \texttt{BLAS}-backed matrix multiplication, eigenvalue/eigenvector decomposition and diagonalization, Pfaffian computation, Pauli matrices, Kronecker tensor products, and formatted output for terminal inspection and Mathematica import:

\lstset{language=[95]Fortran}
\begin{lstlisting}
!------------------------------------------------------------------------------
program hello_world
 use :: array_mod
 complex(8), allocatable :: M(:,:), v(:), v2(:), w(:,:)
 complex(8), allocatable :: evec(:,:), eval(:)
 complex(8), allocatable :: P(:,:), Pinv(:,:)
 complex(8), allocatable :: ldet(:), lpf(:)
 complex(8)              :: z
 complex(8), allocatable :: s0(:,:), s1(:,:), s2(:,:), s3(:,:), gamma(:,:)

 n = 8 ; allocate(M(n,n), v(n))

 ! Overloaded random_number() for complex arrays and Gaussian random numbers.
 ! Create random entries in a Gaussian distribution with sigma=1:
 call matrix_random_init                          ! seed from /dev/urandom
 call random_number(M, sigma=1._8); call random_number(v, sigma=1._8)

 z    = determinant(M)
 ldet = lndet      (M)                            ! log(det(M)) for large M

 v2   = lmatmul(M, v)                             ! use BLAS for M . v

 w    = eigenvectors(M, sortby='RF')              ! Eigenvalues + right evecs,
                                                  ! ordered by ascending Re
 eval = w(:,1 )                                   ! 1st column: eigenvalues
 evec = w(:,2:)                                   ! columns 2..n+1: evecs

 i = 3
 v2 = lmatmul(M, evec(:,i)) - eval(i) * evec(:,i) ! M.v - lambda v = 0
 call print( v2 )                                 ! print allocatable array

 ! Diagonalize M:
 P    = evec ; Pinv = inverse(P)
 M    = lmatmul(Pinv, lmatmul(M, P))              ! P^{-1} . M . P

 call print(M, FMT='F8.3', name='Mdiag')

 ! Compute Pfaffian:
 call random_number    (M)
 call antisymmetric_set(M)
 z    =   Pfaffian(M)                             ! Pfaffian of antisym M
 lpf  = lnPfaffian(M)                             ! log(Pfaffian) for large M
 call print(lpf, form='Mathematica', name='lpf')  ! copy/paste to Mathematica

 ! Tensor product. Use Pauli matrices:
 s0    = PauliMatrix(0); s1 = PauliMatrix(1)
 s2    = PauliMatrix(2); s3 = PauliMatrix(3)
 gamma = tensorprod(s2, s2, s0, s3)
 call print(gamma, form='Gamma', name='gamma')

end program hello_world
\end{lstlisting}


% Section 2: Basic Matrix Operations
%   2.1 Trace, symmetries, sort, norm
% ===================================================================
% Section 2: Basic Matrix Operations
% ===================================================================
\section{Basic Matrix Operations}
\label{sec:basic_matrix_ops}

This section documents the fundamental procedures for matrix and array manipulation, including traces, symmetry transformations and tests, diagonal and identity constructors, Pauli matrices, sorting, random shuffling, norms, and matrix multiplication.

% -------------------------------------------------------------------
\subsection{Trace, symmetries, sort, norm}
\label{subsec:trace_sym_sort_norm}

The procedures in this subsection compute matrix traces ($\Tr A$, $\Tr A^2$, $\Tr(AB)$, and connected quadratic traces), project matrices onto their traceless, Hermitian, symmetric, or antisymmetric parts, measure deviations from exact symmetry, extract or construct diagonal matrices, construct identity and Pauli matrices, sort vectors via Quicksort, shuffle arrays via the Fisher--Yates algorithm, and compute normalized $L_1$ array norms.

\begin{notebox}[Non-Square Input]
The procedures of this subsection that require a square matrix (\procname{trace}, \procname{trace2}, \procname{trace2c}, \procname{traceless}, \procname{traceless\_set}, \procname{hermitian\_set}, \procname{symmetric\_set}, \procname{antisymmetric\_set}, \procname{isHermitian}, \procname{isSymmetric}, and \procname{isAntisymmetric}) do not abort when \texttt{size(C,1) /= size(C,2)}: functions return IEEE \procname{NaN} (a scalar, or an array of the shape of \argname{C} for \procname{traceless}), and the \texttt{*\_set} subroutines overwrite \argname{C} with \procname{NaN}. Test the result with \procname{isNaN}. Empty $0\times 0$ matrices are accepted (traces are $0$).
\end{notebox}

% ===================================================================
% 1. trace
% ===================================================================
\manentry{Pure Generic Function}{trace}{trace}{Computes the trace $\Tr C = \sum_i C_{ii}$ of a real or complex matrix}{proc:trace}

\manheading{Synopsis}
\lstset{language=[95]Fortran}
\begin{lstlisting}
t = trace(C)
\end{lstlisting}

\manheading{Brief}
Returns the sum of the diagonal elements of a rank-2 real or complex array \argname{C}.

\manheading{Description}
The generic function \procname{trace} extracts the main diagonal of the input matrix \argname{C} using the internal diagonal getter (\texttt{array2\_diagonal\_get} or \texttt{array2\_diagonal\_get\_d}) and returns the sum of its elements:
\begin{align}
  \Tr C = \sum_{i=1}^{N} C_{ii}\, , \qquad N = \texttt{size(C,1)}\, .
\end{align}
Both double-precision complex (\typename{complex(8)}) and double-precision real (\typename{real(8)}) matrices are supported. The matrix \argname{C} must be square; for non-square \argname{C}, \procname{trace} returns \procname{NaN}.

\manheading{Dummy Arguments}
\begin{description}[leftmargin=2.2cm, style=nextline]
  \argitem{C}{complex(8) or real(8), dimension(:,:), intent(in)}{Input rank-2 square matrix whose trace is to be evaluated.}
  \argitem{t}{complex(8) or real(8), scalar, return value}{The computed trace $\sum_{i=1}^{\texttt{size(C,1)}} C(i,i)$, matching the data type of \argname{C} (\procname{NaN} if \argname{C} is non-square).}
\end{description}

\manheading{Example}
\lstset{language=[95]Fortran}
\begin{lstlisting}
use array_mod
complex(8) :: m1(4,4), z
real   (8) :: d1(4,4), r

call random_number(m1)
call random_number(d1)

z = trace(m1)   ! complex trace: sum_i m1(i,i)
r = trace(d1)   ! real trace:    sum_i d1(i,i)
\end{lstlisting}

% ===================================================================
% 2. trace2
% ===================================================================
\manentry{Pure Generic Function}{trace2}{trace2}{Computes the quadratic trace $\Tr(C^2)$ or bilinear trace $\Tr(C_1 C_2)$}{proc:trace2}

\manheading{Synopsis}
\lstset{language=[95]Fortran}
\begin{lstlisting}
t = trace2(C, [mtype])
t = trace2(C1, C2, [mtype])
\end{lstlisting}

\manheading{Brief}
Computes $\Tr(C^2) = \sum_{i,j} C_{ij} C_{ji}$ for a single matrix or $\Tr(C_1 C_2) = \sum_{i,j} (C_1)_{ij} (C_2)_{ji}$ for two matrices in $O(N^2)$ operations without performing full $O(N^3)$ matrix multiplication.

\manheading{Description}
The overloaded function \procname{trace2} evaluates the trace of the square of a matrix or the product of two matrices directly from their matrix elements:
\begin{align}
  \texttt{trace2(C)}     &= \Tr(C^2)    = \sum_{i=1}^{M}\sum_{j=1}^{N} C_{ij}\, C_{ji}\, , \\
  \texttt{trace2(C1,C2)} &= \Tr(C_1 C_2) = \sum_{i=1}^{M_1}\sum_{j=1}^{N_1} (C_1)_{ij}\, (C_2)_{ji}\, .
\end{align}
Both the single-matrix form and the two-matrix form validate squareness and conformability before performing any element-wise computation: \argname{C} must be square ($M = N$), and \argname{C1}, \argname{C2} must be square matrices of identical dimensions ($M_1 = N_1 = M_2 = N_2$). If any dimension check fails, \procname{trace2} immediately returns IEEE \procname{NaN}.

When the optional argument \argname{mtype} is supplied to indicate that \argname{C} (or \argname{C2} in the two-matrix call) is Hermitian (\texttt{mtype='ZH'} for complex arrays, or simply \texttt{mtype='H'}) or symmetric (\texttt{mtype='DS'} for real arrays, or \texttt{mtype='S'}), \procname{trace2} uses the identity $C_{ji} = C_{ij}^*$ (or $C_{ji} = C_{ij}$) to evaluate the sum with unit-stride memory access via intrinsic array operations:
\begin{align}
  \Tr(C^2) &= \sum_{i,j} C_{ij} C_{ij}^* = \sum_{i,j} |C_{ij}|^2 \quad (\text{Hermitian } C)\, , \qquad
  \Tr(C_1 C_2) = \sum_{i,j} (C_1)_{ij} (C_2)_{ij}^* \quad (\text{Hermitian } C_2)\, , \\
  \Tr(C^2) &= \sum_{i,j} C_{ij}^2 \quad (\text{Real symmetric } C)\, , \qquad\quad\,\,
  \Tr(C_1 C_2) = \sum_{i,j} (C_1)_{ij} (C_2)_{ij} \quad (\text{Real symmetric } C_2)\, .
\end{align}

\begin{warnbox}[No Runtime Symmetry Verification]
When \texttt{mtype='ZH'} or \texttt{mtype='DS'} is specified, \procname{trace2} does \emph{not} verify that the matrix is actually Hermitian or symmetric. Passing a non-Hermitian (or non-symmetric) matrix with this flag will yield $\Tr(C C^\dagger)$ (or $\Tr(C C^T)$) rather than $\Tr(C^2)$.
\end{warnbox}

\manheading{Dummy Arguments}
\begin{description}[leftmargin=2.2cm, style=nextline]
  \argitem{C}{complex(8) or real(8), dimension(:,:), intent(in)}{Input square matrix for the single-argument form $\Tr(C^2)$.}
  \argitem{C1, C2}{complex(8) or real(8), dimension(:,:), intent(in)}{Input square matrices of the same type and size for the two-argument form $\Tr(C_1 C_2)$.}
  \argitem{mtype}{character(*), optional, intent(in)}{Matrix symmetry hint (case-insensitive). A one-character string is the type itself; a longer string supplies it in its second character: \texttt{'H'} or \texttt{'ZH'} for complex Hermitian matrices, \texttt{'S'} or \texttt{'DS'} for real symmetric matrices, \texttt{'G'}/\texttt{'ZG'}/\texttt{'DG'} (default, also used for an absent or blank \argname{mtype}) for general matrices.}
  \argitem{t}{complex(8) or real(8), scalar, return value}{The trace $\Tr(C^2)$ or $\Tr(C_1 C_2)$ (or \procname{NaN} if \argname{C}, \argname{C1}, or \argname{C2} is non-square, or if \argname{C1} and \argname{C2} have different sizes).}
\end{description}

\manheading{Example}
\lstset{language=[95]Fortran}
\begin{lstlisting}
use array_mod
complex(8) :: m1(8,8), m2(8,8), z
real   (8) :: d1(8,8), d2(8,8), r

call random_number(m1); call random_number(m2)
call random_number(d1); call random_number(d2)

z = trace2(m1)                  ! Tr(m1^2) = sum_ij m1(i,j)*m1(j,i)
r = trace2(d1)                  ! Tr(d1^2) = sum_ij d1(i,j)*d1(j,i)
z = trace2(m1, m2)              ! Tr(m1*m2) = sum_ij m1(i,j)*m2(j,i)
r = trace2(d1, d2)              ! Tr(d1*d2) = sum_ij d1(i,j)*d2(j,i)

! Faster evaluation for Hermitian / Symmetric matrices:
call hermitian_set(m1); call hermitian_set(m2)
call symmetric_set(d1); call symmetric_set(d2)
z = trace2(m1,     mtype='ZH')
z = trace2(m1, m2, mtype='ZH')
r = trace2(d1,     mtype='DS')
r = trace2(d1, d2, mtype='DS')
\end{lstlisting}

% ===================================================================
% 3. trace2c
% ===================================================================
\manentry{Pure Generic Function}{trace2c}{trace2c}{Computes the connected quadratic trace $\Tr(\tilde{C}^2)$ of the traceless part $\tilde{C} = C - \frac{1}{N}(\Tr C) I$}{proc:trace2c}

\manheading{Synopsis}
\lstset{language=[95]Fortran}
\begin{lstlisting}
t = trace2c(C, [mtype])
\end{lstlisting}

\manheading{Brief}
Computes $\Tr(\tilde{C}^2) = \sum_{i,j} \tilde{C}_{ij} \tilde{C}_{ji}$, where $\tilde{C} = C - \frac{1}{N}(\Tr C)\, I_N$ is the traceless projection of \argname{C}.

\manheading{Description}
The generic function \procname{trace2c} creates a local copy $\tilde{C}$ of the input $N\times N$ matrix \argname{C}, removes its trace via \procname{traceless\_set} so that $\tilde{C}_{ij} = C_{ij} - \frac{1}{N}(\Tr C)\delta_{ij}$, and returns \texttt{trace2}($\tilde{C}$, \argname{mtype}):
\begin{align}
  \texttt{trace2c(C)} = \Tr\!\left[\left(C - \frac{\Tr C}{N} I_N\right)^2\right] = \Tr(C^2) - \frac{1}{N}(\Tr C)^2\, .
\end{align}
This quantity is frequently used in $U(N)$ and $SU(N)$ matrix models to measure gauge-invariant extent observables with the center-of-mass mode $\frac{1}{N}\Tr C$ subtracted.

\manheading{Dummy Arguments}
\begin{description}[leftmargin=2.2cm, style=nextline]
  \argitem{C}{complex(8) or real(8), dimension(:,:), intent(in)}{Input square matrix.}
  \argitem{mtype}{character(*), optional, intent(in)}{Matrix symmetry hint forwarded to \procname{trace2}: \texttt{'ZH'} for complex Hermitian \argname{C}, \texttt{'DS'} for real symmetric \argname{C}, or \texttt{'ZG'}/\texttt{'DG'} by default.}
  \argitem{t}{complex(8) or real(8), scalar, return value}{The connected quadratic trace $\Tr(\tilde{C}^2)$ (\procname{NaN} if \argname{C} is non-square).}
\end{description}

\manheading{Example}
\lstset{language=[95]Fortran}
\begin{lstlisting}
use array_mod
complex(8) :: m1(16,16), z
real   (8) :: d1(16,16), r

call random_number(m1); call hermitian_set(m1)
call random_number(d1); call symmetric_set(d1)

z = trace2c(m1, mtype='ZH')   ! Tr((m1 - (Tr(m1)/N)*I)^2)
r = trace2c(d1, mtype='DS')   ! Tr((d1 - (Tr(d1)/N)*I)^2)
\end{lstlisting}

% ===================================================================
% 4. traceless
% ===================================================================
\manentry{Pure Generic Function}{traceless}{traceless}{Returns the traceless part $C - \frac{1}{N}(\Tr C) I$ of a square matrix}{proc:traceless}

\manheading{Synopsis}
\lstset{language=[95]Fortran}
\begin{lstlisting}
B = traceless(C)
\end{lstlisting}

\manheading{Brief}
Returns a newly allocated matrix equal to \argname{C} with $\frac{1}{N}\Tr C$ subtracted from each diagonal entry.

\manheading{Description}
Given an $N\times N$ real or complex matrix \argname{C} (where $N = \texttt{size(C,1)}$), \procname{traceless} allocates an $N\times N$ output array \argname{B}, copies \argname{C} into \argname{B}, and subtracts the mean diagonal value $\frac{1}{N}\Tr C$ from every diagonal element:
\begin{align}
  B_{ij} = C_{ij} - \left(\frac{1}{N}\sum_{k=1}^{N} C_{kk}\right) \delta_{ij}\, ,
\end{align}
so that $\Tr B = 0$ up to machine precision. The input matrix \argname{C} is not modified.

\manheading{Dummy Arguments}
\begin{description}[leftmargin=2.2cm, style=nextline]
  \argitem{C}{complex(8) or real(8), dimension(:,:), intent(in)}{Input square $N\times N$ matrix.}
  \argitem{B}{complex(8) or real(8), dimension(:,:), allocatable, return value}{Output $N\times N$ traceless matrix of the same type as \argname{C}.}
\end{description}

\manheading{Example}
\lstset{language=[95]Fortran}
\begin{lstlisting}
use array_mod
complex(8), allocatable :: m1(:,:), m2(:,:)
real   (8), allocatable :: d1(:,:), d2(:,:)

allocate(m1(8,8), d1(8,8))
call random_number(m1); call random_number(d1)

m2 = traceless(m1)   ! trace(m2) == (0.0_8, 0.0_8)
d2 = traceless(d1)   ! trace(d2) ==  0.0_8
\end{lstlisting}

% ===================================================================
% 5. traceless_set
% ===================================================================
\manentry{Pure Generic Subroutine}{traceless\_set}{traceless\textunderscore set}{Makes a square matrix traceless in-place}{proc:traceless_set}

\manheading{Synopsis}
\lstset{language=[95]Fortran}
\begin{lstlisting}
call traceless_set(C)
\end{lstlisting}

\manheading{Brief}
Modifies the matrix \argname{C} in-place by subtracting $\frac{1}{N}\Tr C$ from its diagonal elements.

\manheading{Description}
The generic subroutine \procname{traceless\_set} projects an $N\times N$ matrix \argname{C} onto the traceless subspace in-place without allocating a temporary matrix:
\begin{align}
  C_{ii} \leftarrow C_{ii} - \frac{1}{N}\sum_{k=1}^{N} C_{kk}\, , \qquad i = 1,\ldots,N\, .
\end{align}
Off-diagonal elements $C_{ij}$ ($i\neq j$) remain unchanged.

\manheading{Dummy Arguments}
\begin{description}[leftmargin=2.2cm, style=nextline]
  \argitem{C}{complex(8) or real(8), dimension(:,:), intent(inout)}{On entry, an $N\times N$ matrix; on exit, its diagonal elements are shifted so that $\Tr C = 0$.}
\end{description}

\manheading{Example}
\lstset{language=[95]Fortran}
\begin{lstlisting}
use array_mod
complex(8) :: m1(8,8)
real   (8) :: d1(8,8)

call random_number(m1); call traceless_set(m1)
call random_number(d1); call traceless_set(d1)
\end{lstlisting}

% ===================================================================
% 6. hermitian_set
% ===================================================================
\manentry{Pure Generic Subroutine}{hermitian\_set}{hermitian\textunderscore set}{Projects or symmetrizes a complex matrix into a Hermitian matrix in-place}{proc:hermitian_set}

\manheading{Synopsis}
\lstset{language=[95]Fortran}
\begin{lstlisting}
call hermitian_set(C, [uplo])
\end{lstlisting}

\manheading{Brief}
Transforms a square complex matrix \argname{C} in-place into a Hermitian matrix ($C = C^\dagger$).

\manheading{Description}
The subroutine \procname{hermitian\_set} makes a square \typename{complex(8)} matrix \argname{C} Hermitian according to the optional character control argument \argname{uplo}:
\begin{itemize}
  \item \textbf{Default (\argname{uplo} absent, or \texttt{uplo(1:1)} $\notin \{\texttt{'U','u','L','l'}\}$)}: Replaces \argname{C} with its Hermitian part:
  \begin{align}
    C \leftarrow \frac{1}{2}\left(C + C^\dagger\right) = \frac{1}{2}\left(C + \texttt{CONJG(TRANSPOSE(C))}\right)\, .
  \end{align}
  \item \textbf{\texttt{uplo='U'} or \texttt{'u'}}: Preserves the strict upper triangular part ($i < j$) and sets the lower triangular part ($i > j$) to its complex conjugate: $C_{ij} \leftarrow C_{ji}^*$ for $i > j$.
  \item \textbf{\texttt{uplo='L'} or \texttt{'l'}}: Preserves the strict lower triangular part ($i > j$) and sets the upper triangular part ($i < j$) to its complex conjugate: $C_{ij} \leftarrow C_{ji}^*$ for $i < j$.
\end{itemize}
In all cases, the imaginary part of every diagonal element is explicitly zeroed out: $C_{ii} \leftarrow \Re(C_{ii})$.

\manheading{Dummy Arguments}
\begin{description}[leftmargin=2.2cm, style=nextline]
  \argitem{C}{complex(8), dimension(:,:), intent(inout)}{Square complex matrix to be made Hermitian in-place.}
  \argitem{uplo}{character(*), optional, intent(in)}{Controls how Hermiticity is enforced: \texttt{'U'} preserves the upper triangle, \texttt{'L'} preserves the lower triangle, and any other character (or omission) averages $\frac{1}{2}(C + C^\dagger)$.}
\end{description}

\manheading{Example}
\lstset{language=[95]Fortran}
\begin{lstlisting}
use array_mod
complex(8) :: m1(8,8), m2(8,8)

call random_number(m1)
call hermitian_set(m1)            ! m1 <- (m1 + m1^\dagger) / 2

call random_number(m2)
call hermitian_set(m2, uplo='U')  ! copy conjg of upper triangle to lower triangle
\end{lstlisting}

% ===================================================================
% 7. symmetric_set
% ===================================================================
\manentry{Pure Generic Subroutine}{symmetric\_set}{symmetric\textunderscore set}{Projects or symmetrizes a real or complex matrix into a symmetric matrix in-place}{proc:symmetric_set}

\manheading{Synopsis}
\lstset{language=[95]Fortran}
\begin{lstlisting}
call symmetric_set(C, [uplo])
\end{lstlisting}

\manheading{Brief}
Transforms a square real or complex matrix \argname{C} in-place into a symmetric matrix ($C = C^T$).

\manheading{Description}
The generic subroutine \procname{symmetric\_set} enforces transpose symmetry ($C_{ij} = C_{ji}$) on a square \typename{complex(8)} or \typename{real(8)} matrix \argname{C} according to the optional argument \argname{uplo}:
\begin{itemize}
  \item \textbf{Default (\argname{uplo} absent, or \texttt{uplo(1:1)} $\notin \{\texttt{'U','u','L','l'}\}$)}: Replaces \argname{C} with its symmetric part:
  \begin{align}
    C \leftarrow \frac{1}{2}\left(C + C^T\right)\, .
  \end{align}
  \item \textbf{\texttt{uplo='U'} or \texttt{'u'}}: Preserves the upper triangular part ($i < j$) and sets $C_{ij} \leftarrow C_{ji}$ for $i > j$.
  \item \textbf{\texttt{uplo='L'} or \texttt{'l'}}: Preserves the lower triangular part ($i > j$) and sets $C_{ij} \leftarrow C_{ji}$ for $i < j$.
\end{itemize}
Diagonal elements $C_{ii}$ are left untouched.

\manheading{Dummy Arguments}
\begin{description}[leftmargin=2.2cm, style=nextline]
  \argitem{C}{complex(8) or real(8), dimension(:,:), intent(inout)}{Square matrix to be made symmetric in-place.}
  \argitem{uplo}{character(*), optional, intent(in)}{Triangular source selector: \texttt{'U'} (upper), \texttt{'L'} (lower), or omitted/other for $\frac{1}{2}(C + C^T)$.}
\end{description}

\manheading{Example}
\lstset{language=[95]Fortran}
\begin{lstlisting}
use array_mod
complex(8) :: m1(8,8)
real   (8) :: d1(8,8)

call random_number(m1); call symmetric_set(m1)
call random_number(d1); call symmetric_set(d1, uplo='L')
\end{lstlisting}

% ===================================================================
% 8. antisymmetric_set
% ===================================================================
\manentry{Pure Generic Subroutine}{antisymmetric\_set}{antisymmetric\textunderscore set}{Projects or skew-symmetrizes a real or complex matrix into an antisymmetric matrix in-place}{proc:antisymmetric_set}

\manheading{Synopsis}
\lstset{language=[95]Fortran}
\begin{lstlisting}
call antisymmetric_set(C, [uplo])
\end{lstlisting}

\manheading{Brief}
Transforms a square real or complex matrix \argname{C} in-place into an antisymmetric (skew-symmetric) matrix ($C = -C^T$).

\manheading{Description}
The generic subroutine \procname{antisymmetric\_set} enforces skew-symmetry ($C_{ij} = -C_{ji}$) on a square \typename{complex(8)} or \typename{real(8)} matrix \argname{C} according to the optional argument \argname{uplo}:
\begin{itemize}
  \item \textbf{Default (\argname{uplo} absent, or \texttt{uplo(1:1)} $\notin \{\texttt{'U','u','L','l'}\}$)}: Replaces \argname{C} with its antisymmetric part:
  \begin{align}
    C \leftarrow \frac{1}{2}\left(C - C^T\right)\, .
  \end{align}
  \item \textbf{\texttt{uplo='U'} or \texttt{'u'}}: Preserves the strict upper triangular part ($i < j$) and sets $C_{ij} \leftarrow -C_{ji}$ for $i > j$.
  \item \textbf{\texttt{uplo='L'} or \texttt{'l'}}: Preserves the strict lower triangular part ($i > j$) and sets $C_{ij} \leftarrow -C_{ji}$ for $i < j$.
\end{itemize}
In all cases, the diagonal entries are explicitly set to zero: $C_{ii} \leftarrow 0$.

\manheading{Dummy Arguments}
\begin{description}[leftmargin=2.2cm, style=nextline]
  \argitem{C}{complex(8) or real(8), dimension(:,:), intent(inout)}{Square matrix to be made antisymmetric in-place.}
  \argitem{uplo}{character(*), optional, intent(in)}{Triangular source selector: \texttt{'U'} (upper), \texttt{'L'} (lower), or omitted/other for $\frac{1}{2}(C - C^T)$.}
\end{description}

\manheading{Example}
\lstset{language=[95]Fortran}
\begin{lstlisting}
use array_mod
complex(8) :: m1(8,8)
real   (8) :: d1(8,8)

call random_number(m1); call antisymmetric_set(m1)
call random_number(d1); call antisymmetric_set(d1, uplo='U')
\end{lstlisting}

% ===================================================================
% 9. hermitian
% ===================================================================
\manentry{Pure Generic Function}{hermitian}{hermitian}{Returns the Hermitian conjugate $C^\dagger$ of a complex matrix}{proc:hermitian}

\manheading{Synopsis}
\lstset{language=[95]Fortran}
\begin{lstlisting}
CH = hermitian(C)
\end{lstlisting}

\manheading{Brief}
Computes and returns the conjugate transpose $C^\dagger = (C^T)^*$ of a complex matrix \argname{C}.

\manheading{Description}
Given an $M\times N$ \typename{complex(8)} matrix \argname{C}, \procname{hermitian} allocates and returns an $N\times M$ \typename{complex(8)} matrix \argname{CH} equal to \texttt{CONJG(TRANSPOSE(C))}, so that $(C_H)_{ij} = C_{ji}^*$. Note that \procname{hermitian} returns the adjoint matrix $C^\dagger$ itself (not the Hermitian projection $\frac{1}{2}(C+C^\dagger)$, which is performed by \procname{hermitian\_set}).

\manheading{Dummy Arguments}
\begin{description}[leftmargin=2.2cm, style=nextline]
  \argitem{C}{complex(8), dimension(:,:), intent(in)}{Input $M\times N$ complex matrix.}
  \argitem{CH}{complex(8), dimension(:,:), allocatable, return value}{Output $N\times M$ Hermitian conjugate matrix $C^\dagger$.}
\end{description}

\manheading{Example}
\lstset{language=[95]Fortran}
\begin{lstlisting}
use array_mod
complex(8), allocatable :: m1(:,:), m2(:,:)

allocate(m1(4,6))
call random_number(m1)
m2 = hermitian(m1)   ! m2 has shape (6,4) and equals m1^\dagger
\end{lstlisting}

% ===================================================================
% 10. symmetric
% ===================================================================
\manentry{Pure Generic Function}{symmetric}{symmetric}{Returns the transpose $C^T$ of a real or complex matrix}{proc:symmetric}

\manheading{Synopsis}
\lstset{language=[95]Fortran}
\begin{lstlisting}
CS = symmetric(C)
\end{lstlisting}

\manheading{Brief}
Computes and returns the matrix transpose $C^T$ of a real or complex matrix \argname{C}.

\manheading{Description}
Given an $M\times N$ \typename{complex(8)} or \typename{real(8)} matrix \argname{C}, \procname{symmetric} allocates and returns an $N\times M$ matrix \argname{CS} equal to \texttt{TRANSPOSE(C)}, so that $(C_S)_{ij} = C_{ji}$.

\manheading{Dummy Arguments}
\begin{description}[leftmargin=2.2cm, style=nextline]
  \argitem{C}{complex(8) or real(8), dimension(:,:), intent(in)}{Input $M\times N$ matrix.}
  \argitem{CS}{complex(8) or real(8), dimension(:,:), allocatable, return value}{Output $N\times M$ transposed matrix $C^T$ of the same type as \argname{C}.}
\end{description}

\manheading{Example}
\lstset{language=[95]Fortran}
\begin{lstlisting}
use array_mod
complex(8), allocatable :: m1(:,:), m2(:,:)
real   (8), allocatable :: d1(:,:), d2(:,:)

allocate(m1(4,6), d1(4,6))
call random_number(m1); call random_number(d1)
m2 = symmetric(m1)   ! returns m1^T
d2 = symmetric(d1)   ! returns d1^T
\end{lstlisting}

% ===================================================================
% 11. antisymmetric
% ===================================================================
\manentry{Pure Generic Function}{antisymmetric}{antisymmetric}{Returns the negative transpose $-C^T$ of a real or complex matrix}{proc:antisymmetric}

\manheading{Synopsis}
\lstset{language=[95]Fortran}
\begin{lstlisting}
CS = antisymmetric(C)
\end{lstlisting}

\manheading{Brief}
Computes and returns the negative transpose $-C^T$ of a real or complex matrix \argname{C}.

\manheading{Description}
Given an $M\times N$ \typename{complex(8)} or \typename{real(8)} matrix \argname{C}, \procname{antisymmetric} allocates and returns an $N\times M$ matrix \argname{CS} equal to \texttt{-TRANSPOSE(C)}, so that $(C_S)_{ij} = -C_{ji}$. For an antisymmetric matrix, \texttt{antisymmetric(C)} equals \argname{C}.

\manheading{Dummy Arguments}
\begin{description}[leftmargin=2.2cm, style=nextline]
  \argitem{C}{complex(8) or real(8), dimension(:,:), intent(in)}{Input $M\times N$ matrix.}
  \argitem{CS}{complex(8) or real(8), dimension(:,:), allocatable, return value}{Output $N\times M$ negative transposed matrix $-C^T$ of the same type as \argname{C}.}
\end{description}

\manheading{Example}
\lstset{language=[95]Fortran}
\begin{lstlisting}
use array_mod
complex(8), allocatable :: m1(:,:), m2(:,:)
real   (8), allocatable :: d1(:,:), d2(:,:)

allocate(m1(4,4), d1(4,4))
call random_number(m1); call random_number(d1)
m2 = antisymmetric(m1)   ! returns -m1^T
d2 = antisymmetric(d1)   ! returns -d1^T
\end{lstlisting}

% ===================================================================
% 12. isHermitian
% ===================================================================
\manentry{Pure Generic Function}{isHermitian}{isHermitian}{Measures the deviation $\norm{C - C^\dagger}$ of a complex matrix from Hermiticity}{proc:isHermitian}

\manheading{Synopsis}
\lstset{language=[95]Fortran}
\begin{lstlisting}
r = isHermitian(C)
\end{lstlisting}

\manheading{Brief}
Returns the normalized $L_1$ norm \texttt{norm(C - CONJG(TRANSPOSE(C)))} as a real number.

\manheading{Description}
The function \procname{isHermitian} quantifies how close a square \typename{complex(8)} matrix \argname{C} is to being Hermitian by computing:
\begin{align}
  r = \texttt{norm}(C - C^\dagger) = \frac{1}{N^2}\sum_{i,j=1}^{N} \left|C_{ij} - C_{ji}^*\right|\, .
\end{align}
An exact Hermitian matrix yields $r = 0.0\_8$ (or $r \lesssim 10^{-15}$ when subject to floating-point rounding).

\manheading{Dummy Arguments}
\begin{description}[leftmargin=2.2cm, style=nextline]
  \argitem{C}{complex(8), dimension(:,:), intent(in)}{Input square complex matrix.}
  \argitem{r}{real(8), scalar, return value}{The non-negative deviation norm $\texttt{norm}(C - C^\dagger)$.}
\end{description}

\manheading{Example}
\lstset{language=[95]Fortran}
\begin{lstlisting}
use array_mod
complex(8) :: m1(8,8)
real   (8) :: r

call random_number(m1)
call hermitian_set(m1)
r = isHermitian(m1)   ! r == 0.0_8
\end{lstlisting}

% ===================================================================
% 13. isSymmetric
% ===================================================================
\manentry{Pure Generic Function}{isSymmetric}{isSymmetric}{Measures the deviation $\norm{C - C^T}$ of a real or complex matrix from symmetry}{proc:isSymmetric}

\manheading{Synopsis}
\lstset{language=[95]Fortran}
\begin{lstlisting}
r = isSymmetric(C)
\end{lstlisting}

\manheading{Brief}
Returns the normalized $L_1$ norm \texttt{norm(C - TRANSPOSE(C))} as a \typename{real(8)} scalar.

\manheading{Description}
The generic function \procname{isSymmetric} quantifies how close a square \typename{complex(8)} or \typename{real(8)} matrix \argname{C} is to being symmetric ($C = C^T$) by evaluating:
\begin{align}
  r = \texttt{norm}(C - C^T) = \frac{1}{N^2}\sum_{i,j=1}^{N} \left|C_{ij} - C_{ji}\right|\, .
\end{align}
An exact symmetric matrix yields $r = 0.0\_8$.

\manheading{Dummy Arguments}
\begin{description}[leftmargin=2.2cm, style=nextline]
  \argitem{C}{complex(8) or real(8), dimension(:,:), intent(in)}{Input square matrix.}
  \argitem{r}{real(8), scalar, return value}{The non-negative deviation norm $\texttt{norm}(C - C^T)$.}
\end{description}

\manheading{Example}
\lstset{language=[95]Fortran}
\begin{lstlisting}
use array_mod
complex(8) :: m1(8,8)
real   (8) :: d1(8,8), r1, r2

call random_number(m1); call symmetric_set(m1)
call random_number(d1); call symmetric_set(d1)
r1 = isSymmetric(m1)   ! r1 == 0.0_8
r2 = isSymmetric(d1)   ! r2 == 0.0_8
\end{lstlisting}

% ===================================================================
% 14. isAntisymmetric
% ===================================================================
\manentry{Pure Generic Function}{isAntisymmetric}{isAntisymmetric}{Measures the deviation $\norm{C + C^T}$ of a real or complex matrix from antisymmetry}{proc:isAntisymmetric}

\manheading{Synopsis}
\lstset{language=[95]Fortran}
\begin{lstlisting}
r = isAntiSymmetric(C)
\end{lstlisting}

\manheading{Brief}
Returns the normalized $L_1$ norm \texttt{norm(C + TRANSPOSE(C))} as a \typename{real(8)} scalar.

\manheading{Description}
The generic function \procname{isAntisymmetric} (case-insensitive in Fortran: \texttt{isAntiSymmetric}) quantifies how close a square \typename{complex(8)} or \typename{real(8)} matrix \argname{C} is to being skew-symmetric ($C = -C^T$) by evaluating:
\begin{align}
  r = \texttt{norm}(C + C^T) = \frac{1}{N^2}\sum_{i,j=1}^{N} \left|C_{ij} + C_{ji}\right|\, .
\end{align}
An exact antisymmetric matrix yields $r = 0.0\_8$.

\manheading{Dummy Arguments}
\begin{description}[leftmargin=2.2cm, style=nextline]
  \argitem{C}{complex(8) or real(8), dimension(:,:), intent(in)}{Input square matrix.}
  \argitem{r}{real(8), scalar, return value}{The non-negative deviation norm $\texttt{norm}(C + C^T)$.}
\end{description}

\manheading{Example}
\lstset{language=[95]Fortran}
\begin{lstlisting}
use array_mod
complex(8) :: m1(8,8)
real   (8) :: d1(8,8), r1, r2

call random_number(m1); call antisymmetric_set(m1)
call random_number(d1); call antisymmetric_set(d1)
r1 = isAntiSymmetric(m1)   ! r1 == 0.0_8
r2 = isAntiSymmetric(d1)   ! r2 == 0.0_8
\end{lstlisting}

% ===================================================================
% 15. diagonal
% ===================================================================
\manentry{Pure Generic Function}{diagonal}{diagonal}{Extracts the diagonal of a matrix as a vector, or constructs a diagonal matrix from a vector or scalar}{proc:diagonal}

\manheading{Synopsis}
\lstset{language=[95]Fortran}
\begin{lstlisting}
d = diagonal(C)      ! Extract diagonal vector from matrix C(:,:)
C = diagonal(d)      ! Construct diagonal matrix from vector d(:)
C = diagonal(z, n)   ! Construct n x n complex diagonal matrix z * I_n
C = diagonal(r, n)   ! Construct n x n real    diagonal matrix r * I_n
\end{lstlisting}

\manheading{Brief}
Overloaded function that either extracts the main diagonal of a rank-2 matrix into a rank-1 vector, or builds a rank-2 diagonal matrix from a rank-1 vector or a constant scalar value.

\manheading{Description}
The generic interface \procname{diagonal} resolves to six module procedures depending on the rank and type of its arguments:
\begin{enumerate}
  \item \textbf{Matrix $\to$ Vector (\texttt{d = diagonal(C)})}: Given a rank-2 $M\times N'$ matrix \argname{C}, allocates and returns a 1D vector \argname{d} of length $N = \min(M, N')$ (\texttt{min(size(C,1), size(C,2))}, so that non-square matrices are accepted) matching the data type of \argname{C} (\typename{complex(8)} for complex \argname{C}, \typename{real(8)} for real \argname{C}), with $d(i) = C(i,i)$ for $i=1,\ldots,N$.
  \item \textbf{Vector $\to$ Matrix (\texttt{C = diagonal(d)})}: Given a 1D \typename{complex(8)} or \typename{real(8)} vector \argname{d} of length $N = \texttt{size(d)}$, allocates and returns an $N\times N$ matrix \argname{C} initialized to zero off the diagonal and $C(i,i) = d(i)$ on the diagonal.
  \item \textbf{Scalar $\to$ Matrix (\texttt{C = diagonal(z, n)} or \texttt{diagonal(r, n)})}: Given a \typename{complex(8)} scalar \argname{z} (or \typename{real(8)} scalar \argname{r}) and an integer dimension \argname{n}, allocates and returns an $n\times n$ \typename{complex(8)} (or \typename{real(8)}) diagonal matrix with $C(i,i) = z$ (or $r$) and $C(i,j) = 0$ for $i\neq j$.
\end{enumerate}

\manheading{Dummy Arguments}
\begin{description}[leftmargin=2.2cm, style=nextline]
  \argitem{C}{complex(8) or real(8), dimension(:,:), intent(in) [or allocatable return value]}{In \texttt{diagonal(C)}, the input rank-2 matrix; in \texttt{diagonal(d)} and \texttt{diagonal(s,n)}, the returned allocatable $N\times N$ diagonal matrix.}
  \argitem{d}{complex(8) or real(8), dimension(:), intent(in) [or allocatable return value]}{In \texttt{diagonal(d)}, the input 1D vector of diagonal entries; in \texttt{diagonal(C)}, the returned 1D vector of length \texttt{min(size(C,1), size(C,2))}.}
  \argitem{z, r}{complex(8) or real(8), scalar, intent(in)}{Constant value placed along the main diagonal in the two-argument constructor form.}
  \argitem{n}{integer, scalar, intent(in)}{Dimension $n$ of the $n\times n$ matrix constructed in the two-argument form.}
\end{description}

\manheading{Example}
\lstset{language=[95]Fortran}
\begin{lstlisting}
use array_mod
complex(8), allocatable :: m1(:,:), w1(:)
real   (8), allocatable :: d1(:,:), v1(:)
complex(8)              :: z = (2.0_8, -1.0_8)
real   (8)              :: r = 3.5_8
integer                 :: n = 4

allocate(m1(n,n), d1(n,n))
call random_number(m1); call random_number(d1)

w1 = diagonal(m1)     ! extract diagonal: w1(i) = m1(i,i)
v1 = diagonal(d1)     ! extract diagonal: v1(i) = d1(i,i)

m1 = diagonal(w1)     ! build diagonal matrix: m1(i,i) = w1(i), m1(i,j)=0
d1 = diagonal(v1)     ! build diagonal matrix: d1(i,i) = v1(i), d1(i,j)=0

m1 = diagonal(z, n)   ! build 4x4 matrix with z on diagonal
d1 = diagonal(r, n)   ! build 4x4 matrix with r on diagonal
\end{lstlisting}

% ===================================================================
% 16. diagonalMatrix
% ===================================================================
\manentry{Pure Generic Function}{diagonalMatrix}{diagonalMatrix}{Constructs a diagonal matrix from the diagonal elements of an existing matrix}{proc:diagonalMatrix}

\manheading{Synopsis}
\lstset{language=[95]Fortran}
\begin{lstlisting}
D = diagonalMatrix(C)
\end{lstlisting}

\manheading{Brief}
Returns a diagonal matrix $D$ of the same shape as $C$ having the same main diagonal as $C$ ($D_{ii} = C_{ii}$) and zero off-diagonal entries ($D_{ij} = 0$ for $i\neq j$).

\manheading{Description}
Given an $M\times N'$ \typename{complex(8)} or \typename{real(8)} matrix \argname{C}, \procname{diagonalMatrix} allocates an $M\times N'$ matrix \argname{D} of the same type and shape, sets all elements to zero, and copies $D(i,i) = C(i,i)$ for $i=1,\ldots,\min(M,N')$. For a square matrix this is equivalent to \texttt{diagonal(diagonal(C))} but avoids allocating an intermediate vector; for a non-square matrix \texttt{diagonal(diagonal(C))} would instead be a square $\min(M,N')\times\min(M,N')$ matrix.

\manheading{Dummy Arguments}
\begin{description}[leftmargin=2.2cm, style=nextline]
  \argitem{C}{complex(8) or real(8), dimension(:,:), intent(in)}{Input matrix whose main diagonal is extracted.}
  \argitem{D}{complex(8) or real(8), dimension(:,:), allocatable, return value}{Output diagonal matrix of the same shape as \argname{C}, with $D_{ij} = C_{ii}\delta_{ij}$.}
\end{description}

\manheading{Example}
\lstset{language=[95]Fortran}
\begin{lstlisting}
use array_mod
complex(8), allocatable :: m1(:,:), m2(:,:)
real   (8), allocatable :: d1(:,:), d2(:,:)

allocate(m1(4,4), d1(4,4))
call random_number(m1); call random_number(d1)

m2 = diagonalMatrix(m1)   ! m2(i,i) = m1(i,i), m2(i,j) = 0 for i /= j
d2 = diagonalMatrix(d1)   ! d2(i,i) = d1(i,i), d2(i,j) = 0 for i /= j
\end{lstlisting}

% ===================================================================
% 17. identitymatrix, cidentitymatrix, didentitymatrix
% ===================================================================
\manentry{Pure Generic Functions}{identitymatrix, cidentitymatrix, didentitymatrix}{identitymatrix, cidentitymatrix, didentitymatrix}{Construct $n\times n$ complex or real identity matrices $I_n$}{proc:identitymatrix}
\addcontentsline{toc}{subsubsection}{\texorpdfstring{\texttt{\textbf{cidentitymatrix}} \textnormal{\footnotesize--- Complex $n\times n$ identity matrix}}{cidentitymatrix - Complex n x n identity matrix}}
\addcontentsline{toc}{subsubsection}{\texorpdfstring{\texttt{\textbf{didentitymatrix}} \textnormal{\footnotesize--- Real $n\times n$ identity matrix}}{didentitymatrix - Real n x n identity matrix}}

\manheading{Synopsis}
\lstset{language=[95]Fortran}
\begin{lstlisting}
m1 =  identitymatrix(n)   ! returns complex(8) n x n identity matrix
m1 = cidentitymatrix(n)   ! returns complex(8) n x n identity matrix
d1 = didentitymatrix(n)   ! returns real   (8) n x n identity matrix
\end{lstlisting}

\manheading{Brief}
Allocate and return an $n\times n$ identity matrix $I_n$ ($C_{ij} = \delta_{ij}$) of type \typename{complex(8)} (\procname{identitymatrix}, \procname{cidentitymatrix}) or \typename{real(8)} (\procname{didentitymatrix}).

\manheading{Description}
These three pure functions construct the $n\times n$ unit matrix $\UnitMatrix_n$:
\begin{itemize}
  \item \procname{identitymatrix(n)} and \procname{cidentitymatrix(n)} both invoke the internal complex identity constructor \texttt{array2\_diagonal\_set\_identity\_complex\_matrix(n)} and return an allocatable \typename{complex(8)} array of shape $(n,n)$ with $(1.0\_8, 0.0\_8)$ on the main diagonal and $(0.0\_8, 0.0\_8)$ elsewhere.
  \item \procname{didentitymatrix(n)} invokes \texttt{array2\_diagonal\_set\_identity\_real\_matrix(n)} and returns an allocatable \typename{real(8)} array of shape $(n,n)$ with $1.0\_8$ on the main diagonal and $0.0\_8$ elsewhere. Note that in Fortran assignment \texttt{m1 = didentitymatrix(n)}, a real identity matrix can also be assigned directly to a complex array.
\end{itemize}

\manheading{Dummy Arguments}
\begin{description}[leftmargin=2.2cm, style=nextline]
  \argitem{n}{integer, scalar, intent(in)}{The matrix dimension $n$ (for $n \le 0$ an empty $0\times 0$ array is returned).}
  \argitem{C}{complex(8) or real(8), dimension(:,:), allocatable, return value}{The allocated $n\times n$ identity matrix $I_n$.}
\end{description}

\manheading{Example}
\lstset{language=[95]Fortran}
\begin{lstlisting}
use array_mod
complex(8), allocatable :: m1(:,:), m2(:,:)
real   (8), allocatable :: d1(:,:)

m1 =  identitymatrix(8)   ! 8x8 complex(8) identity matrix
m2 = cidentitymatrix(8)   ! 8x8 complex(8) identity matrix
d1 = didentitymatrix(8)   ! 8x8 real   (8) identity matrix
\end{lstlisting}

% ===================================================================
% 18. paulimatrix
% ===================================================================
\manentry{Pure Generic Function}{paulimatrix}{paulimatrix}{Returns the $2\times 2$ complex Pauli matrices $\sigma_0, \sigma_1, \sigma_2, \sigma_3$}{proc:paulimatrix}

\manheading{Synopsis}
\lstset{language=[95]Fortran}
\begin{lstlisting}
s = PauliMatrix(n)
\end{lstlisting}

\manheading{Brief}
Returns the $2\times 2$ complex Pauli matrix $\sigma_n$ for $n \in \{0, 1, 2, 3, 4\}$, or a $2\times 2$ matrix of \procname{NaN} values for any other integer $n$.

\manheading{Description}
The pure function \procname{paulimatrix} (case-insensitive: \texttt{PauliMatrix}) returns a fixed-size $2\times 2$ \typename{complex(8)} array corresponding to the standard Pauli basis:
\begin{align}
  \sigma_0 = \sigma_4 = \begin{pmatrix} 1 & 0 \\ 0 & 1 \end{pmatrix}\, , \quad
  \sigma_1 = \begin{pmatrix} 0 & 1 \\ 1 & 0 \end{pmatrix}\, , \quad
  \sigma_2 = \begin{pmatrix} 0 & -i \\ i & 0 \end{pmatrix}\, , \quad
  \sigma_3 = \begin{pmatrix} 1 & 0 \\ 0 & -1 \end{pmatrix}\, .
\end{align}
Notice that both $n=0$ and $n=4$ return the $2\times 2$ unit matrix $1_2$. If $n \notin \{0,1,2,3,4\}$, all four entries of the returned $2\times 2$ matrix are set to IEEE quiet \procname{NaN}.

\manheading{Dummy Arguments}
\begin{description}[leftmargin=2.2cm, style=nextline]
  \argitem{n}{integer, scalar, intent(in)}{Index of the Pauli matrix ($0, 1, 2, 3$, or $4$).}
  \argitem{C}{complex(8), dimension(2,2), return value}{The $2\times 2$ Pauli matrix $\sigma_n$, or \procname{NaN} if $n\notin \{0,1,2,3,4\}$.}
\end{description}

\manheading{Example}
\lstset{language=[95]Fortran}
\begin{lstlisting}
use array_mod
complex(8), allocatable :: s0(:,:), s1(:,:), s2(:,:), s3(:,:)

s0 = PauliMatrix(0)   ! 2x2 identity matrix
s1 = PauliMatrix(1)   ! sigma_1
s2 = PauliMatrix(2)   ! sigma_2
s3 = PauliMatrix(3)   ! sigma_3
\end{lstlisting}

% ===================================================================
% 19. sort
% ===================================================================
\manentry{Generic Function}{sort}{sort}{Sorts a 1D real or complex array using Quicksort with optional permutation index tracking}{proc:sort}

\manheading{Synopsis}
\lstset{language=[95]Fortran}
\begin{lstlisting}
D = sort(C, [P], [by])
\end{lstlisting}

\manheading{Brief}
Sorts a rank-1 \typename{complex(8)} or \typename{real(8)} array \argname{C} in ascending (forward) or descending (reverse) order according to a specified criterion \argname{by}, optionally recording the original indices in the permutation array \argname{P}.

\manheading{Description}
The generic function \procname{sort} allocates and returns a sorted copy \argname{D} of the 1D array \argname{C} by precomputing a real sort key array and running a bounds-guarded, tail-call-eliminated Quicksort (safe under aggressive compiler optimizations and FMA instructions such as \texttt{-march=native}). The sorting criterion and direction are controlled by the two-character string \argname{by}\texttt{(1:2)}:
\begin{itemize}
  \item \textbf{First character \texttt{by(1:1)} (Comparison Key)}:
  \begin{itemize}
    \item For \typename{complex(8)} arrays:
      \begin{itemize}
        \item \texttt{'M'} or \texttt{'m'}: Sort by modulus (absolute value) $|C(i)|$.
        \item \texttt{'R'} or \texttt{'r'}: Sort by real part $\Re(C(i))$.
        \item \texttt{'I'} or \texttt{'i'}: Sort by imaginary part $\Im(C(i))$.
      \end{itemize}
    \item For \typename{real(8)} arrays:
      \begin{itemize}
        \item \texttt{'R'}, \texttt{'r'}, \texttt{'V'}, or \texttt{'v'}: Sort by signed real value $C(i)$.
        \item \texttt{'A'}, \texttt{'a'}, \texttt{'M'}, or \texttt{'m'}: Sort by absolute value (modulus) $|C(i)|$.
      \end{itemize}
  \end{itemize}
  \item \textbf{Second character \texttt{by(2:2)} (Ordering Direction)}:
  \begin{itemize}
    \item \texttt{'F'} or \texttt{'f'}: Forward (ascending) order, from smallest to largest.
    \item \texttt{'R'} or \texttt{'r'}: Reverse (descending) order, from largest to smallest.
  \end{itemize}
  \item \textbf{Default values}: Both characters of \argname{by(1:2)} are handled case-insensitively (including \texttt{'MF'}/\texttt{'MR'} as well as \texttt{'AF'}/\texttt{'AR'} for real arrays). When \argname{by} is omitted, the default is \texttt{by='MF'} (ascending modulus) for \typename{complex(8)} arrays and \texttt{by='RF'} (ascending real value) for \typename{real(8)} arrays. Passing an unrecognized \argname{by} string or a string shorter than 2 characters terminates execution via \texttt{matrix\_error}.
\end{itemize}

When the optional integer array \argname{P} is provided, its size is verified against \argname{C} (\texttt{size(P) == size(C)}, otherwise execution terminates via \texttt{matrix\_error}); \argname{P} is initialized to $[1, 2, \ldots, \texttt{size(C)}]$ and permuted alongside \argname{D} during sorting so that on return:
\begin{align}
  D(i) = C(P(i))\, , \qquad i = 1,\ldots,\texttt{size(C)}\, .
\end{align}
Because \argname{P} is tracked directly during partitioning, all indices $1,\ldots,N$ are preserved without repetition even when \argname{C} contains degenerate (equal) elements. Empty ($N = 0$) and single-element ($N = 1$) arrays return cleanly (with \argname{P} empty or \texttt{[1]}, respectively). Any elements whose sort key is IEEE \texttt{NaN} are partitioned to the end of \argname{D} (preserving their original indices in \argname{P}), while all non-\texttt{NaN} elements are sorted in the requested order.

\manheading{Dummy Arguments}
\begin{description}[leftmargin=2.2cm, style=nextline]
  \argitem{C}{complex(8) or real(8), dimension(:), intent(in)}{Input 1D array of length $N$ to be sorted.}
  \argitem{P}{integer, dimension(:), optional, intent(out)}{Optional integer permutation array of length $N$, pre-allocated by the caller (\texttt{size(P) == size(C)} is checked at runtime), such that $D(i) = C(P(i))$.}
  \argitem{by}{character(*), optional, intent(in)}{Two-character sorting specification: \texttt{'MF'}, \texttt{'MR'}, \texttt{'RF'}, \texttt{'RR'}, \texttt{'IF'}, \texttt{'IR'} for complex arrays; \texttt{'RF'} (or \texttt{'VF'}), \texttt{'RR'} (or \texttt{'VR'}), \texttt{'AF'} (or \texttt{'MF'}), \texttt{'AR'} (or \texttt{'MR'}) for real arrays.}
  \argitem{D}{complex(8) or real(8), dimension(:), allocatable, return value}{Sorted 1D array of length $N$ matching the type of \argname{C}.}
\end{description}

\manheading{Example}
\lstset{language=[95]Fortran}
\begin{lstlisting}
use array_mod
complex(8) :: w(6), w1(6)
real   (8) :: v(6), v1(6)
integer    :: pos(6)

call random_number(w, sigma=1.0_8)
call random_number(v, sigma=1.0_8)

w1 = sort(w)                     ! default 'MF': ascending by |w(i)|
w1 = sort(w, by='MR')            ! descending by |w(i)|
v1 = sort(v, by='RF')            ! ascending by signed value v(i)

w1 = sort(w, by='IR', P=pos)     ! descending by Im(w(i)); w1(i) == w(pos(i))
v1 = sort(v, by='AF', P=pos)     ! ascending by |v(i)|;    v1(i) == v(pos(i))
\end{lstlisting}

% ===================================================================
% 20. random_sort
% ===================================================================
\manentry{Generic Function}{random\_sort}{random\textunderscore sort}{Randomly shuffles a 1D array, character string, or integer range using the Fisher--Yates algorithm}{proc:random_sort}

\manheading{Synopsis}
\lstset{language=[95]Fortran}
\begin{lstlisting}
l1  = random_sort(l)         ! Shuffle integer(4) 1D array l(:)
v1  = random_sort(v)         ! Shuffle real   (8) 1D array v(:)
w1  = random_sort(w)         ! Shuffle complex(8) 1D array w(:)
str = random_sort(string)    ! Shuffle characters of a string
l1  = random_sort(n)         ! Random permutation of integers 1..n
l1  = random_sort(m, n)      ! Random permutation of integers m..n (m <= n)
\end{lstlisting}

\manheading{Brief}
Returns a randomly permuted copy of an integer, real, or complex 1D array, a character string, or an integer range $[1,n]$ or $[m,n]$ via the modern Fisher--Yates shuffle.

\manheading{Description}
The generic function \procname{random\_sort} implements the unbiased $O(N)$ Fisher--Yates (Knuth) shuffle algorithm: iterating from index $i = 1$ to $N-1$, a uniform random index $j \in \{i, i+1, \ldots, N\}$ is drawn via \texttt{j = i + int(r * (N - i + 1))} and elements $i$ and $j$ are swapped. It is overloaded for five distinct input forms:
\begin{enumerate}
  \item \textbf{1D Array Shuffle (\texttt{random\_sort(array\_in)})}: Accepts a 1D \typename{integer}, \typename{real(8)}, or \typename{complex(8)} array and returns an allocatable copy with its elements randomly permuted.
  \item \textbf{Character String Shuffle (\texttt{random\_sort(string)})}: Trims trailing whitespace from \argname{string} and returns an allocatable character string of length \texttt{len(trim(string))} with its characters randomly shuffled.
  \item \textbf{Integer Range Permutation (\texttt{random\_sort(n1, [n2])})}:
  \begin{itemize}
    \item If called with a single integer \texttt{random\_sort(n)}, constructs the sequence $[1, 2, \ldots, n]$ and returns a random permutation of length $n$. If $n < 1$, execution terminates via \texttt{matrix\_error}.
    \item If called with two integers \texttt{random\_sort(m, n)}, constructs the sequence $[m, m+1, \ldots, n]$ of length $n - m + 1$ and returns a random permutation. If $n - m + 1 < 1$, execution terminates via \texttt{matrix\_error}.
  \end{itemize}
\end{enumerate}

\manheading{Dummy Arguments}
\begin{description}[leftmargin=2.2cm, style=nextline]
  \argitem{array\_in}{integer, real(8), or complex(8) dimension(:), or character(*), intent(in)}{Input 1D array or character string to be shuffled (left unchanged).}
  \argitem{n1}{integer, scalar, intent(in)}{Upper bound $n$ of the range $1\ldots n$ when \argname{n2} is absent, or lower bound $m$ of the range $m\ldots n$ when \argname{n2} is present.}
  \argitem{n2}{integer, scalar, optional, intent(in)}{Upper bound $n$ of the integer range $m\ldots n$ ($n \ge m$).}
  \argitem{array}{allocatable 1D array or character(:), return value}{The randomly shuffled array or string.}
\end{description}

\manheading{Example}
\lstset{language=[95]Fortran}
\begin{lstlisting}
use array_mod
integer   , allocatable :: l1(:)
real   (8), allocatable :: v1(:)
complex(8), allocatable :: w1(:)
character(:), allocatable :: str

call matrix_random_init

l1  = random_sort(10)        ! random permutation of 1..10
l1  = random_sort(5, 15)     ! random permutation of 5..15
l1  = random_sort(l1)        ! reshuffle integer array l1
str = random_sort('Hello')   ! random anagram of 'Hello'
\end{lstlisting}

% ===================================================================
% 21. random_sort_array
% ===================================================================
\manentry{Generic Subroutine}{random\_sort\_array}{random\textunderscore sort\textunderscore array}{Randomly shuffles a 1D integer, real, or complex array in-place using the Fisher--Yates algorithm}{proc:random_sort_array}

\manheading{Synopsis}
\lstset{language=[95]Fortran}
\begin{lstlisting}
call random_sort_array(array)
\end{lstlisting}

\manheading{Brief}
Performs an in-place Fisher--Yates random permutation on a rank-1 \typename{integer}, \typename{real(8)}, or \typename{complex(8)} array.

\manheading{Description}
The generic subroutine \procname{random\_sort\_array} applies the same $O(N)$ Fisher--Yates shuffle as \procname{random\_sort}, but mutates the passed 1D array \argname{array} directly in-place without allocating a temporary output array.

\manheading{Dummy Arguments}
\begin{description}[leftmargin=2.2cm, style=nextline]
  \argitem{array}{integer, real(8), or complex(8), dimension(:), intent(inout)}{1D array whose elements are randomly permuted in-place.}
\end{description}

\manheading{Example}
\lstset{language=[95]Fortran}
\begin{lstlisting}
use array_mod
integer    :: l(5) = [1, 2, 3, 4, 5]
real   (8) :: v(5)
complex(8) :: w(5)

call random_number(v); call random_number(w)
call random_sort_array(l)   ! l is shuffled in-place
call random_sort_array(v)   ! v is shuffled in-place
call random_sort_array(w)   ! w is shuffled in-place
\end{lstlisting}

% ===================================================================
% 22. norm
% ===================================================================
\manentry{Pure Generic Function}{norm}{norm}{Computes the normalized $L_1$ norm (mean absolute value of elements) of a rank-1, rank-2, or rank-3 array}{proc:norm}

\manheading{Synopsis}
\lstset{language=[95]Fortran}
\begin{lstlisting}
r = norm(C)
\end{lstlisting}

\manheading{Brief}
Returns $\texttt{SUM(ABS(C)) / SIZE(C)}$ as a \typename{real(8)} scalar for any 1D, 2D, or 3D real or complex array \argname{C}.

\manheading{Description}
The pure generic function \procname{norm} computes the average element-wise absolute value (normalized $L_1$ norm) of an array \argname{C} of rank 1, 2, or 3:
\begin{align}
  \text{Rank 1:}\quad \texttt{norm}(v) &= \frac{1}{M}\sum_{i=1}^{M} |v_i|\, , \\
  \text{Rank 2:}\quad \texttt{norm}(A) &= \frac{1}{M N}\sum_{i=1}^{M}\sum_{j=1}^{N} |A_{ij}|\, , \\
  \text{Rank 3:}\quad \texttt{norm}(T) &= \frac{1}{M N K}\sum_{i=1}^{M}\sum_{j=1}^{N}\sum_{k=1}^{K} |T_{ijk}|\, ,
\end{align}
where $|z| = \sqrt{(\Re z)^2 + (\Im z)^2}$ for complex arrays. Normalizing by the total number of elements $\texttt{SIZE(C)}$ makes \procname{norm} scale-invariant with respect to array size, which is convenient for convergence tests and numerical residual checks across different matrix sizes $N$.

\manheading{Dummy Arguments}
\begin{description}[leftmargin=2.2cm, style=nextline]
  \argitem{C}{complex(8) or real(8), dimension(:), (:,:), or (:,:,:), intent(in)}{Input array of rank 1, 2, or 3.}
  \argitem{r}{real(8), scalar, return value}{The mean absolute element value $\texttt{SUM(ABS(C)) / SIZE(C)} \ge 0$.}
\end{description}

\manheading{Example}
\lstset{language=[95]Fortran}
\begin{lstlisting}
use array_mod
complex(8) :: t1(4,4,3), m1(4,4), w1(4)
real   (8) :: s1(4,4,3), d1(4,4), v1(4), r

call random_number(t1); call random_number(m1); call random_number(w1)
call random_number(s1); call random_number(d1); call random_number(v1)

r = norm(t1) ; r = norm(s1)   ! rank-3 arrays: sum_ijk |C_ijk| / (4*4*3)
r = norm(m1) ; r = norm(d1)   ! rank-2 arrays: sum_ij  |C_ij|  / (4*4)
r = norm(w1) ; r = norm(v1)   ! rank-1 arrays: sum_i   |C_i|   / 4
\end{lstlisting}

%   2.2 Matrix Multiplication
% ===================================================================
% Subsection 2.2: Matrix Multiplication
% ===================================================================
\subsection{Matrix Multiplication}
\label{subsec:matrix_multiplication}

\modname{array\_mod} provides five complementary interfaces for matrix--matrix and matrix--vector multiplication backed by Level-2 and Level-3 \texttt{BLAS} routines (\texttt{zgemm}, \texttt{dgemm}, \texttt{zhemm}, \texttt{dsymm}, \texttt{zherk}, \texttt{dsyrk}, \texttt{zgemv}, \texttt{dgemv}, \texttt{zhemv}, and \texttt{dsymv}):
\begin{enumerate}
  \item \procname{operator(.mm.)}: Infix operator overloading for matrix--matrix, matrix--vector, and vector--matrix products across all combinations of \typename{complex(8)} and \typename{real(8)} operands.
  \item \procname{lmatmul}: Functional interface supporting general products, Hermitian/symmetric matrix multiplication, rank-$k$ products ($A A^\dagger, A^\dagger A, A A^T, A^T A$), and transposed/conjugate-transposed matrix--vector products.
  \item \procname{mmmult}: High-performance subroutine interface for matrix--matrix and rank-$k$ products without temporary array allocation overhead.
  \item \procname{mvmult}: High-performance subroutine interface for right matrix--vector multiplication ($w = A v$, $w = A^T v$, $w = A^\dagger v$).
  \item \procname{vmmult}: High-performance subroutine interface for left vector--matrix multiplication ($w = v A$, i.e., $w_i = \sum_j v_j A_{ji}$).
\end{enumerate}

\begin{notebox}[Preservation of Custom Array Bounds]
In Fortran, when the target array on the left-hand side of an assignment (e.g., \texttt{m3 = lmatmul(m1, m2)} or \texttt{m3 = (m1 .mm. m2)}) is already allocated with conformable shape and custom lower/upper bounds \texttt{m3(is:ie, ks:ke)}, those bounds are preserved on assignment (unless compiler flags or shape mismatches trigger automatic reallocation, in which case the lower bounds reset to \texttt{1}). In the subroutine forms (\procname{mmmult}, \procname{mvmult}, \procname{vmmult}), the target array is passed as an \intentname{intent(inout)} argument and its bounds are always preserved.
\end{notebox}

% ===================================================================
% 23. operator(.mm.)
% ===================================================================
\manentry{Pure Generic Operator}{operator(.mm.)}{operator(.mm.)}{Infix operator for BLAS-accelerated matrix--matrix, matrix--vector, and vector--matrix multiplication}{proc:operator_mm}

\manheading{Synopsis}
\lstset{language=[95]Fortran}
\begin{lstlisting}
! Matrix-Matrix Multiplication:
C = (A .mm. B)

! Matrix-Vector and Vector-Matrix Multiplication:
w = (A .mm. v)      ! w_i = sum_j A_ij v_j
w = (v .mm. A)      ! w_i = sum_j v_j  A_ji = (A^T v)_i
\end{lstlisting}

\manheading{Brief}
Evaluates matrix--matrix ($A B$), matrix--vector ($A v$), or vector--matrix ($v^T A$) products for any combination of \typename{complex(8)} and \typename{real(8)} arrays by delegating to \procname{lmatmul}.

\manheading{Description}
The binary operator \procname{.mm.} is overloaded across 12 module procedures covering all combinations of \typename{complex(8)} and \typename{real(8)} rank-2 matrices and rank-1 vectors:
\begin{itemize}
  \item \textbf{Matrix--Matrix (\texttt{C1 .mm. C2})}: Calls \texttt{lmatmul(C1, C2)}. If one operand is \typename{real(8)} and the other is \typename{complex(8)}, the real array is promoted to a temporary \typename{complex(8)} array before calling \texttt{lmatmul}, and a \typename{complex(8)} matrix of shape $(\texttt{size(C1,1)}, \texttt{size(C2,2)})$ is returned.
  \item \textbf{Matrix--Vector (\texttt{C1 .mm. v2})}: Calls \texttt{lmatmul(A=C1, v=v2)} to compute $v_{3,i} = \sum_j (C_1)_{ij} v_{2,j}$ of length \texttt{size(C1,1)}. Mixed real/complex operands are promoted to \typename{complex(8)}.
  \item \textbf{Vector--Matrix (\texttt{v1 .mm. C2})}: Calls \texttt{lmatmul(v=v1, A=C2, type='T')} to compute $v_{3,i} = \sum_j v_{1,j} (C_2)_{ji}$ of length \texttt{size(C2,2)}. Mixed real/complex operands are promoted to \typename{complex(8)}.
\end{itemize}
If the inner dimensions do not match, all elements of the returned array are set to IEEE \procname{NaN}.

\begin{warnbox}[Operator Precedence and Specialized BLAS Calls]
\begin{enumerate}
  \item \textbf{Always enclose \texttt{.mm.} expressions in parentheses!} In the Fortran standard, user-defined operators like \texttt{.mm.} have the \emph{lowest} operator precedence---lower than addition (\texttt{+}), subtraction (\texttt{-}), and even logical operators. Thus, \lcodef{m1 .mm. m2 - m3} is parsed by the compiler as \lcodef{m1 .mm. (m2 - m3)}, \emph{not} \lcodef{(m1 .mm. m2) - m3}.
  \item The procedures overloaded by \texttt{.mm.} always call the general \texttt{BLAS} routines (\texttt{zgemm}/\texttt{dgemm}/\texttt{zgemv}/\texttt{dgemv}) and do not invoke the Hermitian or symmetric \texttt{BLAS} routines. Use \procname{lmatmul} or \procname{mmmult} when you want to pass \texttt{mtype='ZH'} or \texttt{mtype='DS'}.
\end{enumerate}
\end{warnbox}

\manheading{Dummy Arguments}
\begin{description}[leftmargin=2.2cm, style=nextline]
  \argitem{C1}{complex(8) or real(8), dimension(:,:) or dimension(:), intent(in)}{Left operand: an $M\times K$ matrix or a length-$M$ vector.}
  \argitem{C2}{complex(8) or real(8), dimension(:,:) or dimension(:), intent(in)}{Right operand: a $K\times N$ matrix (or $M\times N$ when \argname{C1} is a vector) or a length-$K$ vector.}
  \argitem{C3 / v3}{complex(8) or real(8), allocatable rank-2 or rank-1 array, return value}{Resulting matrix or vector product (\typename{real(8)} if both operands are real, \typename{complex(8)} otherwise).}
\end{description}

\manheading{Example}
\lstset{language=[95]Fortran}
\begin{lstlisting}
use array_mod
complex(8) :: m1(4,6), m2(6,5), m3(4,5), v1(6), v2(4)
real   (8) :: d1(4,6), d2(6,5), d3(4,5), u1(6), u2(4)

call random_number(m1); call random_number(m2); call random_number(v1); call random_number(v2)
call random_number(d1); call random_number(d2); call random_number(u1); call random_number(u2)

! Matrix-matrix products (pure and mixed real/complex):
m3 = (m1 .mm. m2)
m3 = (m1 .mm. d2)
m3 = (d1 .mm. m2)
d3 = (d1 .mm. d2)

! Matrix-vector and vector-matrix products:
v2 = (m1 .mm. v1) ; v1 = (v2 .mm. m1)
v2 = (m1 .mm. u1) ; v1 = (u2 .mm. m1)
v2 = (d1 .mm. v1) ; v1 = (v2 .mm. d1)
u2 = (d1 .mm. u1) ; u1 = (u2 .mm. d1)
\end{lstlisting}

% ===================================================================
% 24. lmatmul
% ===================================================================
\manentry{Pure Generic Function}{lmatmul}{lmatmul}{BLAS-backed matrix--matrix, rank-$k$, and matrix--vector multiplication function}{proc:lmatmul}

\manheading{Synopsis}
\lstset{language=[95]Fortran}
\begin{lstlisting}
! 1. Two-Matrix Multiplication (A . B):
C = lmatmul(A, B, [mtype])

! 2. Single-Matrix Rank-k Multiplication (A . A^\dagger, A^\dagger . A, A . A^T, A^T . A):
C = lmatmul(A, [side])

! 3. Matrix-Vector Multiplication (A . v, A^T . v, A^\dagger . v):
w = lmatmul(A, v, [type])
w = lmatmul(v=v, A=A, [type])
\end{lstlisting}

\manheading{Brief}
Computes matrix--matrix products, Hermitian/symmetric rank-$k$ products, or matrix--vector products using \texttt{BLAS} routines \texttt{zgemm}, \texttt{dgemm}, \texttt{zhemm}, \texttt{dsymm}, \texttt{zherk}, \texttt{dsyrk}, \texttt{zgemv}, \texttt{dgemv}, \texttt{zhemv}, and \texttt{dsymv}.

\manheading{Description}
The generic function \procname{lmatmul} provides three distinct overloading families for contiguous \typename{complex(8)} or \typename{real(8)} arrays:
\begin{enumerate}
  \item \textbf{Matrix--Matrix Multiplication (\texttt{C = lmatmul(A, B, [mtype])})}:
  Given an $M\times K$ matrix \argname{A} and a $K'\times N$ matrix \argname{B} of the same type (\typename{complex(8)} or \typename{real(8)}), allocates and returns the $M\times N$ product $C = A B$:
  \begin{itemize}
    \item The matrix type is read from \argname{mtype} case-insensitively: a one-character string is the type itself (\texttt{'H'}, \texttt{'S'}, \texttt{'G'}), a longer string supplies it in its second character (\texttt{'ZH'}, \texttt{'DS'}, \texttt{'ZG'}, \texttt{'DG'}).
    \item By default (\argname{mtype} absent, or a type other than \texttt{'H'}/\texttt{'S'}), calls \texttt{zgemm} (complex) or \texttt{dgemm} (real).
    \item If \texttt{mtype='ZH'} (or \texttt{'H'}) for complex matrices, \argname{A} is assumed to be an $M\times M$ Hermitian matrix and \texttt{zhemm('L','U',...)} is called (referencing the upper triangle of \argname{A}; \argname{B} does not need to be Hermitian).
    \item If \texttt{mtype='DS'} (or \texttt{'S'}) for real matrices, \argname{A} is assumed to be an $M\times M$ symmetric matrix and \texttt{dsymm('L','U',...)} is called (referencing the upper triangle of \argname{A}).
    \item If $K \neq K'$ (or if $M \neq K$ when \texttt{'H'}/\texttt{'S'} is requested), the returned $M\times N$ array is filled with IEEE \procname{NaN}.
  \end{itemize}

  \item \textbf{Single-Matrix Rank-$k$ Product (\texttt{C = lmatmul(A, [side])})}:
  Given an $M\times N$ matrix \argname{A}, uses the specialized rank-$k$ update \texttt{BLAS} routines \texttt{zherk} (complex) or \texttt{dsyrk} (real) to compute the upper triangle of the Hermitian/symmetric product in roughly half the operations of a full matrix multiply, and then mirrors the upper triangle into the lower triangle:
  \begin{itemize}
    \item If \argname{side} is omitted, blank, or \texttt{side='N'}: returns the $M\times M$ matrix $C = A A^\dagger$ (complex) or $C = A A^T$ (real).
    \item If \texttt{side(1:1)} $\in \{\texttt{'C','c','T','t','H','h','S','s','L','l'}\}$ (e.g., \texttt{side='H'} or \texttt{side='S'}): returns the $N\times N$ matrix $C = A^\dagger A$ (complex) or $C = A^T A$ (real).
    \item Any other value of \texttt{side(1:1)} returns an $M\times M$ array filled with IEEE \procname{NaN}.
  \end{itemize}

  \item \textbf{Matrix--Vector Multiplication (\texttt{w = lmatmul(A, v, [type])})}:
  Given an $M\times N$ matrix \argname{A} and a 1D vector \argname{v} of the same type, computes (all \argname{type} selectors case-insensitive):
  \begin{itemize}
    \item \texttt{type='N'} (or \texttt{'n'}, default): $w = A v$, i.e., $w_i = \sum_{j=1}^{N} A_{ij} v_j$ of length $M$ (requires $\texttt{size(v)} == N$), via \texttt{zgemv('N',...)} or \texttt{dgemv('N',...)}.
    \item \texttt{type='T'} (or \texttt{'t'}): $w = A^T v = v^T A$, i.e., $w_i = \sum_{j=1}^{M} v_j A_{ji}$ of length $N$ (requires $\texttt{size(v)} == M$), via \texttt{zgemv('T',...)} or \texttt{dgemv('T',...)}. Note that using Fortran keyword arguments \lcodef{lmatmul(v=v2, A=m1, type='T')} provides a natural left-multiplication syntax.
    \item \texttt{type='C'} (or \texttt{'c'}): For complex arrays, $w = A^\dagger v$, i.e., $w_i = \sum_{j=1}^{M} A_{ji}^* v_j$ of length $N$ (requires $\texttt{size(v)} == M$), via \texttt{zgemv('C',...)}. (For real arrays, \texttt{'C'}/\texttt{'c'} is equivalent to \texttt{'T'}.)
    \item \texttt{type='H'}, \texttt{'h'}, \texttt{'S'}, or \texttt{'s'}: For complex arrays, \texttt{type='H'} or \texttt{'h'} assumes \argname{A} is an $M\times M$ Hermitian matrix and computes $w = A v$ via \texttt{zhemv('U',...)} (referencing only the upper triangle of \argname{A}). For real arrays, \texttt{type='S'}, \texttt{'s'}, \texttt{'H'}, or \texttt{'h'} assumes \argname{A} is an $M\times M$ real symmetric matrix and computes $w = A v$ via \texttt{dsymv('U',...)} (referencing only the upper triangle of \argname{A}, requiring $M = N = \texttt{size(v)}$).
    \item If \argname{A} and \argname{v} are non-conformable, if \argname{A} is non-square when \texttt{'H'}/\texttt{'S'} is selected, or if \argname{type} is unrecognized (including an empty \texttt{type=''}), the returned vector \argname{w} is filled with IEEE \procname{NaN}. Conformable 0-sized inputs (real or complex) return a 0-sized or zero vector safely without calling \texttt{BLAS}.
  \end{itemize}
\end{enumerate}

\manheading{Dummy Arguments}
\begin{description}[leftmargin=2.2cm, style=nextline]
  \argitem{A}{complex(8) or real(8), dimension(:,:), contiguous, intent(in)}{Input matrix ($M\times K$ in two-matrix form; $M\times N$ in single-matrix and matrix--vector forms).}
  \argitem{B}{complex(8) or real(8), dimension(:,:), contiguous, intent(in)}{Right-hand $K\times N$ input matrix in the two-matrix form.}
  \argitem{v}{complex(8) or real(8), dimension(:), contiguous, intent(in)}{Input vector in the matrix--vector form.}
  \argitem{mtype}{character(*), optional, intent(in)}{In \texttt{lmatmul(A,B,mtype)}: \texttt{'ZH'} (or \texttt{'H'}) if complex \argname{A} is Hermitian, \texttt{'DS'} (or \texttt{'S'}) if real \argname{A} is symmetric, or \texttt{'ZG'}/\texttt{'DG'} (default) for general matrices.}
  \argitem{side}{character(*), optional, intent(in)}{In \texttt{lmatmul(A,side)}: \texttt{'N'} (default) for $A A^\dagger$ ($A A^T$), or \texttt{'H'}/\texttt{'S'}/\texttt{'C'}/\texttt{'T'}/\texttt{'L'} for $A^\dagger A$ ($A^T A$); any other character gives a \procname{NaN}-filled result.}
  \argitem{type}{character(*), optional, intent(in)}{In \texttt{lmatmul(A,v,type)} (case-insensitive): \texttt{'N'} (default) for $A v$, \texttt{'T'} for $A^T v$, \texttt{'C'} for $A^\dagger v$ (or $A^T v$ for real), \texttt{'H'} for Hermitian $A v$ (complex \texttt{zhemv}) or symmetric $A v$ (real \texttt{dsymv}), and \texttt{'S'} for real symmetric $A v$ (\texttt{dsymv}).}
  \argitem{C / w}{complex(8) or real(8), allocatable rank-2 or rank-1 array, return value}{The resulting product matrix \argname{C} or vector \argname{w} (or \procname{NaN} on dimension mismatch).}
\end{description}

\manheading{Example}
\lstset{language=[95]Fortran}
\begin{lstlisting}
use array_mod
complex(8) :: m1(4,6), m2(6,5), m3(4,5), m4(4,4), m5(6,6), m6(4,6), mh(4,4)
complex(8) :: v1(6), v2(4), v3(4), v4(6)
real   (8) :: d1(4,6), d2(6,5), d3(4,5), d4(4,4), d5(6,6), d6(4,6), ds(4,4)
real   (8) :: u1(6), u2(4), u3(4), u4(6)

call random_number(m1); call random_number(m2); call random_number(mh); call hermitian_set(mh)
call random_number(d1); call random_number(d2); call random_number(ds); call symmetric_set(ds)
call random_number(v1); call random_number(v2); call random_number(u1); call random_number(u2)

! Matrix-matrix and rank-k products:
m3 = lmatmul(m1, m2)
m6 = lmatmul(mh, m1, mtype='ZH')   ! mh is 4x4 Hermitian, m1 is 4x6 general: m6 is 4x6
m4 = lmatmul(m1)                   ! m4 = m1 . m1^\dagger (4x4)
m5 = lmatmul(m1, side='H')         ! m5 = m1^\dagger . m1 (6x6)

d3 = lmatmul(d1, d2)
d6 = lmatmul(ds, d1, mtype='DS')   ! ds is 4x4 symmetric, d1 is 4x6 general: d6 is 4x6
d4 = lmatmul(d1)                   ! d4 = d1 . d1^T (4x4)
d5 = lmatmul(d1, side='S')         ! d5 = d1^T . d1 (6x6)

! Matrix-vector products:
v3 = lmatmul(m1, v1)               ! v3 = m1 . v1
v4 = lmatmul(v=v2, A=m1, type='T') ! v4 = m1^T . v2
v4 = lmatmul(m1, v2, type='C')     ! v4 = m1^\dagger . v2
v3 = lmatmul(mh, v2, type='H')     ! v3 = mh . v2 (mh assumed Hermitian)

u3 = lmatmul(d1, u1)               ! u3 = d1 . u1
u4 = lmatmul(v=u2, A=d1, type='T') ! u4 = d1^T . u2
u3 = lmatmul(ds, u2, type='S')     ! u3 = ds . u2 (ds assumed symmetric)
\end{lstlisting}

% ===================================================================
% 25. mmmult
% ===================================================================
\manentry{Pure Generic Subroutine}{mmmult}{mmmult}{Zero-allocation BLAS subroutine for matrix--matrix and rank-$k$ multiplication}{proc:mmmult}

\manheading{Synopsis}
\lstset{language=[95]Fortran}
\begin{lstlisting}
! 1. Two-Matrix Multiplication (C = A . B):
call mmmult(A, B, C, [mtype])

! 2. Single-Matrix Rank-k Multiplication (C = A . A^\dagger, A^\dagger . A, A . A^T, A^T . A):
call mmmult(A, C, [side])
\end{lstlisting}

\manheading{Brief}
Subroutine counterpart of \procname{lmatmul} that writes the matrix product $C = A B$ or rank-$k$ product $C = A A^\dagger, A^\dagger A, A A^T, A^T A$ directly into a pre-allocated output array \argname{C}.

\manheading{Description}
For performance-critical inner loops, \procname{mmmult} avoids dynamic memory allocation of the result by writing the result directly into the caller-supplied \intentname{intent(inout)} array \argname{C}. As shown in the benchmarks in Section~\ref{sec:performance}, \procname{mmmult} achieves the exact same execution speed as raw Fortran~77 calls to \texttt{zgemm}/\texttt{dgemm}.
\begin{itemize}
  \item \textbf{Three-matrix call \texttt{call mmmult(A, B, C, [mtype])}}: Computes $C = A B$ for $M\times K$ matrix \argname{A}, $K\times N$ matrix \argname{B}, and $M\times N$ matrix \argname{C} using \texttt{zgemm}/\texttt{dgemm} (or \texttt{zhemm}/\texttt{dsymm} when \texttt{mtype='ZH'} or \texttt{'DS'}).
  \item \textbf{Two-matrix call \texttt{call mmmult(A, C, [side])}}: Computes $C = A A^\dagger$ (complex) or $C = A A^T$ (real) of size $M\times M$ by default, or $C = A^\dagger A$ / $C = A^T A$ of size $N\times N$ when \texttt{side(1:1)} $\in \{\texttt{'C','T','H','S','L'}\}$ (case-insensitive), using \texttt{zherk} or \texttt{dsyrk}. Any other value of \texttt{side(1:1)} fills \argname{C} with IEEE \procname{NaN}.
\end{itemize}
\argname{C} does not need to be initialized: the \texttt{BLAS} routines are called with $\beta=0$, so the previous contents of \argname{C} (even \procname{NaN}) are never read, and \argname{C} is set to zero explicitly only when the inner dimension is $0$. If the dimensions of \argname{A}, \argname{B}, and \argname{C} are not conformable, \argname{C} is filled with IEEE \procname{NaN} and the subroutine returns immediately.

\manheading{Dummy Arguments}
\begin{description}[leftmargin=2.2cm, style=nextline]
  \argitem{A}{complex(8) or real(8), dimension(:,:), contiguous, intent(in)}{Input matrix ($M\times K$ in 3-matrix form; $M\times N$ in 2-matrix form).}
  \argitem{B}{complex(8) or real(8), dimension(:,:), contiguous, intent(in)}{Second input matrix ($K\times N$) in the 3-matrix form.}
  \argitem{C}{complex(8) or real(8), dimension(:,:), contiguous, intent(inout)}{Pre-allocated output matrix ($M\times N$ in 3-matrix form; $M\times M$ or $N\times N$ in 2-matrix form).}
  \argitem{mtype}{character(*), optional, intent(in)}{Symmetry hint for \argname{A} in the 3-matrix form: \texttt{'ZH'} (Hermitian) or \texttt{'DS'} (real symmetric).}
  \argitem{side}{character(*), optional, intent(in)}{Orientation selector in the 2-matrix form: \texttt{'N'} (default) for $A A^\dagger$ ($A A^T$), or \texttt{'H'}/\texttt{'S'}/\texttt{'C'}/\texttt{'T'}/\texttt{'L'} for $A^\dagger A$ ($A^T A$); any other character fills \argname{C} with \procname{NaN}.}
\end{description}

\begin{warnbox}[Output Must Not Alias an Input]
In \procname{mmmult}, \procname{mvmult}, and \procname{vmmult} the output argument must be a different array from the inputs. Calls such as \lcodef{call mmmult(A, B, A)} or \lcodef{call mvmult(A, v, v)} are invalid Fortran (an \intentname{intent(inout)} argument aliasing an \intentname{intent(in)} argument) and give wrong results; use a temporary array instead.
\end{warnbox}

\manheading{Example}
\lstset{language=[95]Fortran}
\begin{lstlisting}
use array_mod
complex(8) :: m1(8,8), m2(8,8), m3(8,8), m4(8,8), m5(8,8)
real   (8) :: d1(8,8), d2(8,8), d3(8,8), d4(8,8), d5(8,8)

call random_number(m1); call random_number(m2)
call random_number(d1); call random_number(d2)

call mmmult(m1, m2, m3)              ! m3 = m1 . m2
call mmmult(m1, m4)                  ! m4 = m1 . m1^\dagger
call mmmult(m1, m5, side='H')        ! m5 = m1^\dagger . m1
call mmmult(m4, m2, m3, mtype='ZH')  ! m4 is Hermitian

call mmmult(d1, d2, d3)              ! d3 = d1 . d2
call mmmult(d1, d4)                  ! d4 = d1 . d1^T
call mmmult(d1, d5, side='S')        ! d5 = d1^T . d1
call mmmult(d4, d2, d3, mtype='DS')  ! d4 is Symmetric
\end{lstlisting}

% ===================================================================
% 26. mvmult
% ===================================================================
\manentry{Pure Generic Subroutine}{mvmult}{mvmult}{Zero-allocation BLAS subroutine for matrix--vector multiplication}{proc:mvmult}

\manheading{Synopsis}
\lstset{language=[95]Fortran}
\begin{lstlisting}
call mvmult(A, v, w, [type])
\end{lstlisting}

\manheading{Brief}
Computes the matrix--vector product $w = A v$, $w = A^T v$, or $w = A^\dagger v$ into a pre-allocated output vector \argname{w} using \texttt{BLAS} \texttt{zgemv}, \texttt{zhemv}, \texttt{dgemv}, or \texttt{dsymv}.

\manheading{Description}
The generic subroutine \procname{mvmult} multiplies an $M\times N$ matrix \argname{A} by a 1D vector \argname{v} and stores the result in the pre-allocated 1D vector \argname{w} (both \typename{complex(8)} or both \typename{real(8)}; all \argname{type} selectors case-insensitive):
\begin{itemize}
  \item \texttt{type='N'} (or \texttt{'n'}, default): $w_i = \sum_{j=1}^{N} A_{ij} v_j$ ($\texttt{size(v)}=N$, $\texttt{size(w)}=M$).
  \item \texttt{type='T'} (or \texttt{'t'}): $w_i = \sum_{j=1}^{M} v_j A_{ji}$ ($\texttt{size(v)}=M$, $\texttt{size(w)}=N$).
  \item \texttt{type='C'} (or \texttt{'c'}): $w_i = \sum_{j=1}^{M} v_j A_{ji}^*$ for complex arrays, or $w_i = \sum_{j=1}^{M} v_j A_{ji}$ for real arrays ($\texttt{size(v)}=M$, $\texttt{size(w)}=N$).
  \item \texttt{type='H'}, \texttt{'h'}, \texttt{'S'}, or \texttt{'s'}: $w_i = \sum_{j=1}^{M} A_{ij} v_j$ requiring $M = N = \texttt{size(v)} = \texttt{size(w)}$ and referencing only the upper triangle of \argname{A}: \texttt{type='H'}/\texttt{'h'} calls \texttt{zhemv('U',...)} assuming Hermitian \argname{A} for complex arrays, while \texttt{type='S'}/\texttt{'s'}/\texttt{'H'}/\texttt{'h'} calls \texttt{dsymv('U',...)} assuming symmetric \argname{A} for real arrays.
\end{itemize}
If any array dimension is non-conformable (including non-square \argname{A} for \texttt{'H'}/\texttt{'S'}) or \argname{type} is unrecognized (including an empty \texttt{type=''}), \argname{w} is filled with IEEE \procname{NaN}. Conformable 0-sized inputs (real or complex) return safely without calling \texttt{BLAS}; when the inner dimension is $0$, \argname{w} is set to zero.

\manheading{Dummy Arguments}
\begin{description}[leftmargin=2.2cm, style=nextline]
  \argitem{A}{complex(8) or real(8), dimension(:,:), contiguous, intent(in)}{Input $M\times N$ matrix.}
  \argitem{v}{complex(8) or real(8), dimension(:), contiguous, intent(in)}{Input vector of length $N$ (for \texttt{'N'}, \texttt{'H'}, \texttt{'S'}) or $M$ (for \texttt{'T'}, \texttt{'C'}).}
  \argitem{w}{complex(8) or real(8), dimension(:), contiguous, intent(inout)}{Pre-allocated output vector of length $M$ (for \texttt{'N'}, \texttt{'H'}, \texttt{'S'}) or $N$ (for \texttt{'T'}, \texttt{'C'}).}
  \argitem{type}{character(*), optional, intent(in)}{Case-insensitive operation selector: \texttt{'N'} (default), \texttt{'T'}, \texttt{'C'}, \texttt{'H'} (Hermitian for complex, symmetric for real), or \texttt{'S'} (symmetric for real).}
\end{description}

\manheading{Example}
\lstset{language=[95]Fortran}
\begin{lstlisting}
use array_mod
complex(8) :: m1(8,8), v1(8), v3(8)
real   (8) :: ds(8,8), u1(8), u3(8)

call random_number(m1); call random_number(v1)
call random_number(ds); call random_number(u1)

call mvmult(m1, v1, v3)            ! v3_i = m1_ij v1_j
call mvmult(m1, v1, v3, type='T')  ! v3_j = v1_i  m1_ij
call mvmult(m1, v1, v3, type='C')  ! v3_j = v1_i (m1_ij)*

call hermitian_set(m1)
call mvmult(m1, v1, v3, type='H')  ! v3_i = m1_ij v1_j (m1 assumed Hermitian)

call symmetric_set(ds)
call mvmult(ds, u1, u3, type='S')  ! u3_i = ds_ij u1_j (ds assumed symmetric)
\end{lstlisting}

% ===================================================================
% 27. vmmult
% ===================================================================
\manentry{Pure Generic Subroutine}{vmmult}{vmmult}{Zero-allocation BLAS subroutine for left vector--matrix multiplication}{proc:vmmult}

\manheading{Synopsis}
\lstset{language=[95]Fortran}
\begin{lstlisting}
call vmmult(v, A, w)
\end{lstlisting}

\manheading{Brief}
Computes the left vector--matrix product $w = v^T A$ ($w_i = \sum_{j=1}^{M} v_j A_{ji}$) into a pre-allocated output vector \argname{w} using \texttt{BLAS} \texttt{zgemv('T',...)} or \texttt{dgemv('T',...)}.

\manheading{Description}
The generic subroutine \procname{vmmult} multiplies a length-$M$ row vector \argname{v} from the left with an $M\times N$ matrix \argname{A} and writes the resulting length-$N$ vector into \argname{w} (all \typename{complex(8)} or all \typename{real(8)}):
\begin{align}
  w_i = \sum_{j=1}^{M} v_j\, A_{ji} = (A^T v)_i\, , \qquad i = 1,\ldots,N\, .
\end{align}
Internally, it invokes \texttt{zgemv('T', M, N, ...)} or \texttt{dgemv('T', M, N, ...)}. Because Fortran stores arrays in column-major order, computing $w_i = \sum_j v_j A_{ji}$ via \texttt{gemv('T')} accesses each column $A_{:,i}$ contiguously in memory and computes dot products along contiguous memory addresses. On multi-threaded supercomputer architectures (such as ARIS, see Table~\ref{tab:b7}), \procname{vmmult} can be up to $2\times$ faster than \procname{mvmult} for large $N \ge 250$.

If $\texttt{size(v)} \neq M$ or $\texttt{size(w)} \neq N$, \argname{w} is filled with IEEE \procname{NaN}.

\manheading{Dummy Arguments}
\begin{description}[leftmargin=2.2cm, style=nextline]
  \argitem{v}{complex(8) or real(8), dimension(:), contiguous, intent(in)}{Input left-multiplying vector of length $M = \texttt{size(A,1)}$.}
  \argitem{A}{complex(8) or real(8), dimension(:,:), contiguous, intent(in)}{Input $M\times N$ matrix.}
  \argitem{w}{complex(8) or real(8), dimension(:), contiguous, intent(inout)}{Pre-allocated output vector of length $N = \texttt{size(A,2)}$.}
\end{description}

\manheading{Example}
\lstset{language=[95]Fortran}
\begin{lstlisting}
use array_mod
complex(8) :: m1(4,6), v2(4), v4(6)
real   (8) :: d1(4,6), u2(4), u4(6)

call random_number(m1); call random_number(v2)
call random_number(d1); call random_number(u2)

call vmmult(v2, m1, v4)   ! v4_i = sum_j v2_j m1_ji
call vmmult(u2, d1, u4)   ! u4_i = sum_j u2_j d1_ji
\end{lstlisting}

%   2.3 Linear Algebra (Inverse, determinant, Pfaffian; Eigenvalues and eigenvectors)
% ===================================================================
% Subsection 2.3: Linear Algebra
% ===================================================================
\subsection{Linear Algebra}
\label{subsec:linear_algebra}

This subsection documents the high-level linear algebra solvers of \modname{array\_mod} for computing matrix inverses, determinants, Pfaffians, eigenvalues, and right eigenvectors. Under the hood, these procedures query optimal \texttt{LAPACK} workspace sizes dynamically (\texttt{LWORK = -1}) and invoke \texttt{LAPACK} and \texttt{PFAPACK} routines.

% -------------------------------------------------------------------
\subsubsection{Inverse, determinant, Pfaffian}
\label{subsubsec:inv_det_pf}

% ===================================================================
% 28. inverse
% ===================================================================
\manentry{Generic Function}{inverse}{inverse}{Computes the matrix inverse $C^{-1}$ of a real or complex square matrix via LU factorization}{proc:inverse}

\manheading{Synopsis}
\lstset{language=[95]Fortran}
\begin{lstlisting}
CI = inverse(C)
\end{lstlisting}

\manheading{Brief}
Computes and returns the inverse $C^{-1}$ of a non-singular square \typename{complex(8)} or \typename{real(8)} matrix \argname{C} such that $C^{-1} C = C C^{-1} = I$.

\manheading{Description}
Given a contiguous $N\times N$ \typename{complex(8)} or \typename{real(8)} matrix \argname{C}, \procname{inverse} computes its inverse \argname{CI} in two steps using \texttt{LAPACK}:
\begin{enumerate}
  \item Performs an $LU$ decomposition with partial pivoting, $C = P L U$, by calling \texttt{zgetrf} (complex) or \texttt{dgetrf} (real).
  \item Queries the optimal workspace size (\texttt{LWORK = -1}) and inverts the $LU$-factored matrix by calling \texttt{zgetri} (complex) or \texttt{dgetri} (real).
\end{enumerate}
If \argname{C} is not square ($\texttt{size(C,1)} \neq \texttt{size(C,2)}$), \procname{inverse} returns a \texttt{size(C,1)}$\times$\texttt{size(C,2)} array filled with IEEE \procname{NaN}. If \argname{C} is singular (\texttt{zgetrf}/\texttt{dgetrf} or \texttt{zgetri}/\texttt{dgetri} returns \texttt{INFO > 0}), the program does \emph{not} abort: the returned $N\times N$ array is filled with \procname{NaN}, which can be detected with \lcodef{isNaN(CI)}. Only an illegal-argument failure (\texttt{INFO < 0}) or a failed workspace query prints an error message to \texttt{stderr} and terminates execution via \texttt{matrix\_error} with a non-zero exit status (\texttt{error stop 1}).

\manheading{Dummy Arguments}
\begin{description}[leftmargin=2.2cm, style=nextline]
  \argitem{C}{complex(8) or real(8), dimension(:,:), contiguous, intent(in)}{Input $N\times N$ non-singular square matrix.}
  \argitem{CI}{complex(8) or real(8), dimension(:,:), allocatable, return value}{Allocated $N\times N$ inverse matrix $C^{-1}$ of the same type as \argname{C} (filled with \procname{NaN} if \argname{C} is singular or non-square).}
\end{description}

\manheading{Example}
\lstset{language=[95]Fortran}
\begin{lstlisting}
use array_mod
complex(8), allocatable :: m1(:,:), m2(:,:)
real   (8), allocatable :: d1(:,:), d2(:,:)

allocate(m1(8,8), d1(8,8))
call random_number(m1, sigma=1.0_8)
call random_number(d1, sigma=1.0_8)

m2 = inverse(m1)   ! lmatmul(m1, m2) == identitymatrix(8)
d2 = inverse(d1)   ! lmatmul(d1, d2) == didentitymatrix(8)
if (isNaN(d2)) print *, 'd1 is singular'
\end{lstlisting}

% ===================================================================
% 29. determinant
% ===================================================================
\manentry{Generic Function}{determinant}{determinant}{Computes the determinant $\det C$ of a real or complex square matrix via LU factorization}{proc:determinant}

\manheading{Synopsis}
\lstset{language=[95]Fortran}
\begin{lstlisting}
det = determinant(C)
\end{lstlisting}

\manheading{Brief}
Computes the determinant $\det C$ of an $N\times N$ \typename{complex(8)} or \typename{real(8)} matrix \argname{C} using \texttt{LAPACK}'s \texttt{zgetrf} or \texttt{dgetrf}.

\manheading{Description}
The generic function \procname{determinant} copies the input $N\times N$ matrix \argname{C} into a temporary workspace and computes its $LU$ factorization with partial pivoting, $C = P L U$, via \texttt{zgetrf} (complex) or \texttt{dgetrf} (real). Because $L$ has unit diagonal elements, the determinant is evaluated as the product of the diagonal entries of $U$ multiplied by the sign of the row permutation $P$:
\begin{align}
  \det C = (-1)^{S} \prod_{i=1}^{N} U_{ii}\, ,
\end{align}
where $S$ is the number of pivot indices satisfying $\texttt{ipiv}(i) \neq i$.
\begin{itemize}
  \item For singular matrices (\texttt{INFO > 0} returned by \texttt{zgetrf}/\texttt{dgetrf}), \procname{determinant} does \emph{not} abort; it returns $0$ (\texttt{(0.0\_8, 0.0\_8)} for complex \argname{C}, \texttt{0.0\_8} for real \argname{C}).
  \item For an empty $0\times 0$ matrix ($N = 0$), \procname{determinant} returns $1$ (\texttt{(1.0\_8, 0.0\_8)} for complex \argname{C}, \texttt{1.0\_8} for real \argname{C}).
  \item For non-square input ($\texttt{size(C,1)} \neq \texttt{size(C,2)}$), \procname{determinant} returns \procname{NaN} (\texttt{(NaN, NaN)} for complex \argname{C}).
  \item Only an illegal-argument failure from \texttt{LAPACK} (\texttt{INFO < 0}) terminates execution via \texttt{matrix\_error}.
\end{itemize}

\begin{warnbox}[Floating-Point Overflow Limit for Large Matrices]
Because $\det C$ grows or decays exponentially with matrix size $N$, \procname{determinant} is suitable only for moderately sized matrices. For example, for an $N\times N$ matrix filled with Gaussian random numbers with $\sigma = 1$, IEEE double-precision overflow ($|\det C| > 10^{308}$) occurs above $N \sim 192$. For larger matrices, always use \procname{lndet}.
\end{warnbox}

\manheading{Dummy Arguments}
\begin{description}[leftmargin=2.2cm, style=nextline]
  \argitem{C}{complex(8) or real(8), dimension(:,:), contiguous, intent(in)}{Input $N\times N$ square matrix.}
  \argitem{det}{complex(8) or real(8), scalar, return value}{The determinant $\det C$, matching the data type of \argname{C} ($0$ for singular matrices, $1$ for $N = 0$).}
\end{description}

\manheading{Example}
\lstset{language=[95]Fortran}
\begin{lstlisting}
use array_mod
complex(8) :: m1(16,16), zdet
real   (8) :: d1(16,16), rdet

call random_number(m1, sigma=1.0_8)
call random_number(d1, sigma=1.0_8)

zdet = determinant(m1)
rdet = determinant(d1)
\end{lstlisting}

% ===================================================================
% 30. lndet
% ===================================================================
\manentry{Generic Function}{lndet}{lndet}{Computes the logarithm of the modulus and the phase/sign of $\det C$ for large matrices}{proc:lndet}

\manheading{Synopsis}
\lstset{language=[95]Fortran}
\begin{lstlisting}
ldet = lndet(C)
! Reconstruct determinant: det = exp(ldet(1)) * ldet(2)
\end{lstlisting}

\manheading{Brief}
Returns a 2-element array $\texttt{ldet(1:2)}$ containing $\ln|\det C|$ in \texttt{ldet(1)} and the complex phase $\e{i\arg(\det C)}$ (or real sign $\pm 1.0$) in \texttt{ldet(2)} (or $\texttt{ldet(1)} = -\infty$ and $\texttt{ldet(2)} = 0$ when $\det C = 0$).

\manheading{Description}
To prevent floating-point overflow or underflow when $N \gtrsim 192$, \procname{lndet} performs the $LU$ decomposition $C = P L U$ via \texttt{zgetrf} or \texttt{dgetrf} and accumulates the logarithm of the absolute values and the unit-modulus phase of the diagonal elements $U_{ii}$ separately:
\begin{align}
  \texttt{ldet(1)} &= \ln|\det C| = \sum_{i=1}^{N} \ln|U_{ii}|\, , \\
  \texttt{ldet(2)} &= \e{i\arg(\det C)} = (-1)^{S} \prod_{i=1}^{N} \frac{U_{ii}}{|U_{ii}|}\, .
\end{align}
The full determinant is related to \texttt{ldet} by:
\begin{align}
  \det C = \exp(\texttt{ldet(1)}) \times \texttt{ldet(2)}\, .
\end{align}
\begin{itemize}
  \item For a non-singular \typename{complex(8)} matrix \argname{C}, \procname{lndet} returns a \typename{complex(8)} array \texttt{ldet(2)} where \texttt{ldet(1)} has zero imaginary part ($\Re(\texttt{ldet(1)}) = \ln|\det C|$) and \texttt{ldet(2)} is the unit-modulus complex phase $\e{i\theta}$.
  \item For a non-singular \typename{real(8)} matrix \argname{C}, \procname{lndet} returns a \typename{real(8)} array \texttt{ldet(2)} where \texttt{ldet(1)} $= \ln|\det C|$ and \texttt{ldet(2)} $= \pm 1.0$ (\texttt{+/-1.0\_8}).
  \item If \argname{C} is singular (\texttt{INFO > 0} from \texttt{zgetrf}/\texttt{dgetrf}, or any $|U_{ii}| = 0$ so that $\det C = 0$), \procname{lndet} does \emph{not} abort; it returns \texttt{ldet(1) = (}$-\infty$\texttt{, 0.0\_8)} (complex) or $-\infty$ (real) via \texttt{IEEE\_VALUE(0.0\_dp, IEEE\_NEGATIVE\_INF)}, and \texttt{ldet(2) = (0.0\_8, 0.0\_8)} (complex) or \texttt{0.0\_8} (real), so that $\Re(\texttt{ldet(1)}) = \ln|0| = -\infty$ and $\exp(\texttt{ldet(1)}) \times \texttt{ldet(2)} = 0$ unambiguously.
  \item For an empty $0\times 0$ matrix ($N = 0$), \procname{lndet} returns $\texttt{ldet(1)} = 0$ and $\texttt{ldet(2)} = 1$.
  \item For non-square input, both elements \texttt{ldet(1)} and \texttt{ldet(2)} are set to \procname{NaN}.
  \item Execution terminates via \texttt{matrix\_error} only if \texttt{zgetrf}/\texttt{dgetrf} reports an illegal argument (\texttt{INFO < 0}).
\end{itemize}

\manheading{Dummy Arguments}
\begin{description}[leftmargin=2.2cm, style=nextline]
  \argitem{C}{complex(8) or real(8), dimension(:,:), contiguous, intent(in)}{Input $N\times N$ square matrix.}
  \argitem{det}{complex(8) or real(8), dimension(2), return value}{Two-element array: \texttt{det(1)} $= \ln|\det C|$ and \texttt{det(2)} $= \e{i\arg(\det C)}$ (or $\pm 1.0$ for real \argname{C}), or \texttt{det(1)} $= -\infty$ and \texttt{det(2)} $= 0$ for singular \argname{C} (so that $\exp(\texttt{det(1)}) \times \texttt{det(2)} = 0$).}
\end{description}

\manheading{Example}
\lstset{language=[95]Fortran}
\begin{lstlisting}
use array_mod
complex(8) :: m1(32,32), zdet, zldet(2)
real   (8) :: d1(32,32), rdet, rldet(2)

call random_number(m1, sigma=1.0_8)
call random_number(d1, sigma=1.0_8)

zldet = lndet(m1)
zdet  = exp(zldet(1)) * zldet(2)   ! equals determinant(m1)

rldet = lndet(d1)
rdet  = exp(rldet(1)) * rldet(2)   ! equals determinant(d1)
\end{lstlisting}

% ===================================================================
% 31. pfaffian
% ===================================================================
\manentry{Generic Function}{pfaffian}{pfaffian}{Computes the Pfaffian $\Pf C$ of a complex antisymmetric matrix}{proc:pfaffian}

\manheading{Synopsis}
\lstset{language=[95]Fortran}
\begin{lstlisting}
Pf = Pfaffian(C, [alg])
\end{lstlisting}

\manheading{Brief}
Computes the Pfaffian $\Pf C$ (satisfying $(\Pf C)^2 = \det C$ for even dimension $N = 2n$, and $\Pf C = 0$ for odd $N$ in both algorithms) of an $N\times N$ complex skew-symmetric matrix \argname{C} using the \texttt{PFAPACK} library (default) or Simon Catterall's Gaussian elimination algorithm.

\manheading{Description}
For an $N\times N$ antisymmetric matrix $C = -C^T$, the Pfaffian vanishes identically when $N$ is odd, while for even dimension $N = 2n$ the determinant is the square of the Pfaffian, $(\Pf C)^2 = \det C$. The function \procname{pfaffian} (case-insensitive: \texttt{Pfaffian}) supports two algorithms selected via the optional argument \argname{alg}:
\begin{itemize}
  \item \textbf{\texttt{alg='P'} (Default --- \texttt{PFAPACK} library)}: Calls the \texttt{skpfa} subroutine from Michael Wimmer's \texttt{PFAPACK} library~\cite{Wimmer_2012,pfapack} (incorporated into \filename{libMatrix.a} as of April 22, 2026). Note that all \texttt{PFAPACK} subroutines can also be accessed directly by adding \lcodef{use f95_pfapack}. \texttt{PFAPACK} is more than $10\times$ faster than Catterall's algorithm for large matrices.
  \item \textbf{\texttt{alg='C'} or \texttt{'c'} (Simon Catterall's algorithm)}: Calls \texttt{array2\_pfaffian(C)}, an optimized Fortran implementation of Simon Catterall's skew-symmetric Gaussian elimination algorithm with pivoting~\cite{Catterall_2003}. Zero pivots terminate elimination early and return $0$ without division by zero.
\end{itemize}
Both \texttt{alg='P'} and \texttt{alg='C'} share the following edge-case behavior:
\begin{itemize}
  \item If \argname{C} is odd-dimensional (\texttt{MOD(N,2) /= 0}) or singular, \procname{pfaffian} returns \texttt{(0.0\_8, 0.0\_8)}.
  \item If \argname{C} is an empty $0\times 0$ matrix ($N = 0$), \procname{pfaffian} returns \texttt{(1.0\_8, 0.0\_8)}.
  \item If \argname{C} is non-square, \procname{pfaffian} returns \texttt{(NaN, NaN)}.
  \item If \texttt{skpfa} reports a non-zero \texttt{INFO}, execution terminates via \texttt{matrix\_error}.
  \item Omitting \argname{alg} or passing an empty string \texttt{alg=''} selects the default \texttt{'P'}.
\end{itemize}

\begin{warnbox}[Overflow Limit for Large Matrices]
For a $2n\times 2n$ complex antisymmetric matrix filled with Gaussian $\sigma = 1$ entries, \procname{pfaffian} can be used up to matrix dimensions $2n \sim 312$ before double-precision floating-point overflow occurs. For larger matrices ($2n \gtrsim 312$), use \procname{lnPfaffian}.
\end{warnbox}

\manheading{Dummy Arguments}
\begin{description}[leftmargin=2.2cm, style=nextline]
  \argitem{C}{complex(8), dimension(:,:), contiguous, intent(in)}{Input $N\times N$ complex antisymmetric matrix ($C^T = -C$).}
  \argitem{alg}{character(*), optional, intent(in)}{Algorithm selector: \texttt{'P'} (default, also used if empty \texttt{''}) uses \texttt{PFAPACK}'s \texttt{skpfa}; \texttt{'C'} or \texttt{'c'} uses Simon Catterall's algorithm.}
  \argitem{pfaffian}{complex(8), scalar, return value}{The computed Pfaffian $\Pf C$ (\texttt{(0.0\_8, 0.0\_8)} for odd $N$ or singular \argname{C}; \texttt{(1.0\_8, 0.0\_8)} for $N = 0$).}
\end{description}

\manheading{Example}
\lstset{language=[95]Fortran}
\begin{lstlisting}
use array_mod
complex(8) :: m1(16,16), pf1, pf2

call random_number(m1, sigma=1.0_8)
call antisymmetric_set(m1)

pf1 = Pfaffian(m1)            ! uses PFAPACK (default, >10x faster)
pf2 = Pfaffian(m1, alg='C')   ! uses Catterall's algorithm
! Note: pf1**2 == determinant(m1)
\end{lstlisting}

% ===================================================================
% 32. lnPfaffian
% ===================================================================
\manentry{Generic Function}{lnPfaffian}{lnPfaffian}{Computes the logarithm of the modulus and the complex phase of $\Pf C$ for large antisymmetric matrices}{proc:lnPfaffian}

\manheading{Synopsis}
\lstset{language=[95]Fortran}
\begin{lstlisting}
lnPf = lnPfaffian(C, [alg])
! Reconstruct Pfaffian: Pf = exp(lnPf(1)) * lnPf(2)
\end{lstlisting}

\manheading{Brief}
Returns a 2-element \typename{complex(8)} array $\texttt{lnPf(1:2)}$ containing $\ln|\Pf C|$ in \texttt{lnPf(1)} and the unit-modulus complex phase $\e{i\arg(\Pf C)}$ in \texttt{lnPf(2)} (or \texttt{lnPf(1) = (}$-\infty$\texttt{, 0.0\_8)} and \texttt{lnPf(2) = (0.0\_8, 0.0\_8)} when $\Pf C = 0$).

\manheading{Description}
To prevent floating-point overflow on large antisymmetric matrices ($2n \gtrsim 312$), \procname{lnPfaffian} decomposes the Pfaffian into its logarithmic magnitude and complex phase:
\begin{align}
  \Pf C = \exp(\texttt{lnPf(1)}) \times \texttt{lnPf(2)}\, ,
\end{align}
where:
\begin{align}
  \texttt{lnPf(1)} &= \ln|\Pf C|\, , \\
  \texttt{lnPf(2)} &= \e{i\arg(\Pf C)}\, .
\end{align}
\begin{itemize}
  \item \textbf{Default (\texttt{alg='P'})}: Calls \texttt{PFAPACK}'s overflow-safe subroutine \texttt{skpf10(C\_, lpf, info=info)} (aborting via \texttt{matrix\_error} if \texttt{info /= 0}), which returns the base-10 mantissa \texttt{lpf(1)} and exponent \texttt{lpf(2)} such that $\Pf C = \texttt{lpf(1)} \times 10^{\texttt{lpf(2)}}$. If \texttt{abs(lpf(1)) == 0.0\_dp} (singular matrix), \texttt{log} and \texttt{atan2} are skipped and $\texttt{[(}-\infty\texttt{, 0), (0, 0)]}$ is returned directly; otherwise it converts them to natural logarithm and phase via:
  \begin{align}
    \texttt{lnPf(1)} &= \ln|\texttt{lpf(1)}| + \Re(\texttt{lpf(2)})\ln 10\, , \\
    \texttt{lnPf(2)} &= \exp\!\left(i\,\texttt{atan2}(\Im(\texttt{lpf(1)}), \Re(\texttt{lpf(1)}))\right)\, .
  \end{align}
  \item \textbf{\texttt{alg='C'} or \texttt{'c'}}: Calls \texttt{array2\_log\_pfaffian(C)} (Simon Catterall's algorithm), accumulating $\sum_{i=1,3,\ldots}\ln|M_{i+1,i}|$ and the product of unit phases $M_{i+1,i}/|M_{i+1,i}|$ at each $2\times 2$ pivot step; zero pivots terminate elimination early and return $\texttt{[(}-\infty\texttt{, 0), (0, 0)]}$ without division by zero.
\end{itemize}
Both \texttt{alg='P'} and \texttt{alg='C'} handle edge cases identically:
\begin{itemize}
  \item If \argname{C} is odd-dimensional (\texttt{MOD(N,2) /= 0}) or singular ($\Pf C = 0$), \procname{lnPfaffian} returns \texttt{lnPf(1) = (}$-\infty$\texttt{, 0.0\_8)} (via \texttt{IEEE\_VALUE(0.0\_dp, IEEE\_NEGATIVE\_INF)}) and \texttt{lnPf(2) = (0.0\_8, 0.0\_8)}, satisfying $\exp(\texttt{lnPf(1)}) \times \texttt{lnPf(2)} = \texttt{(0.0\_8, 0.0\_8)}$.
  \item If \argname{C} is an empty $0\times 0$ matrix ($N = 0$), \procname{lnPfaffian} returns \texttt{lnPf(1) = (0.0\_8, 0.0\_8)} and \texttt{lnPf(2) = (1.0\_8, 0.0\_8)}.
  \item If \argname{C} is non-square, both elements \texttt{lnPf(1)} and \texttt{lnPf(2)} are set to \texttt{(NaN, NaN)}.
  \item If \texttt{skpf10} fails with \texttt{info /= 0}, execution terminates via \texttt{matrix\_error}.
  \item Empty \texttt{alg=''} defaults to \texttt{'P'}.
\end{itemize}

\manheading{Dummy Arguments}
\begin{description}[leftmargin=2.2cm, style=nextline]
  \argitem{C}{complex(8), dimension(:,:), contiguous, intent(in)}{Input $N\times N$ complex antisymmetric matrix ($C^T = -C$).}
  \argitem{alg}{character(*), optional, intent(in)}{Algorithm selector: \texttt{'P'} (default, also used if empty \texttt{''}; \texttt{PFAPACK} \texttt{skpf10}) or \texttt{'C'}/\texttt{'c'} (Simon Catterall's algorithm).}
  \argitem{pfaffian}{complex(8), dimension(2), return value}{Two-element complex array: \texttt{pfaffian(1)} $= \ln|\Pf C|$ and \texttt{pfaffian(2)} $= \e{i\arg(\Pf C)}$, or \texttt{pfaffian(1) = (}$-\infty$\texttt{, 0.0\_8)} and \texttt{pfaffian(2) = (0.0\_8, 0.0\_8)} when \argname{C} is odd-dimensional or singular ($\exp(\texttt{pfaffian(1)}) \times \texttt{pfaffian(2)} = 0$).}
\end{description}

\manheading{Example}
\lstset{language=[95]Fortran}
\begin{lstlisting}
use array_mod
complex(8) :: m1(64,64), pf, lnPf(2), lnPf_cat(2)

call random_number(m1, sigma=1.0_8)
call antisymmetric_set(m1)

lnPf     = lnPfaffian(m1)            ! PFAPACK skpf10 (default)
lnPf_cat = lnPfaffian(m1, alg='C')   ! Catterall algorithm
pf       = exp(lnPf(1)) * lnPf(2)    ! equals Pfaffian(m1)
\end{lstlisting}

% -------------------------------------------------------------------
\subsubsection{Eigenvalues and eigenvectors:}
\label{subsubsec:eigenvalues_eigenvectors}

% ===================================================================
% 33. eigenvalues
% ===================================================================
\manentry{Generic Function}{eigenvalues}{eigenvalues}{Computes the eigenvalues of a general, Hermitian, or real symmetric square matrix}{proc:eigenvalues}

\manheading{Synopsis}
\lstset{language=[95]Fortran}
\begin{lstlisting}
eval = eigenvalues(C, [mtype], [sortby])
\end{lstlisting}

\manheading{Brief}
Computes all $N$ eigenvalues of an $N\times N$ \typename{complex(8)} or \typename{real(8)} matrix \argname{C} using \texttt{LAPACK} (\texttt{zgeev}, \texttt{zheev}, \texttt{dgeev}, or \texttt{dsyev}), with optional sorting.

\manheading{Description}
Given a contiguous $N\times N$ \typename{complex(8)} or \typename{real(8)} matrix \argname{C}, \procname{eigenvalues} allocates and returns a 1D \typename{complex(8)} array \argname{eigenval} of length $N$ containing the eigenvalues $\{\lambda_1, \ldots, \lambda_N\}$ of \argname{C}:
\begin{itemize}
  \item \textbf{General Complex Matrix (\argname{C} is \typename{complex(8)}, default \texttt{mtype='ZG'})}: Calls \texttt{LAPACK}'s \texttt{zgeev} with \texttt{JOBVL='N'}, \texttt{JOBVR='N'}.
  \item \textbf{Hermitian Complex Matrix (\argname{C} is \typename{complex(8)}, \texttt{mtype='ZH'} or \texttt{'H'})}: Calls \texttt{LAPACK}'s \texttt{zheev} with \texttt{JOBZ='N'}, \texttt{UPLO='U'}. This is significantly faster than \texttt{zgeev} and returns real eigenvalues (stored in the \typename{complex(8)} return vector with zero imaginary parts) ordered by \texttt{zheev} in ascending algebraic value.
  \item \textbf{General Real Matrix (\argname{C} is \typename{real(8)}, default \texttt{mtype='DG'})}: Calls \texttt{LAPACK}'s \texttt{dgeev} with \texttt{JOBVL='N'}, \texttt{JOBVR='N'}. Because a general real matrix can have complex-conjugate pairs of eigenvalues, the returned array is \typename{complex(8)}.
  \item \textbf{Real Symmetric Matrix (\argname{C} is \typename{real(8)}, \texttt{mtype='DS'} or \texttt{'S'})}: Calls \texttt{LAPACK}'s \texttt{dsyev} with \texttt{JOBZ='N'}, \texttt{UPLO='U'}. The real eigenvalues (ordered ascending by \texttt{dsyev}) are returned as a \typename{complex(8)} array (which can be assigned directly to a \typename{real(8)} vector \texttt{deval = eigenvalues(B, mtype='DS')} in Fortran).
\end{itemize}
If \argname{C} is non-square, or if the \texttt{LAPACK} routine fails to converge (\texttt{INFO > 0}), the returned vector of length \texttt{size(C,1)} is filled with \procname{NaN} (test with \lcodef{isNaN(eval)}); only an illegal-argument failure (\texttt{INFO < 0}) or an invalid internal \texttt{job} value terminates execution via \texttt{matrix\_error}.

If the optional argument \argname{sortby} is present, the eigenvalues are sorted prior to return via \texttt{sort(eigenval, by=sortby)} using any of the two-character complex sorting codes (\texttt{'RF'}, \texttt{'RR'}, \texttt{'IF'}, \texttt{'IR'}, \texttt{'MF'}, \texttt{'MR'}).

\begin{warnbox}[Eigenvalue Ordering Subtlety in LAPACK \texttt{zgeev} / \texttt{dgeev}]
\begin{enumerate}
  \item When \texttt{mtype='ZH'} (or \texttt{'DS'}) is used, \texttt{zheev} (or \texttt{dsyev}) automatically sorts the real eigenvalues in ascending algebraic order ($\lambda_1 \le \lambda_2 \le \cdots \le \lambda_N$), which differs from the unsorted order returned by \texttt{zgeev} (or \texttt{dgeev}).
  \item Even for the same general matrix \argname{A}, for large enough dimensions ($N \gtrsim 100$), \texttt{LAPACK}'s \texttt{zgeev}/\texttt{dgeev} does \emph{not} return the eigenvalues in the exact same order when called with \texttt{job='N'} (\procname{eigenvalues(A)}) versus \texttt{job='V'} (\procname{eigenvectors(A)}). Always pass \texttt{sortby='RF'} (or sort both results with \procname{sort}) before comparing them element-by-element!
\end{enumerate}
\end{warnbox}

\manheading{Dummy Arguments}
\begin{description}[leftmargin=2.2cm, style=nextline]
  \argitem{C}{complex(8) or real(8), dimension(:,:), contiguous, intent(in)}{Input $N\times N$ square matrix.}
  \argitem{mtype}{character(*), optional, intent(in)}{Matrix symmetry type: \texttt{'ZH'} for complex Hermitian, \texttt{'DS'} for real symmetric, or \texttt{'ZG'}/\texttt{'DG'} (default) for general matrices.}
  \argitem{sortby}{character(*), optional, intent(in)}{Two-character sorting code forwarded to \procname{sort}: \texttt{'RF'} / \texttt{'RR'} (real part forward/reverse), \texttt{'IF'} / \texttt{'IR'} (imaginary part forward/reverse), \texttt{'MF'} / \texttt{'MR'} (modulus forward/reverse).}
  \argitem{eigenval}{complex(8), dimension(:), allocatable, return value}{Allocated 1D array of length $N$ containing the eigenvalues of \argname{C}.}
\end{description}

\manheading{Example}
\lstset{language=[95]Fortran}
\begin{lstlisting}
use array_mod
complex(8) :: A(8,8), eval(8)
real   (8) :: B(8,8), deval(8)

call random_number(A); call random_number(B)

eval  = eigenvalues(A)                          ! general complex matrix
eval  = eigenvalues(A, sortby='RF')             ! sorted ascending by Re(lambda)

call hermitian_set(A)
eval  = eigenvalues(A, mtype='ZH')              ! faster for Hermitian A

eval  = eigenvalues(B)                          ! general real B has complex eigenvalues
call symmetric_set(B)
deval = eigenvalues(B, mtype='DS', sortby='RR') ! real symmetric B, sorted descending
\end{lstlisting}

% ===================================================================
% 34. eigenvectors
% ===================================================================
\manentry{Generic Function}{eigenvectors}{eigenvectors}{Computes both the eigenvalues and normalized right eigenvectors of a square matrix}{proc:eigenvectors}

\manheading{Synopsis}
\lstset{language=[95]Fortran}
\begin{lstlisting}
W    = eigenvectors(C, [mtype], [sortby])
eval = W(:, 1 )   ! Column 1 contains the N eigenvalues
evec = W(:, 2:)   ! Columns 2..N+1 contain the N right eigenvectors
\end{lstlisting}

\manheading{Brief}
Computes the eigenvalues $\lambda_i$ and normalized right eigenvectors $v_i$ ($C v_i = \lambda_i v_i$) of an $N\times N$ \typename{complex(8)} or \typename{real(8)} matrix \argname{C}, packing them into an $N\times (N+1)$ complex array \argname{W}.

\manheading{Description}
Given a contiguous $N\times N$ \typename{complex(8)} or \typename{real(8)} matrix \argname{C}, \procname{eigenvectors} allocates and returns an $N\times (N+1)$ \typename{complex(8)} matrix \argname{evs} (\argname{W}) packed as follows:
\begin{itemize}
  \item \textbf{First column \texttt{W(:,1)}}: Contains the $N$ eigenvalues $\lambda_i = \texttt{eval}(i)$ for $i=1,\ldots,N$.
  \item \textbf{Remaining $N$ columns \texttt{W(:,2:N+1)}}: Contains the $N\times N$ eigenvector matrix $V = \texttt{evec}(:,1:N)$, where the $i$-th column $v_i = \texttt{evec}(:,i) = \texttt{W}(:,i+1)$ is the Euclidean-normalized right eigenvector corresponding to $\lambda_i$:
  \begin{align}
    C\, v_i = \lambda_i\, v_i \quad (i = 1,\ldots,N) \quad \Longleftrightarrow \quad C\, V = V\, \Lambda\, , \quad \Lambda = \texttt{diagonal(eval)}\, .
  \end{align}
\end{itemize}
Depending on the type of \argname{C} and the optional argument \argname{mtype}, \procname{eigenvectors} calls:
\begin{itemize}
  \item \texttt{zgeev('N','V',...)} for a general \typename{complex(8)} matrix (default \texttt{mtype='ZG'}).
  \item \texttt{zheev('V','U',...)} for a Hermitian \typename{complex(8)} matrix (\texttt{mtype='ZH'}), yielding an unitary eigenvector matrix $V$ ($V^\dagger V = I$).
  \item \texttt{dgeev('N','V',...)} for a general \typename{real(8)} matrix (default \texttt{mtype='DG'}). When \texttt{dgeev} encounters a complex-conjugate eigenvalue pair $\lambda_j = \lambda_{j+1}^*$ with $\Im(\lambda_j) > 0$, \procname{eigenvectors} automatically reconstructs the complex right eigenvectors from the real and imaginary parts stored in columns $j$ and $j+1$ of \texttt{VR}:
  \begin{align}
    v_j = \texttt{VR}(:,j) + i\,\texttt{VR}(:,j+1)\, , \qquad v_{j+1} = \texttt{VR}(:,j) - i\,\texttt{VR}(:,j+1)\, .
  \end{align}
  \item \texttt{dsyev('V','U',...)} for a real symmetric \typename{real(8)} matrix (\texttt{mtype='DS'}), yielding real eigenvalues and an orthogonal real eigenvector matrix $V$ ($V^T V = I$). Because Fortran converts \typename{complex(8)} to \typename{real(8)} on assignment by taking the real part, the returned array can be assigned directly to a \typename{real(8)} array \texttt{V(n,n+1)} when \argname{C} is real symmetric.
\end{itemize}

When the optional argument \argname{sortby} is supplied (\texttt{'RF'}, \texttt{'RR'}, \texttt{'IF'}, \texttt{'IR'}, \texttt{'MF'}, or \texttt{'MR'}), \procname{eigenvectors} sorts the eigenvalues using \texttt{sort(eigenval, P=pos, by=sortby)} and reorders the columns of the eigenvector matrix using the permutation array \texttt{pos} so that column \texttt{evec(:,i)} always corresponds to eigenvalue \texttt{eval(i)}.

As for \procname{eigenvalues}, non-square input or a \texttt{LAPACK} convergence failure (\texttt{INFO > 0}) does not abort: the whole $N\times(N+1)$ result (with $N = \texttt{size(C,1)}$) is filled with \procname{NaN}.

\manheading{Dummy Arguments}
\begin{description}[leftmargin=2.2cm, style=nextline]
  \argitem{C}{complex(8) or real(8), dimension(:,:), contiguous, intent(in)}{Input $N\times N$ square matrix.}
  \argitem{mtype}{character(*), optional, intent(in)}{Matrix symmetry type: \texttt{'ZH'} for complex Hermitian, \texttt{'DS'} for real symmetric, or \texttt{'ZG'}/\texttt{'DG'} (default) for general matrices.}
  \argitem{sortby}{character(*), optional, intent(in)}{Two-character sorting specification (\texttt{'RF'}, \texttt{'RR'}, \texttt{'IF'}, \texttt{'IR'}, \texttt{'MF'}, \texttt{'MR'}) applied simultaneously to eigenvalues and eigenvector columns.}
  \argitem{evs}{complex(8), dimension(:,:), allocatable, return value}{Allocated $N\times(N+1)$ matrix where \texttt{evs(:,1)} holds the eigenvalues and \texttt{evs(:,2:N+1)} holds the corresponding right eigenvectors.}
\end{description}

\manheading{Example}
\lstset{language=[95]Fortran}
\begin{lstlisting}
use array_mod
integer, parameter      :: n = 8
complex(8)              :: A(n,n), eval(n), evec(n,n), W(n,n+1), L(n,n)
real   (8)              :: B(n,n), deval(n), devec(n,n), V(n,n+1)
integer                 :: i

call random_number(A, sigma=1.0_8)
call random_number(B, sigma=1.0_8)

! 1. General complex matrix A, sorted ascending by real part:
W    = eigenvectors(A, sortby='RF')
eval = W(:, 1 )
evec = W(:, 2:)

! Check A . v_i = lambda_i v_i and A . V = V . Lambda:
L = diagonal(eval)
print *, 'Max residual:', norm(lmatmul(A, evec) - lmatmul(evec, L))

! 2. Real symmetric matrix B, sorted descending by real part:
call symmetric_set(B)
V     = eigenvectors(B, mtype='DS', sortby='RR')
deval = V(:, 1 )
devec = V(:, 2:)
print *, 'Real sym residual:', &
         norm([(norm(matmul(B, devec(:,i)) - deval(i)*devec(:,i)), i=1,n)])
\end{lstlisting}

%   2.4 Tensor Products - Pauli Matrices & 2.5 Random Numbers
% ===================================================================
% Subsection 2.4: Tensor Products - Pauli Matrices
% ===================================================================
\subsection{Tensor Products --- Pauli Matrices}
\label{subsec:tensor_products}

The matrix Kronecker (tensor) product $Z = X \otimes Y$ and associated vectorization utilities are implemented in the standalone module \modname{tensorprod\_mod} (\filename{tensorprod\_mod.f90}). Because \modname{array\_mod} re-exports all public procedures of \modname{tensorprod\_mod} (\procname{tensorprod}, \procname{vec2col}, \procname{vec2row}, and \procname{vectorize}), a program containing \lcodef{use array_mod} automatically has access to them without needing a separate \lcodef{use tensorprod_mod} statement (though \modname{tensorprod\_mod} can also be used independently without \modname{array\_mod}).

Given an $N_m\times N_n$ matrix $X$ and an $N_p\times N_q$ matrix $Y$, their Kronecker product $Z = X \otimes Y$ is an $N_i\times N_j$ matrix with $N_i = N_m N_p$ and $N_j = N_n N_q$, whose elements are defined by:
\begin{align}
  Z_{ij} = X_{mn}\, Y_{pq}\, , \qquad
  i = N_p (m - 1) + p\, , \quad
  j = N_q (n - 1) + q\, ,
\end{align}
where $m=1,\ldots,N_m$, $n=1,\ldots,N_n$, $p=1,\ldots,N_p$, and $q=1,\ldots,N_q$.

% ===================================================================
% 35. tensorprod
% ===================================================================
\manentry{Pure Generic Function}{tensorprod}{tensorprod}{Computes the Kronecker tensor product $X_1 \otimes X_2 \otimes \cdots$ of 2 to 6 matrices or two 1D vectors}{proc:tensorprod}

\manheading{Synopsis}
\lstset{language=[95]Fortran}
\begin{lstlisting}
! Kronecker product of 2 to 6 rank-2 matrices:
Z = tensorprod(X1, X2)
Z = tensorprod(X1, X2, X3)
Z = tensorprod(X1, X2, X3, X4)
Z = tensorprod(X1, X2, X3, X4, X5)
Z = tensorprod(X1, X2, X3, X4, X5, X6)

! Outer tensor product of two rank-1 vectors (u^T \otimes v):
Z = tensorprod(u, v)
\end{lstlisting}

\manheading{Brief}
Computes the Kronecker product $X_1 \otimes X_2 \otimes \cdots \otimes X_k$ ($2 \le k \le 6$) of real or complex rank-2 matrices, or the outer product $u^T \otimes v$ of two rank-1 vectors.

\manheading{Description}
The pure generic function \procname{tensorprod} is overloaded across 12 module procedures in \modname{tensorprod\_mod}:
\begin{enumerate}
  \item \textbf{Two Matrices (\texttt{Z = tensorprod(X, Y)})}: For \typename{complex(8)} or \typename{real(8)} rank-2 arrays $X(N_m,N_n)$ and $Y(N_p,N_q)$, allocates and returns the $(N_m N_p)\times (N_n N_q)$ matrix $Z = X \otimes Y$.
  \item \textbf{Three to Six Matrices (\texttt{Z = tensorprod(X1, X2, ..., Xk)}, $3 \le k \le 6$)}: Evaluates the iterated Kronecker product $X_1 \otimes (X_2 \otimes \cdots \otimes X_k)$ right-to-left by successive binary calls to \texttt{tensorprod\_complex\_complex} or \texttt{tensorprod\_real\_real}. All input matrices in a call must have the same data type (\typename{complex(8)} or \typename{real(8)}).
  \item \textbf{Two 1D Vectors (\texttt{Z = tensorprod(u, v)})}: Given a 1D vector \argname{u} of length $N_n$ and a 1D vector \argname{v} of length $N_p$, treats \argname{u} as a $1\times N_n$ row matrix ($u^T$) and \argname{v} as an $N_p\times 1$ column matrix ($v$), returning the $N_p\times N_n$ rank-2 matrix $Z = u^T \otimes v = v u^T$:
  \begin{align}
    Z_{pn} = u_n\, v_p\, , \qquad p = 1,\ldots,N_p\, , \quad n = 1,\ldots,N_n\, .
  \end{align}
  \item \textbf{Mixed Vector--Matrix Tensor Products}: Because a 1D Fortran array does not distinguish between a row vector and a column vector, mixed tensor products of a vector and a matrix are performed by first converting the 1D vector into an $N\times 1$ column matrix via \procname{vec2col} or a $1\times N$ row matrix via \procname{vec2row}:
  \begin{itemize}
    \item \lcodef{tensorprod(vec2col(w), m1)} computes $w \otimes m_1$ of shape $(N_w M_1, N_1)$.
    \item \lcodef{tensorprod(vec2row(w), m1)} computes $w^T \otimes m_1$ of shape $(M_1, N_w N_1)$.
    \item \lcodef{tensorprod(m1, vec2col(w))} computes $m_1 \otimes w$ of shape $(M_1 N_w, N_1)$.
    \item \lcodef{tensorprod(m1, vec2row(w))} computes $m_1 \otimes w^T$ of shape $(M_1, N_1 N_w)$.
  \end{itemize}
\end{enumerate}

\manheading{Dummy Arguments}
\begin{description}[leftmargin=2.2cm, style=nextline]
  \argitem{X, Y, X1..X6}{complex(8) or real(8), dimension(:,:), intent(in)}{From 2 up to 6 rank-2 matrices of the same type (\typename{complex(8)} or \typename{real(8)}).}
  \argitem{u, v}{complex(8) or real(8), dimension(:), intent(in)}{Two rank-1 vectors of the same type (\argname{u} of length $N_n$, \argname{v} of length $N_p$).}
  \argitem{Z}{complex(8) or real(8), dimension(:,:), allocatable, return value}{Allocated rank-2 Kronecker product matrix.}
\end{description}

\manheading{Example}
\lstset{language=[95]Fortran}
\begin{lstlisting}
use array_mod
complex(8), allocatable :: s0(:,:), s1(:,:), s2(:,:), s3(:,:)
complex(8), allocatable :: g1(:,:), g3(:,:), m(:,:), m1(:,:)
complex(8)              :: w(4)
complex(8), parameter   :: IMU = (0.0_8, 1.0_8)

! Construct 16x16 10D Majorana-Weyl Gamma matrices from 2x2 Pauli matrices:
s0 = PauliMatrix(0); s1 = PauliMatrix(1); s2 = PauliMatrix(2); s3 = PauliMatrix(3)
g1 = IMU * tensorprod(s2, s2, s2, s2)
g3 = IMU * tensorprod(s2, s2, s0, s3)
call print(g1, form='Gamma', name='g1')
call print(g3, form='Gamma', name='g3')

! Vector-matrix tensor products using vec2col / vec2row:
allocate(m1(2,2)); call random_number(m1); call random_number(w)
m = tensorprod(vec2col(w), m1)   ! w   \otimes m1 (8 x 2)
m = tensorprod(vec2row(w), m1)   ! w^T \otimes m1 (2 x 8)
\end{lstlisting}

% ===================================================================
% 36. vec2col
% ===================================================================
\manentry{Pure Generic Function}{vec2col}{vec2col}{Reshapes a 1D vector of length $N_p$ into an $N_p \times 1$ rank-2 column matrix}{proc:vec2col}

\manheading{Synopsis}
\lstset{language=[95]Fortran}
\begin{lstlisting}
Z = vec2col(u)
\end{lstlisting}

\manheading{Brief}
Converts a rank-1 \typename{complex(8)} or \typename{real(8)} array \argname{u} of size $N_p$ into an allocatable rank-2 column matrix \argname{Z} of shape $(N_p, 1)$.

\manheading{Description}
The pure generic function \procname{vec2col} allocates an $N_p\times 1$ matrix \argname{Z} (where $N_p = \texttt{size(u)}$) and sets $\texttt{Z(:,1) = u}$. This allows a 1D vector to participate unambiguously as a column vector in \procname{tensorprod} or matrix multiplication routines.

\manheading{Dummy Arguments}
\begin{description}[leftmargin=2.2cm, style=nextline]
  \argitem{u}{complex(8) or real(8), dimension(:), intent(in)}{Input 1D vector of length $N_p$.}
  \argitem{Z}{complex(8) or real(8), dimension(:,:), allocatable, return value}{Output rank-2 column matrix of shape $(N_p, 1)$ with $Z(p,1) = u(p)$.}
\end{description}

\manheading{Example}
\lstset{language=[95]Fortran}
\begin{lstlisting}
use array_mod
complex(8)              :: w(4)
complex(8), allocatable :: m(:,:)

call random_number(w)
m = vec2col(w)   ! shape(m) is [4, 1], m(:,1) == w(:)
\end{lstlisting}

% ===================================================================
% 37. vec2row
% ===================================================================
\manentry{Pure Generic Function}{vec2row}{vec2row}{Reshapes a 1D vector of length $N_n$ into a $1 \times N_n$ rank-2 row matrix}{proc:vec2row}

\manheading{Synopsis}
\lstset{language=[95]Fortran}
\begin{lstlisting}
Z = vec2row(u)
\end{lstlisting}

\manheading{Brief}
Converts a rank-1 \typename{complex(8)} or \typename{real(8)} array \argname{u} of size $N_n$ into an allocatable rank-2 row matrix \argname{Z} of shape $(1, N_n)$.

\manheading{Description}
The pure generic function \procname{vec2row} allocates a $1\times N_n$ matrix \argname{Z} (where $N_n = \texttt{size(u)}$) and sets $\texttt{Z(1,:) = u}$. Notice that \texttt{vec2row(u)} represents the transpose $u^T$ of the column vector $\texttt{vec2col(u)}$ (without complex conjugation).

\manheading{Dummy Arguments}
\begin{description}[leftmargin=2.2cm, style=nextline]
  \argitem{u}{complex(8) or real(8), dimension(:), intent(in)}{Input 1D vector of length $N_n$.}
  \argitem{Z}{complex(8) or real(8), dimension(:,:), allocatable, return value}{Output rank-2 row matrix of shape $(1, N_n)$ with $Z(1,n) = u(n)$.}
\end{description}

\manheading{Example}
\lstset{language=[95]Fortran}
\begin{lstlisting}
use array_mod
real(8)              :: v(4)
real(8), allocatable :: d(:,:)

call random_number(v)
d = vec2row(v)   ! shape(d) is [1, 4], d(1,:) == v(:)
\end{lstlisting}

% ===================================================================
% 38. vectorize
% ===================================================================
\manentry{Pure Generic Function}{vectorize}{vectorize}{Flattens a rank-2 matrix into a 1D vector by stacking its columns in column-major order}{proc:vectorize}

\manheading{Synopsis}
\lstset{language=[95]Fortran}
\begin{lstlisting}
u = vectorize(Z)
\end{lstlisting}

\manheading{Brief}
Constructs a 1D vector $\operatorname{vec}(Z)$ of length $N_m N_n$ by stacking the successive columns of an $N_m\times N_n$ matrix \argname{Z} one below the other.

\manheading{Description}
Given an $N_m\times N_n$ \typename{complex(8)} or \typename{real(8)} matrix \argname{Z}, \procname{vectorize} allocates and returns a 1D array \argname{u} of length $N_i = N_m N_n$ defined by the standard column-major vectorization mapping:
\begin{align}
  u(i) = Z(m,n)\, , \qquad i = (n - 1) N_m + m\, , \quad m = 1,\ldots,N_m\, , \quad n = 1,\ldots,N_n\, .
\end{align}
In linear algebra and tensor calculus, $\operatorname{vec}(Z)$ satisfies the fundamental Kronecker-product identity $\operatorname{vec}(A X B) = (B^T \otimes A)\operatorname{vec}(X)$.

\manheading{Dummy Arguments}
\begin{description}[leftmargin=2.2cm, style=nextline]
  \argitem{Z}{complex(8) or real(8), dimension(:,:), intent(in)}{Input $N_m\times N_n$ rank-2 matrix.}
  \argitem{u}{complex(8) or real(8), dimension(:), allocatable, return value}{Allocated 1D vector of length $N_m N_n$ containing the stacked columns of \argname{Z}.}
\end{description}

\manheading{Example}
\lstset{language=[95]Fortran}
\begin{lstlisting}
use array_mod
complex(8)              :: m(3,4)
complex(8), allocatable :: w(:)
real   (8)              :: d(3,4)
real   (8), allocatable :: v(:)

call random_number(m); call random_number(d)

w = vectorize(m)   ! 1D complex vector of size 12: [m(:,1), m(:,2), m(:,3), m(:,4)]
v = vectorize(d)   ! 1D real    vector of size 12: [d(:,1), d(:,2), d(:,3), d(:,4)]
\end{lstlisting}

% ===================================================================
% Subsection 2.5: Random Numbers
% ===================================================================
\subsection{Random Numbers}
\label{subsec:random_numbers}

\modname{array\_mod} extends the intrinsic Fortran \texttt{random\_number} subroutine to support \typename{complex(8)} scalars and rank-1, rank-2, and rank-3 arrays, as well as Gaussian (normal) random variates via the Box--Muller transform for both \typename{real(8)} and \typename{complex(8)} scalars and arrays. It also provides \procname{matrix\_random\_init} to seed the pseudorandom number generator from the operating system's entropy source \filename{/dev/urandom}.

% ===================================================================
% 39. random_number
% ===================================================================
\manentry{Generic Subroutine}{random\_number}{random\textunderscore number}{Overloaded uniform and Box--Muller Gaussian random number generator for real and complex scalars and arrays}{proc:random_number}

\manheading{Synopsis}
\lstset{language=[95]Fortran}
\begin{lstlisting}
! 1. Uniform random numbers in [0, 1) (real) or [0, 1) + i [0, 1) (complex):
call random_number(x)

! 2. Gaussian (normal) random numbers with standard deviation sigma:
call random_number(x, sigma=s)
\end{lstlisting}

\manheading{Brief}
Populates a \typename{real(8)} or \typename{complex(8)} scalar, 1D vector, 2D matrix, or 3D array \argname{x} with uniform random numbers (when \argname{sigma} is omitted) or zero-mean Gaussian random numbers of standard deviation $\sigma$ (when \argname{sigma} is provided).

\manheading{Description}
The generic interface \procname{random\_number} overloads the Fortran intrinsic \texttt{random\_number} with 12 module procedures:
\begin{enumerate}
  \item \textbf{Uniform Distribution (\argname{sigma} omitted)}:
  \begin{itemize}
    \item For any \typename{real} scalar or array of any rank, resolves to the standard Fortran intrinsic \texttt{random\_number(x)}, drawing independent uniform variates in $[0,1)$.
    \item For a \typename{complex(8)} scalar $z$, rank-1 array $w_1(:)$, rank-2 array $m_1(:,:)$, or rank-3 array $t_1(:,:,:)$, draws independent uniform variates $X, Y \in [0,1)$ for the real and imaginary parts and sets $C = X + i Y$.
  \end{itemize}
  \item \textbf{Gaussian Distribution (\texttt{sigma=s} present)}:
  Generates independent zero-mean Gaussian random numbers using the Box--Muller transform from two uniform variates $x, y \in [0,1)$ (using $1 - x \in (0,1]$ inside the logarithm to avoid the $\ln 0$ singularity at $x = 0$):
  \begin{itemize}
    \item For a \typename{real(8)} scalar $r$, 1D array $v_1(:)$, 2D array $d_1(:,:)$, or 3D array $s_1(:,:,:)$:
    \begin{align}
      C = \sigma \sqrt{-2\ln(1 - x)}\,\cos(2\pi y) \sim \mathcal{N}(0, \sigma^2)\, , \qquad \vev{C} = 0\, , \quad \vev{C^2} = \sigma^2\, .
    \end{align}
    \item For a \typename{complex(8)} scalar $z$, 1D array $w_1(:)$, 2D array $m_1(:,:)$, or 3D array $t_1(:,:,:)$:
    \begin{align}
      C = \sigma \sqrt{-2\ln(1 - x)}\,\left[\cos(2\pi y) + i\sin(2\pi y)\right]\, ,
    \end{align}
    so that the real and imaginary parts $\Re C$ and $\Im C$ are independent normal variates each with standard deviation $\sigma$ ($\vev{(\Re C)^2} = \vev{(\Im C)^2} = \sigma^2$, and $\vev{|C|^2} = 2\sigma^2$). (To obtain $\vev{|C|^2} = 1$, pass $\texttt{sigma = 1.0\_8 / sqrt(2.0\_8)}$.)
  \end{itemize}
\end{enumerate}

\manheading{Dummy Arguments}
\begin{description}[leftmargin=2.2cm, style=nextline]
  \argitem{x}{complex(8) or real(8), scalar or dimension(:), (:,:), or (:,:,:), intent(out)}{Target scalar or array of rank 1, 2, or 3 to be populated with random numbers.}
  \argitem{sigma}{real(8), scalar, optional, intent(in)}{When present, selects the Box--Muller Gaussian distribution with standard deviation $\sigma = \argname{sigma}$ for the real (and, for complex targets, imaginary) part.}
\end{description}

\manheading{Example}
\lstset{language=[95]Fortran}
\begin{lstlisting}
use array_mod
complex(8) :: t1(4,4,3), m1(4,4), w1(4), z
real   (8) :: s1(4,4,3), d1(4,4), v1(4), r, s

call matrix_random_init
s = 1.0_8

! Uniform complex random numbers in [0,1) + i[0,1):
call random_number(t1) ; call random_number(m1)
call random_number(w1) ; call random_number(z )

! Gaussian complex random numbers with Re, Im ~ N(0, s^2):
call random_number(t1, sigma=s) ; call random_number(m1, sigma=s)
call random_number(w1, sigma=s) ; call random_number(z , sigma=s)

! Gaussian real random numbers ~ N(0, s^2):
call random_number(s1, sigma=s) ; call random_number(d1, sigma=s)
call random_number(v1, sigma=s) ; call random_number(r , sigma=s)
\end{lstlisting}

% ===================================================================
% 40. matrix_random_init
% ===================================================================
\manentry{Subroutine}{matrix\_random\_init}{matrix\textunderscore random\textunderscore init}{Seeds the Fortran intrinsic random number generator from \texttt{/dev/urandom}}{proc:matrix_random_init}

\manheading{Synopsis}
\lstset{language=[95]Fortran}
\begin{lstlisting}
call matrix_random_init
\end{lstlisting}

\manheading{Brief}
Initializes the intrinsic \texttt{random\_seed} state vector using cryptographic-grade OS entropy read from \filename{/dev/urandom}.

\manheading{Description}
The subroutine \procname{matrix\_random\_init} (defined in \modname{matrix\_mod\_common}) queries the required seed array size \texttt{nseeds} via \texttt{call random\_seed(size=nseeds)}, allocates an integer array \texttt{seeds(nseeds)}, opens \filename{/dev/urandom} as an unformatted stream file, reads \texttt{nseeds} random integers, clears their sign bit with \texttt{seeds = iand(seeds, huge(seeds))} (so that all seeds are non-negative without the overflow of \texttt{ABS(-huge-1)}), makes sure the seed vector is not all zeros, and seeds the Fortran runtime PRNG via \texttt{call random\_seed(put=seeds)}. If \filename{/dev/urandom} cannot be opened or read (e.g., on Windows or in restricted containers), it prints a warning to \procname{f\_mout} and falls back to \texttt{call random\_seed()}, whose seeding is processor dependent (non-repeatable with \texttt{gfortran} $\ge 7$). It also logs a status comment line to the module output unit \procname{f\_mout}:
\begin{verbatim}
# matrix_random_init: Initializing random_number from /dev/urandom
\end{verbatim}

\manheading{Dummy Arguments}
None. (Reads from \filename{/dev/urandom} and writes an initialization banner to the public unit variable \procname{f\_mout}.)

\manheading{Example}
\lstset{language=[95]Fortran}
\begin{lstlisting}
use array_mod
complex(8) :: m(4,4)

call matrix_random_init           ! seed PRNG from /dev/urandom
call random_number(m, sigma=1._8) ! draw fresh non-repeatable Gaussian matrix
\end{lstlisting}

%   2.6 Utilities (Saving and printing, Printing arrays, NaN stuff, Timing your program)
% ===================================================================
% Subsection 2.6: Utilities
% ===================================================================
\subsection{Utilities}
\label{subsec:utilities}

The utility procedures of \modname{array\_mod} handle formatted file I/O (\procname{save}, \procname{read}) controlled by the default I/O unit variables \procname{f\_mout} and \procname{f\_minput}, rich terminal/file/Mathematica/Gamma-matrix array printing (\procname{print}, \procname{printna}), IEEE quiet \procname{NaN} construction and checking (\procname{isNaN}), and program execution timing exported from \modname{timer\_mod} (\procname{timer\_mod\_init}, \procname{timer\_mod\_now}, \procname{timer\_mod\_print}).

% -------------------------------------------------------------------
\subsubsection{Saving and printing}
\label{subsubsec:saving_printing}

% ===================================================================
% 41. save
% ===================================================================
\manentry{Generic Subroutine}{save}{save}{Writes a rank-1, rank-2, or rank-3 real or complex array to an output unit in plain text format}{proc:save}
\addcontentsline{toc}{subsubsection}{\texorpdfstring{\texttt{\textbf{f\_mout}} \textnormal{\footnotesize--- Default output unit variable (\texttt{stdout})}}{f\textunderscore mout}}

\manheading{Synopsis}
\lstset{language=[95]Fortran}
\begin{lstlisting}
call save(C, [unit], [fmt])
\end{lstlisting}

\manheading{Brief}
Writes all elements of a 1D, 2D, or 3D \typename{complex(8)} or \typename{real(8)} array \argname{C} sequentially (one element per line) to a Fortran I/O unit so that the array can be restored later with \procname{read}.

\manheading{Description}
The generic subroutine \procname{save} outputs the elements of a rank-1, rank-2, or rank-3 array \argname{C} in a deterministic loop order:
\begin{itemize}
  \item \textbf{Rank-1 array \texttt{C(1:m)}}: Loops $i = 1, \ldots, m$.
  \item \textbf{Rank-2 array \texttt{C(1:m, 1:n)}}: Loops $i = 1, \ldots, m$ (outer) and $j = 1, \ldots, n$ (inner), writing one matrix element $C(i,j)$ per line in row-major line order.
  \item \textbf{Rank-3 array \texttt{C(1:m, 1:n, 1:k)}}: Loops $l = 1, \ldots, k$ (outermost), $i = 1, \ldots, m$, and $j = 1, \ldots, n$ (innermost), writing $C(i,j,l)$ (representing $A^l_{ij}$) one element per line.
\end{itemize}
For \typename{complex(8)} arrays, each line is written with format \texttt{'(2('//fmt//'))'} (real and imaginary parts on the same line); for \typename{real(8)} arrays, each line is written with format \texttt{'('//fmt//')'}. When a user format \argname{fmt} is given, it is prefixed with \texttt{1X}, i.e., \texttt{'(2(1X,'//fmt//'))'} and \texttt{'(1X,'//fmt//')'}, so that consecutive fields are always separated by a blank and can be read back by \procname{read} even when a value fills its whole field (e.g., \texttt{-1234.567} written with \texttt{fmt='F9.3'}). The default format \texttt{G28.17} always leaves leading blanks and is used unchanged. (Files written by external programs with a fixed-column format should be read with whitespace-separated parsing.)

When the optional argument \argname{unit} is omitted, \procname{save} writes to the public module integer variable \procname{f\_mout} (defined in \modname{matrix\_mod\_common} and initialized to \texttt{stdout = output\_unit}). Users can redirect all default module output globally by setting \lcodef{f_mout = my_unit}. When \argname{fmt} is omitted, the full double-precision format \texttt{'G28.17'} is used.

\manheading{Dummy Arguments}
\begin{description}[leftmargin=2.2cm, style=nextline]
  \argitem{C}{complex(8) or real(8), dimension(:), (:,:), or (:,:,:), intent(in)}{Input rank-1, rank-2, or rank-3 array to be saved.}
  \argitem{unit}{integer, scalar, optional, intent(in)}{Fortran I/O unit number for output. Defaults to \procname{f\_mout} (\texttt{stdout}).}
  \argitem{fmt}{character(*), optional, intent(in)}{Fortran numeric format descriptor per real component. Defaults to \texttt{'G28.17'}.}
\end{description}

\manheading{Example}
\lstset{language=[95]Fortran}
\begin{lstlisting}
use array_mod
complex(8) :: m1(4,4)
real   (8) :: d1(4,4)
integer    :: un

call random_number(m1); call random_number(d1)
open(newunit=un, file='matrices.dat', action='write')
call save(m1, unit=un, fmt='G28.17')
call save(d1, unit=un, fmt='G28.17')
close(un)
\end{lstlisting}

% ===================================================================
% 42. read
% ===================================================================
\manentry{Generic Subroutine}{read}{read}{Reads a rank-1, rank-2, or rank-3 real or complex array from an input unit}{proc:read}
\addcontentsline{toc}{subsubsection}{\texorpdfstring{\texttt{\textbf{f\_minput}} \textnormal{\footnotesize--- Default input unit variable (\texttt{stdin})}}{f\textunderscore minput}}

\manheading{Synopsis}
\lstset{language=[95]Fortran}
\begin{lstlisting}
call read(C, [unit])
\end{lstlisting}

\manheading{Brief}
Reads the elements of a pre-allocated 1D, 2D, or 3D \typename{complex(8)} or \typename{real(8)} array \argname{C} from a Fortran I/O unit in the exact order written by \procname{save}.

\manheading{Description}
The generic subroutine \procname{read} populates the array \argname{C} using list-directed input (\texttt{read(un, *)}) in the same loop order as \procname{save}:
\begin{itemize}
  \item For \typename{complex(8)} arrays, reads two real numbers \texttt{x, y} per element and assigns \texttt{C(...) = CMPLX(x, y, kind=dp)}.
  \item For \typename{real(8)} arrays, reads one real number \texttt{C(...)} per element.
\end{itemize}
Because the input is list-directed, any whitespace-separated layout is accepted; files written by \procname{save} with any \argname{fmt} (with or without the \texttt{1X} prefix of newer versions) are read identically. If a read fails (end of file, or a token that is not a number), execution terminates via \texttt{matrix\_error} with a message giving the procedure, the element indices, the unit, and the run-time error text.

When \argname{unit} is omitted, \procname{read} reads from the public module integer variable \procname{f\_minput} (defined in \modname{matrix\_mod\_common} and initialized to \texttt{stdin = input\_unit}).

\manheading{Dummy Arguments}
\begin{description}[leftmargin=2.2cm, style=nextline]
  \argitem{C}{complex(8) or real(8), dimension(:), (:,:), or (:,:,:), intent(inout)}{Target array of rank 1, 2, or 3 (must be allocated with the desired dimensions prior to the call).}
  \argitem{unit}{integer, scalar, optional, intent(in)}{Fortran I/O unit number to read from. Defaults to \procname{f\_minput} (\texttt{stdin}).}
\end{description}

\manheading{Example}
\lstset{language=[95]Fortran}
\begin{lstlisting}
use array_mod
complex(8) :: m1(4,4)
real   (8) :: d1(4,4)
integer    :: un

open(newunit=un, file='matrices.dat', action='read')
call read(m1, unit=un)
call read(d1, unit=un)
close(un)
\end{lstlisting}

% -------------------------------------------------------------------
\subsubsection{Printing arrays}
\label{subsubsec:printing_arrays}

% ===================================================================
% 43. print
% ===================================================================
\manentry{Generic Subroutine}{print}{print}{Prints an allocatable 1D or 2D real or complex array in Default, Mathematica, Info, or Gamma format}{proc:print}

\manheading{Synopsis}
\lstset{language=[95]Fortran}
\begin{lstlisting}
! Rank-2 allocatable matrix C(:,:):
call print(C, [unit], [fmt], [form], [name], [ips], [ipe], [jps], [jpe])

! Rank-1 allocatable vector C(:):
call print(C, [unit], [fmt], [form], [name], [ips], [ipe])
\end{lstlisting}

\manheading{Brief}
Formatted printing subroutine for \emph{allocatable} rank-1 and rank-2 \typename{complex(8)} or \typename{real(8)} arrays, preserving and displaying custom lower and upper array bounds, with support for subarray slicing and direct export to Wolfram Mathematica or Dirac $\Gamma$-matrix notation.

\manheading{Description}
In Fortran 2003+, declaring a dummy array argument with the \typename{allocatable} attribute passes its actual lower and upper bounds (\texttt{lbound(C)}, \texttt{ubound(C)}) into the subroutine. The generic subroutine \procname{print} exploits this feature to report the true index ranges of \argname{C} and to allow slicing by actual indices via \argname{ips}, \argname{ipe}, \argname{jps}, and \argname{jpe}. (To print a non-allocatable or static array, use \procname{printna}.)

The output style is controlled by the first character of the optional argument \argname{form} (case-insensitive):
\begin{itemize}
  \item \textbf{\texttt{form='Default'} or \texttt{'d'} (default)}:
  Prints a header banner displaying the array name (default \texttt{'AMAT'}), its dimensions $m\times n$ (or vector length $m$), and the active printing bounds \texttt{is : ie} and \texttt{js : je}, followed by the formatted array rows between \texttt{\#pmstart:} / \texttt{\#pmend:} (for matrices) or \texttt{\#pvstart:} / \texttt{\#pvend:} (for vectors) delimiters.
  \item \textbf{\texttt{form='Mathematica'} or \texttt{'m'}}:
  Outputs valid Wolfram Mathematica assignment syntax \texttt{<name>=\{...\};} where each floating-point number is wrapped inside \texttt{(Interpreter["Number"]["..."])} so that Fortran exponential notation (e.g., \texttt{0.123E-04}) is parsed seamlessly when copy-pasted or imported into Mathematica.
  \item \textbf{\texttt{form='Info'} or \texttt{'i'}}:
  Prints only a single summary line prefixed by \texttt{\#pminfo:} reporting the array name, dimensions, and index bounds without printing the array elements.
  \item \textbf{\texttt{form='Gamma'} or \texttt{'g'} (for rank-2 \typename{complex(8)} matrices only)}:
  Designed for visually inspecting Pauli and Dirac $\Gamma$ matrices whose entries lie in $\{0, \pm 1, \pm i\}$ (within a tolerance of $10^{-10}$). Prints a compact grid displaying \texttt{0}, \texttt{1}, \texttt{-1}, \texttt{i}, \texttt{-i} (or \texttt{NN} for any other value).
\end{itemize}

\manheading{Dummy Arguments}
\begin{description}[leftmargin=2.2cm, style=nextline]
  \argitem{C}{complex(8) or real(8), allocatable, dimension(:,:) or (:), intent(in)}{The allocatable 2D matrix or 1D vector to print.}
  \argitem{unit}{integer, scalar, optional, intent(in)}{Output Fortran unit number. Defaults to \procname{f\_mout} (\texttt{stdout}).}
  \argitem{fmt}{character(*), optional, intent(in)}{Numeric format descriptor for each real value. Defaults to \texttt{'G28.17'}.}
  \argitem{form}{character(*), optional, intent(in)}{Output format mode: \texttt{'Default'} (\texttt{'d'}), \texttt{'Mathematica'} (\texttt{'m'}), \texttt{'Info'} (\texttt{'i'}), or \texttt{'Gamma'} (\texttt{'g'}, rank-2 complex only).}
  \argitem{name}{character(*), optional, intent(in)}{Display variable name (up to 20 characters). Defaults to \texttt{'AMAT'}.}
  \argitem{ips, ipe}{integer, scalar, optional, intent(in)}{Start and end row (or vector) indices to print, clamped to $[\texttt{lbound(C,1)}, \texttt{ubound(C,1)}]$.}
  \argitem{jps, jpe}{integer, scalar, optional, intent(in)}{Start and end column indices to print (for rank-2 matrices), clamped to $[\texttt{lbound(C,2)}, \texttt{ubound(C,2)}]$.}
\end{description}

\manheading{Example}
\lstset{language=[95]Fortran}
\begin{lstlisting}
use array_mod
complex(8), allocatable :: m1(:,:)
real   (8), allocatable :: d1(:,:)

allocate(m1(4,4), d1(0:5, 0:5))
call random_number(m1); call random_number(d1)

call print(m1)
call print(m1, fmt='F8.3', form='Mathematica', name='MAT')

m1 = PauliMatrix(2)   ! sigma_2 Pauli matrix
call print(m1, form='Gamma', name='s2')

! Print only a 3x3 subblock of d1(0:5, 0:5) and its bounds info:
call print(d1, fmt='F8.3', name='d1', ips=1, ipe=3, jps=2, jpe=4)
call print(d1, form='Info', name='d1')
\end{lstlisting}

% ===================================================================
% 44. printna
% ===================================================================
\manentry{Generic Subroutine}{printna}{printna}{Prints a non-allocatable (static or assumed-shape) 1D or 2D real or complex array}{proc:printna}

\manheading{Synopsis}
\lstset{language=[95]Fortran}
\begin{lstlisting}
call printna(C, [unit], [fmt], [form], [name], [ips], [ipe], [jps], [jpe])
\end{lstlisting}

\manheading{Brief}
Wrapper around \procname{print} for non-allocatable (static, explicit-shape, or array-slice) rank-1 and rank-2 \typename{complex(8)} or \typename{real(8)} arrays.

\manheading{Description}
Because \procname{print} requires its dummy argument \argname{C} to have the \typename{allocatable} attribute, static arrays such as \lcodef{real(8) :: d2(4,4)} or array expressions cannot be passed directly to \procname{print}. The subroutine \procname{printna} accepts any non-allocatable rank-1 or rank-2 array \argname{C}, creates a temporary allocatable copy via \lcodef{ALLOCATE(Ca, source=C)}, invokes \procname{print} on \texttt{Ca} with all optional arguments forwarded, and deallocates \texttt{Ca}.

\begin{notebox}[1-Based Indexing in \texttt{printna}]
In Fortran, passing an array with custom lower bounds (such as \lcodef{real(8) :: sd(3:5, 6:8)}) to an assumed-shape dummy argument \lcodef{C(:,:)} resets the lower bounds inside \procname{printna} to $1$ (so \texttt{Ca} has bounds \texttt{1:3, 1:3}). Consequently, the optional section arguments \argname{ips}, \argname{ipe}, \argname{jps}, and \argname{jpe} in \procname{printna} always refer to 1-based offsets $1,\ldots,m$ and $1,\ldots,n$.
\end{notebox}

\manheading{Dummy Arguments}
\begin{description}[leftmargin=2.2cm, style=nextline]
  \argitem{C}{complex(8) or real(8), dimension(:,:) or (:), intent(in)}{Non-allocatable 2D matrix or 1D vector to print.}
  \argitem{unit, fmt, form, name, ips, ipe, jps, jpe}{optional, intent(in)}{Identical to the optional arguments of \procname{print} (note that \argname{jps}, \argname{jpe} are also accepted and ignored in the rank-1 overload).}
\end{description}

\manheading{Example}
\lstset{language=[95]Fortran}
\begin{lstlisting}
use array_mod
real   (8) :: sd(3:5, 6:8)   ! static 3x3 array with custom lower bounds
complex(8) :: sm(4, 4)

call random_number(sd); call random_number(sm)

call printna(sm, fmt='F8.3', name='sm')
! In printna, bounds are seen as 1:3, 1:3, so ips=2, jps=2 prints sd(4:5, 7:8):
call printna(sd, fmt='F8.3', name='d1', ips=2, jps=2)
\end{lstlisting}

% -------------------------------------------------------------------
\subsubsection{NaN stuff}
\label{subsubsec:nan_stuff}

% ===================================================================
% 45. NaN
% ===================================================================
\manentry{Pure Generic Function}{NaN}{NaN}{Returns an IEEE quiet NaN (Not-a-Number) value for real or complex variables and arrays}{proc:NaN}

\manheading{Synopsis}
\lstset{language=[95]Fortran}
\begin{lstlisting}
z = NaN()    ! Returns complex(8) (NaN, NaN)
z = NaN(z0)  ! Returns complex(8) (NaN, NaN)
r = NaN(r0)  ! Returns real   (8)  NaN
\end{lstlisting}

\manheading{Brief}
Generates an IEEE~754 quiet \texttt{NaN} (\texttt{IEEE\_QUIET\_NAN}) of type \typename{complex(8)} or \typename{real(8)} using the intrinsic \texttt{ieee\_arithmetic} module.

\manheading{Description}
The pure generic function \procname{NaN} (case-insensitive: \texttt{nan}), defined in \modname{matrix\_mod\_common}, constructs quiet Not-a-Number values for error signaling and uninitialized array poisoning:
\begin{itemize}
  \item \textbf{Zero-argument call \texttt{NaN()}}: Calls \texttt{IEEE\_VALUE(0.0\_dp, IEEE\_QUIET\_NAN)} and returns a \typename{complex(8)} value $r = (\texttt{NaN}, \texttt{NaN})$ whose real and imaginary parts are both quiet \texttt{NaN}. Note that because Fortran converts \typename{complex(8)} to \typename{real(8)} on assignment by extracting the real part, \lcodef{r = NaN()} or \lcodef{d1 = NaN()} also works for \typename{real(8)} scalars and arrays!
  \item \textbf{Complex argument \texttt{NaN(z)}}: Given a \typename{complex(8)} scalar \argname{z}, returns \typename{complex(8)} $(\texttt{NaN}, \texttt{NaN})$.
  \item \textbf{Real argument \texttt{NaN(r)}}: Given a \typename{real(8)} scalar \argname{r}, returns a \typename{real(8)} quiet \texttt{NaN}. When assigned to a complex variable \lcodef{z = NaN(r)}, the real part $\Re(z)$ is set to \texttt{NaN} and $\Im(z)$ to $0.0\_8$.
\end{itemize}

\manheading{Dummy Arguments}
\begin{description}[leftmargin=2.2cm, style=nextline]
  \argitem{z}{complex(8) or real(8), scalar, optional, intent(in)}{Optional mold scalar selecting a \typename{complex(8)} or \typename{real(8)} return value.}
  \argitem{r}{complex(8) or real(8), scalar, return value}{IEEE quiet \texttt{NaN} scalar (broadcastable by assignment to an entire array or array element).}
\end{description}

\manheading{Example}
\lstset{language=[95]Fortran}
\begin{lstlisting}
use array_mod
complex(8) :: m1(4,4), z
real   (8) :: d1(4,4), r

z  = nan()              ! sets Re(z) = NaN and Im(z) = NaN
z  = nan(r)             ! sets Re(z) = NaN and Im(z) = 0.0_8
z  = cmplx(real(z), nan(r), kind=8) ! sets Im(z) = NaN
r  = nan(r)             ! real(8) NaN (r = nan() also works)

m1 = nan()              ! sets all m1(i,j) = (NaN, NaN)
d1 = nan(r)             ! sets all d1(i,j) =  NaN
\end{lstlisting}

% ===================================================================
% 46. isNaN
% ===================================================================
\manentry{Pure Generic Function}{isNaN}{isNaN}{Tests whether a real or complex scalar or rank-1, rank-2, or rank-3 array contains any NaN value}{proc:isNaN}

\manheading{Synopsis}
\lstset{language=[95]Fortran}
\begin{lstlisting}
flag = isNaN(x)
\end{lstlisting}

\manheading{Brief}
Returns \texttt{.true.} if a \typename{real(8)} or \typename{complex(8)} scalar, or any element of a rank-1, rank-2, or rank-3 array \argname{x}, is an IEEE \texttt{NaN}.

\manheading{Description}
The pure generic function \procname{isNaN} combines the scalar interfaces in \modname{matrix\_mod\_common} and the array interfaces in \modname{matrix\_mod\_array} (8 module procedures in total):
\begin{itemize}
  \item \textbf{Real scalar \texttt{isNaN(r)}}: Returns \texttt{IEEE\_IS\_NAN(r)}.
  \item \textbf{Complex scalar \texttt{isNaN(z)}}: Returns \texttt{IEEE\_IS\_NAN(REAL(z)) .or. IEEE\_IS\_NAN(AIMAG(z))}.
  \item \textbf{Rank-1, Rank-2, or Rank-3 Array \texttt{isNaN(C)}}: Loops over all elements of the \typename{real(8)} or \typename{complex(8)} array \argname{C} and returns \texttt{.true.} if \emph{any} element is \texttt{NaN}, or \texttt{.false.} otherwise. Because non-conformable operations in \modname{array\_mod} (such as \procname{lmatmul}, \procname{mmmult}, \procname{trace2}, and \procname{paulimatrix}) return \procname{NaN}-filled arrays on dimension errors, and \procname{inverse}, \procname{determinant}, \procname{lndet}, \procname{pfaffian}, \procname{lnPfaffian}, \procname{eigenvalues}, and \procname{eigenvectors} return \procname{NaN} for non-square input (and \procname{inverse} also for singular matrices, the eigen-solvers also on \texttt{LAPACK} convergence failure), \procname{isNaN} provides a convenient way to detect invalid computations.
\end{itemize}

\manheading{Dummy Arguments}
\begin{description}[leftmargin=2.2cm, style=nextline]
  \argitem{x}{complex(8) or real(8), scalar or dimension(:), (:,:), or (:,:,:), intent(in)}{Scalar or rank-1, rank-2, or rank-3 array to test for \texttt{NaN} entries.}
  \argitem{itis}{logical, scalar, return value}{\texttt{.true.} if at least one component of \argname{x} is \texttt{NaN}; \texttt{.false.} otherwise.}
\end{description}

\manheading{Example}
\lstset{language=[95]Fortran}
\begin{lstlisting}
use array_mod
complex(8) :: m1(4,4), z
real   (8) :: d1(4,4), r

call random_number(m1); call random_number(d1)
print *, isNaN(m1), isNaN(d1)   ! .false.  .false.

m1(2,3) = nan()
d1(1,4) = nan(r)
print *, isNaN(m1), isNaN(d1)   ! .true.   .true.
\end{lstlisting}

% -------------------------------------------------------------------
\subsubsection{Timing your program}
\label{subsubsec:timing_program}

\modname{array\_mod} re-exports all public procedures and state variables of \modname{timer\_mod} (\filename{timer\_mod.f90}), enabling accurate tracking of both cumulative and incremental CPU time (via \texttt{cpu\_time}) and wall-clock time (via \texttt{system\_clock} with 64-bit integer counters), together with formatted timestamps (via \texttt{date\_and\_time}).

% ===================================================================
% 47. timer_mod_init
% ===================================================================
\manentry{Subroutine}{timer\_mod\_init}{timer\textunderscore mod\textunderscore init}{Initializes the timer\_mod CPU and wall-clock timers and prints a header banner}{proc:timer_mod_init}

\manheading{Synopsis}
\lstset{language=[95]Fortran}
\begin{lstlisting}
call timer_mod_init([un])
\end{lstlisting}

\manheading{Brief}
Records the initial CPU and wall-clock reference times, resets the call counter \procname{timer\_count} to zero, updates \procname{timer\_date}, and writes a two-line \texttt{\#TIMERINIT::} header to unit \argname{un}.

\manheading{Description}
Calling \procname{timer\_mod\_init} establishes the $t = 0$ origin for subsequent calls to \procname{timer\_mod\_now} and \procname{timer\_mod\_print}:
\begin{enumerate}
  \item Calls \texttt{cpu\_time(cpu\_start)} and initializes \texttt{cpu\_mark = cpu\_start}, \texttt{cpu\_now = cpu\_start}, and \texttt{timer\_count = 0}.
  \item Calls \texttt{system\_clock(wtime\_start, wtime\_rate)} using 64-bit integers (\typename{int64}) and sets \texttt{wtime\_mark = wtime\_start}, \texttt{wtime\_now = wtime\_start}.
  \item Formats the current date and time into \procname{timer\_date} and writes a two-line header to \argname{un} (or to the public variable \procname{timer\_unit}, initially \texttt{output\_unit}, if \argname{un} is omitted):
\begin{verbatim}
#TIMERINIT:: Date: YYYY-MM-DD    HH:MM:SS   ms:  MMM
#TIMERINIT::   i MESSG      cpu_tot   cpu_del   wtime_tot   wtime_del   Date
\end{verbatim}
\end{enumerate}

\manheading{Dummy Arguments}
\begin{description}[leftmargin=2.2cm, style=nextline]
  \argitem{un}{integer, scalar, optional, intent(in)}{Fortran I/O unit number for the initialization banner. Defaults to \procname{timer\_unit} (initially \texttt{output\_unit}, i.e., \texttt{stdout}).}
\end{description}

\manheading{Example}
\lstset{language=[95]Fortran}
\begin{lstlisting}
use array_mod
call timer_mod_init(un=f_mout)
\end{lstlisting}

% ===================================================================
% 48. timer_mod_now
% ===================================================================
\manentry{Subroutine}{timer\_mod\_now}{timer\textunderscore mod\textunderscore now}{Updates the public CPU time, wall-clock time, counter, and timestamp variables without printing}{proc:timer_mod_now}
\addcontentsline{toc}{subsubsection}{\texorpdfstring{\texttt{\textbf{timer\_count, timer\_cputot, timer\_cpudel, timer\_wtimetot, timer\_wtimedel, timer\_date}} \textnormal{\footnotesize--- Public timing variables}}{timer\textunderscore count, timer\textunderscore cputot, timer\textunderscore cpudel, timer\textunderscore wtimetot, timer\textunderscore wtimedel, timer\textunderscore date}}

\manheading{Synopsis}
\lstset{language=[95]Fortran}
\begin{lstlisting}
call timer_mod_now
\end{lstlisting}

\manheading{Brief}
Increments \procname{timer\_count} and updates \procname{timer\_cputot}, \procname{timer\_cpudel}, \procname{timer\_wtimetot}, \procname{timer\_wtimedel}, and \procname{timer\_date} relative to \procname{timer\_mod\_init} and the previous checkpoint.

\manheading{Description}
The subroutine \procname{timer\_mod\_now} samples the current CPU clock (\texttt{cpu\_time}) and system wall clock (\texttt{system\_clock}) and updates six public module variables exported by \modname{timer\_mod} and \modname{array\_mod}:
\begin{description}[leftmargin=3.2cm, style=nextline]
  \argitem{timer\_count}{integer, public variable}{Number of times \procname{timer\_mod\_now} (or \procname{timer\_mod\_print}) has been called since \procname{timer\_mod\_init}.}
  \argitem{timer\_cputot}{real(real64), public variable}{Total elapsed CPU time in seconds since the last call to \procname{timer\_mod\_init}.}
  \argitem{timer\_cpudel}{real(real64), public variable}{Incremental CPU time in seconds elapsed since the previous call to \procname{timer\_mod\_now} or \procname{timer\_mod\_print} (or \procname{timer\_mod\_init} on the first call).}
  \argitem{timer\_wtimetot}{real(real64), public variable}{Total elapsed wall-clock time in seconds since the last call to \procname{timer\_mod\_init}.}
  \argitem{timer\_wtimedel}{real(real64), public variable}{Incremental wall-clock time in seconds elapsed since the previous call to \procname{timer\_mod\_now} or \procname{timer\_mod\_print}.}
  \argitem{timer\_date}{character(100), public variable}{Timestamp string at the most recent call to \procname{timer\_mod\_init}, \procname{timer\_mod\_now}, or \procname{timer\_mod\_print}, formatted as \texttt{'Date: YYYY-MM-DD    HH:MM:SS   ms:  MMM'}.}
  \argitem{timer\_unit}{integer, public variable}{Default output unit of \procname{timer\_mod\_init} and \procname{timer\_mod\_print} when \argname{un} is omitted. Initialized to \texttt{output\_unit} (\texttt{stdout}); \modname{timer\_mod} does not depend on \modname{matrix\_mod\_common}, so it does not follow \procname{f\_mout} automatically. Set \lcodef{timer_unit = f_mout} (or any open unit) once to redirect all timer output.}
\end{description}

\manheading{Dummy Arguments}
None. (Updates the public module variables listed above in-place.)

\manheading{Example}
\lstset{language=[95]Fortran}
\begin{lstlisting}
use array_mod
complex(8) :: A(500,500), B(500,500)

call timer_mod_init(un=f_mout)
call random_number(A, sigma=1.0_8)
B = lmatmul(A)

call timer_mod_now
print *, 'timer:', timer_count, timer_cputot, timer_cpudel, &
                   timer_wtimetot, timer_wtimedel, trim(timer_date)
\end{lstlisting}

% ===================================================================
% 49. timer_mod_print
% ===================================================================
\manentry{Subroutine}{timer\_mod\_print}{timer\textunderscore mod\textunderscore print}{Updates all timing variables via \texttt{timer\_mod\_now} and prints a formatted timing report line}{proc:timer_mod_print}

\manheading{Synopsis}
\lstset{language=[95]Fortran}
\begin{lstlisting}
call timer_mod_print([un], [messg])
\end{lstlisting}

\manheading{Brief}
Invokes \procname{timer\_mod\_now} and writes a single formatted \texttt{\#TIMERMOD} checkpoint line containing the call counter, an optional user message \argname{messg}, cumulative and incremental CPU and wall-clock times, and the current date/time.

\manheading{Description}
The subroutine \procname{timer\_mod\_print} calls \procname{timer\_mod\_now} and outputs a structured log line with format \texttt{'(A,I6,A,4G28.16,A)'} to unit \argname{un}:
\begin{verbatim}
#TIMERMOD 1 Flag this point <cpu_tot> <cpu_del> <wtime_tot> <wtime_del> Date:...
\end{verbatim}
Because every output line starts with the tag \texttt{\#TIMERMOD}, timing records can be easily filtered from simulation log files using \texttt{grep '\#TIMERMOD'}.

\manheading{Dummy Arguments}
\begin{description}[leftmargin=2.2cm, style=nextline]
  \argitem{un}{integer, scalar, optional, intent(in)}{Fortran I/O unit number for output. Defaults to \procname{timer\_unit} (initially \texttt{output\_unit}, i.e., \texttt{stdout}).}
  \argitem{messg}{character(*), optional, intent(in)}{Optional descriptive label printed next to the checkpoint counter \procname{timer\_count}. Defaults to \texttt{' '}.}
\end{description}

\manheading{Example}
\lstset{language=[95]Fortran}
\begin{lstlisting}
use array_mod
complex(8) :: A(500,500), B(500,500)

call timer_mod_init(un=f_mout)
call random_number(A, sigma=1.0_8)
call timer_mod_print(un=f_mout, messg='Random generation complete')

B = lmatmul(A)
call timer_mod_print(un=f_mout, messg='End of program')
\end{lstlisting}


% Section 3: File Structure & Section 4: Performance
% ===================================================================
% Section 3: File Structure
% ===================================================================
\section{File Structure}
\label{sec:file_structure}

The top-level module \modname{array\_mod} is defined in the file \filename{array\_mod.f90}. It serves as a clean public facade over the internal implementation modules \modname{matrix\_mod\_array}, \modname{matrix\_mod\_common}, \modname{tensorprod\_mod}, and \modname{timer\_mod}, using \lcodef{private} by default and explicitly exporting only the documented public generic interfaces and variables.

The source files in \filename{Software/Tools/Matrix/} that comprise \modname{array\_mod} and \filename{libMatrix.a} are organized as follows:
\lstset{language=[95]Fortran}
\begin{lstlisting}
array_mod.f90                ! Top-level module exporting public interfaces
matrix_mod_array.f90         ! Generic interfaces and core array procedures
matrix_mod_array_lapack.f90  ! Included in matrix_mod_array.f90: BLAS/LAPACK/PFAPACK routines
matrix_mod_common.f90        ! Common constants (dp, PI, ONE, ZERO, IMU), f_mout, NaN, isNaN
tensorprod_mod.f90           ! Kronecker products (tensorprod, vec2col, vec2row, vectorize)
timer_mod.f90                ! CPU and wall-clock timing utilities (timer_mod_init, etc.)
pfapack/fortran/             ! PFAPACK library source code (f95_pfapack, skpfa, skpf10)
\end{lstlisting}

If desired, one can \lcodef{use matrix_mod_array} directly, in which case all internal module procedures (such as \texttt{array2\_zgeev}, \texttt{array2\_pfaffian}, and the raw \texttt{BLAS}/\texttt{LAPACK} interface blocks) are also public. Similarly, \lcodef{use f95_pfapack} provides direct access to all \texttt{PFAPACK} routines.

The library is built using the provided \filename{Makefile}:
\lstset{language=bash}
\begin{lstlisting}
make gnu   ! Compiles PFAPACK, builds *.mod and *.o with gfortran, creates libMatrix.a in ./ and lib/gnu/
make intel ! Compiles PFAPACK, builds *.mod and *.o with ifx/MKL, creates libMatrix.a in ./ and lib/intel/
make mods  ! Compiles the Fortran modules and produces object files
make lib   ! Archives object files into libMatrix.a
make FC=ifort FFLAGS='-O3 -heap-arrays -xHOST -parallel' LIBS='-mkl=parallel' lib
\end{lstlisting}

The remaining source files in the directory (\filename{matrix\_mod\_matrix*.f90}) implement the object-oriented derived-type \modname{matrix\_mod\_matrix} module, which is intended primarily for educational use rather than high-performance computing.

% ===================================================================
% Section 4: Performance
% ===================================================================
\section{Performance}
\label{sec:performance}

A comprehensive series of benchmarks for matrix--matrix and matrix--vector multiplication was performed on a desktop workstation\footnote{Desktop configuration: Intel(R) Core(TM) i7-4790 CPU @ 3.60\,GHz with Intel \texttt{ifort} Version 19.1.1.217 Build 20200306, compiled with \texttt{-O3 -heap-arrays -xHost -parallel -mkl=parallel}, \texttt{nthreads=8} (4 physical cores).} and on the ARIS national supercomputer.\footnote{Supercomputer configuration (ARIS): Intel Xeon E5-2680v2, Intel \texttt{ifort} Version 19.0.1.144 Build 20181018, compiled with \texttt{-O3 -heap-arrays -xCORE-AVX-I -parallel -mkl=parallel}, \texttt{nthreads=20}.}

\begin{notebox}[Summary of Performance Findings]
\begin{itemize}
  \item \textbf{Matrix--Matrix Multiplication}: Calling \texttt{BLAS} \texttt{zgemm} via \modname{array\_mod} is always faster than the Fortran intrinsic \texttt{MATMUL}, especially for small and medium matrices ($N \le 250$).
  \begin{itemize}
    \item The subroutine \procname{mmmult} is just as fast as direct Fortran~77 calls to \texttt{zgemm} on statically allocated arrays.
    \item The function \procname{lmatmul(m1, m2)} incurs a negligible overhead compared to \procname{mmmult}.
    \item The overloaded operator \procname{.mm.} has a noticeable temporary-array overhead for $N \le 250$ and a small overhead for larger $N$.
    \item Computing $A A^\dagger$ or $A^\dagger A$ via \procname{lmatmul(m1)} or \procname{mmmult(m1, m4)} (which calls \texttt{zherk}) yields a substantial $\sim 35\%$--$40\%$ speedup for large matrices ($N \ge 1000$). By contrast, using \texttt{zhemm} (\texttt{mtype='ZH'}) instead of \texttt{zgemm} for general $A B$ with Hermitian $A$ shows no significant gain.
    \item Power-of-two matrix sizes ($N = 32, 64, 128, 256, 512, 1024$) were also benchmarked against nearby non-power-of-two sizes ($N = 30, 60, 120, 250, 500, 1000$) to test for cache-associativity conflicts with Intel MKL (Table~\ref{tab:b4}); only minor variations were observed.
  \end{itemize}
  \item \textbf{Matrix--Vector Multiplication}: Avoid the intrinsic \texttt{MATMUL} for $N \le 200$ (where \texttt{BLAS} is up to $2\times$--$3\times$ faster); for larger $N$, intrinsic \texttt{MATMUL} pays only a small penalty. Between \procname{lmatmul}, \procname{mvmult}, and \procname{vmmult}, performance on the desktop is nearly identical, whereas on the 20-thread Xeon supercomputer node (Table~\ref{tab:b7}), left vector--matrix multiplication \procname{vmmult} (\texttt{zgemv('T',...)}) is up to $2\times$--$3\times$ faster than \procname{mvmult} (\texttt{zgemv('N',...)}) due to contiguous column-major dot-product access and more efficient multi-threaded vectorization.
\end{itemize}
\end{notebox}

% -------------------------------------------------------------------
% Table B.2: Legend for Matrix-Matrix Benchmarks
% -------------------------------------------------------------------
\begin{table}[htbp]
\centering
\begin{tabular}{c l}
\toprule
\textbf{Column} & \textbf{Matrix--Matrix Multiplication Method} \\
\midrule
\texttt{(5)}  & Direct calls to \texttt{zgemm} for statically allocated arrays \\
\texttt{(11)} & Calls to the subroutine \lcodef{call mmmult(m1, m2, m3)} \\
\texttt{(4)}  & Calls to the function \lcodef{m3 = lmatmul(m1, m2)} \\
\texttt{(3)}  & Calls to the Fortran intrinsic \lcodef{m3 = MATMUL(m1, m2)} \\
\texttt{(6)}  & Operator overloading: \lcodef{m3 = (m1 .mm. m2)} \\
\texttt{(9)}  & Computation of $m_1^\dagger m_1$ via \lcodef{lmatmul(m1, side='H')} (\texttt{zherk}) \\
\bottomrule
\end{tabular}
\caption{\label{tab:b2}Legend of the methods used for benchmarking $N\times N$ complex matrix--matrix multiplication in Tables~\ref{tab:b1}, \ref{tab:b3}, and \ref{tab:b4}.}
\end{table}

% -------------------------------------------------------------------
% Table B.1: Desktop Matrix-Matrix Benchmarks
% -------------------------------------------------------------------
\begin{table}[htbp]
\centering
\begin{tabular}{r c c c c c c}
\toprule
$N$ & \texttt{(5) zgemm} & \texttt{(11) mmmult} & \texttt{(4) lmatmul} & \texttt{(3) MATMUL} & \texttt{(6) .mm.} & \texttt{(9) zherk} \\
\midrule
30   & \texttt{3.09e-06} & \texttt{3.46e-06} & \texttt{3.52e-06} & \texttt{6.56e-06} & \texttt{4.40e-06} & \texttt{4.34e-06} \\
60   & \texttt{1.30e-05} & \texttt{1.36e-05} & \texttt{1.37e-05} & \texttt{2.48e-05} & \texttt{1.74e-05} & \texttt{1.43e-05} \\
120  & \texttt{7.71e-05} & \texttt{8.02e-05} & \texttt{8.00e-05} & \texttt{1.57e-04} & \texttt{9.66e-05} & \texttt{7.29e-05} \\
190  & \texttt{2.86e-04} & \texttt{2.91e-04} & \texttt{3.04e-04} & \texttt{5.47e-04} & \texttt{3.26e-04} & \texttt{2.44e-04} \\
250  & \texttt{6.47e-04} & \texttt{6.58e-04} & \texttt{6.69e-04} & \texttt{7.13e-04} & \texttt{7.20e-04} & \texttt{5.11e-04} \\
500  & \texttt{5.33e-03} & \texttt{5.33e-03} & \texttt{5.35e-03} & \texttt{5.54e-03} & \texttt{5.91e-03} & \texttt{3.71e-03} \\
750  & ---               & \texttt{1.72e-02} & \texttt{1.73e-02} & \texttt{1.82e-02} & \texttt{1.92e-02} & \texttt{1.18e-02} \\
1000 & \texttt{4.12e-02} & \texttt{4.05e-02} & \texttt{4.08e-02} & \texttt{4.27e-02} & \texttt{4.23e-02} & \texttt{2.63e-02} \\
2000 & \texttt{3.08e-01} & \texttt{3.07e-01} & \texttt{3.38e-01} & \texttt{3.49e-01} & \texttt{3.15e-01} & \texttt{1.94e-01} \\
4000 & \texttt{2.51e+00} & \texttt{2.44e+00} & \texttt{2.47e+00} & \texttt{2.41e+00} & \texttt{2.51e+00} & \texttt{1.48e+00} \\
\bottomrule
\end{tabular}
\caption{\label{tab:b1}Time per complex matrix--matrix multiplication in seconds on the desktop workstation (Intel Core i7-4790, \texttt{array\_benchmark\_v01.dat}). Column definitions are given in Table~\ref{tab:b2}.}
\end{table}

% -------------------------------------------------------------------
% Table B.3: Supercomputer Matrix-Matrix Benchmarks
% -------------------------------------------------------------------
\begin{table}[htbp]
\centering
\begin{tabular}{r c c c c c c}
\toprule
$N$ & \texttt{(5) zgemm} & \texttt{(11) mmmult} & \texttt{(4) lmatmul} & \texttt{(3) MATMUL} & \texttt{(6) .mm.} & \texttt{(9) zherk} \\
\midrule
30   & \texttt{1.37e-05} & \texttt{1.47e-05} & \texttt{4.75e-05} & \texttt{3.25e-05} & \texttt{1.58e-05} & \texttt{3.02e-05} \\
60   & \texttt{2.35e-05} & \texttt{2.26e-05} & \texttt{2.29e-05} & \texttt{2.57e-05} & \texttt{3.27e-05} & \texttt{1.18e-04} \\
120  & \texttt{6.72e-05} & \texttt{6.48e-05} & \texttt{6.72e-05} & \texttt{1.10e-04} & \texttt{9.93e-05} & \texttt{6.78e-04} \\
190  & \texttt{2.29e-04} & \texttt{2.30e-04} & \texttt{2.29e-04} & \texttt{3.41e-04} & \texttt{3.22e-04} & \texttt{3.29e-04} \\
250  & \texttt{4.36e-04} & \texttt{4.37e-04} & \texttt{4.36e-04} & \texttt{5.47e-04} & \texttt{5.85e-04} & \texttt{6.67e-04} \\
500  & \texttt{2.69e-03} & \texttt{2.69e-03} & \texttt{2.69e-03} & \texttt{3.18e-03} & \texttt{3.36e-03} & \texttt{2.87e-03} \\
750  & ---               & \texttt{9.48e-03} & \texttt{9.38e-03} & \texttt{1.05e-02} & \texttt{1.13e-02} & \texttt{8.69e-03} \\
1000 & \texttt{2.21e-02} & \texttt{2.29e-02} & \texttt{2.17e-02} & \texttt{2.42e-02} & \texttt{2.56e-02} & \texttt{1.73e-02} \\
2000 & \texttt{1.71e-01} & \texttt{1.71e-01} & \texttt{1.69e-01} & \texttt{1.79e-01} & \texttt{1.84e-01} & \texttt{1.11e-01} \\
4000 & \texttt{1.28e+00} & \texttt{1.28e+00} & \texttt{1.28e+00} & \texttt{1.30e+00} & \texttt{1.34e+00} & \texttt{7.55e-01} \\
\bottomrule
\end{tabular}
\caption{\label{tab:b3}Time per complex matrix--matrix multiplication in seconds on the ARIS supercomputer (20 threads, Intel Xeon E5-2680v2, \texttt{array\_benchmark\_v01\_aris.dat}).}
\end{table}

% -------------------------------------------------------------------
% Table B.4: Power-of-2 Matrix-Matrix Benchmarks on Supercomputer
% -------------------------------------------------------------------
\begin{table}[htbp]
\centering
\begin{tabular}{r c c c c c c}
\toprule
$N$ & \texttt{(5) zgemm} & \texttt{(11) mmmult} & \texttt{(4) lmatmul} & \texttt{(3) MATMUL} & \texttt{(6) .mm.} & \texttt{(9) zherk} \\
\midrule
30   & \texttt{1.32e-05} & \texttt{1.26e-05} & \texttt{1.35e-05} & \texttt{1.39e-05} & \texttt{1.54e-05} & \texttt{2.97e-05} \\
32   & ---               & \texttt{1.28e-05} & \texttt{1.26e-05} & \texttt{1.34e-05} & \texttt{1.56e-05} & \texttt{3.14e-05} \\
60   & \texttt{2.13e-05} & \texttt{2.16e-05} & \texttt{1.92e-05} & \texttt{2.55e-05} & \texttt{3.04e-05} & \texttt{1.18e-04} \\
64   & ---               & \texttt{2.19e-05} & \texttt{2.08e-05} & \texttt{3.34e-05} & \texttt{3.40e-05} & \texttt{1.34e-04} \\
120  & \texttt{6.27e-05} & \texttt{6.38e-05} & \texttt{5.88e-05} & \texttt{1.08e-04} & \texttt{9.69e-05} & \texttt{6.66e-04} \\
128  & ---               & \texttt{7.65e-05} & \texttt{4.50e-03} & \texttt{1.36e-04} & \texttt{1.23e-04} & \texttt{8.02e-04} \\
250  & \texttt{4.29e-04} & \texttt{4.27e-04} & \texttt{4.37e-04} & \texttt{5.42e-04} & \texttt{5.81e-04} & \texttt{6.68e-04} \\
256  & ---               & \texttt{4.85e-04} & \texttt{4.96e-04} & \texttt{6.71e-04} & \texttt{6.38e-04} & \texttt{6.48e-04} \\
500  & \texttt{2.67e-03} & \texttt{2.71e-03} & \texttt{2.67e-03} & \texttt{3.07e-03} & \texttt{3.35e-03} & \texttt{2.75e-03} \\
512  & ---               & \texttt{3.08e-03} & \texttt{3.11e-03} & \texttt{3.55e-03} & \texttt{3.78e-03} & \texttt{3.01e-03} \\
1000 & \texttt{2.18e-02} & \texttt{2.22e-02} & \texttt{2.22e-02} & \texttt{2.38e-02} & \texttt{2.51e-02} & \texttt{1.71e-02} \\
1024 & ---               & \texttt{2.40e-02} & \texttt{2.38e-02} & \texttt{2.58e-02} & \texttt{2.76e-02} & \texttt{1.90e-02} \\
\bottomrule
\end{tabular}
\caption{\label{tab:b4}Comparison of matrix--matrix multiplication times (in seconds) on the ARIS supercomputer for matrix sizes $N$ that are powers of 2 versus nearby non-powers of 2 (\texttt{array\_benchmark\_v02*.dat}).}
\end{table}

% -------------------------------------------------------------------
% Table B.6: Legend for Matrix-Vector Benchmarks
% -------------------------------------------------------------------
\begin{table}[htbp]
\centering
\begin{tabular}{c l}
\toprule
\textbf{Column} & \textbf{Matrix--Vector Multiplication Method} \\
\midrule
\texttt{(3)} & Call to the Fortran intrinsic \lcodef{v3 = MATMUL(m1, v2)} \\
\texttt{(4)} & Call to the function \lcodef{v3 = lmatmul(m1, v2)} \\
\texttt{(5)} & Call to the subroutine \lcodef{call mvmult(m1, v2, v3)} \\
\texttt{(6)} & Call to the subroutine \lcodef{call vmmult(v1, m1, v3)} \\
\bottomrule
\end{tabular}
\caption{\label{tab:b6}Legend of the methods used for benchmarking complex matrix--vector multiplication in Table~\ref{tab:b5} (sub-tables B.5 and B.7).}
\end{table}

% -------------------------------------------------------------------
% Table B.5 & B.7: Desktop and Supercomputer Matrix-Vector Benchmarks
% -------------------------------------------------------------------
\begin{table}[htbp]
\centering
\scriptsize
\setlength{\tabcolsep}{4pt}
\begin{tabular}{r c c c c @{\hspace{1.2em}} c c c c}
\toprule
 & \multicolumn{4}{c}{\textbf{Table B.5: Desktop (Core i7-4790)}} & \multicolumn{4}{c}{\textbf{Table B.7: Supercomputer (Xeon E5-2680v2)}} \\
\cmidrule(lr){2-5} \cmidrule(lr){6-9}
$N$ & \texttt{(3) MATMUL} & \texttt{(4) lmatmul} & \texttt{(5) mvmult} & \texttt{(6) vmmult} & \texttt{(3) MATMUL} & \texttt{(4) lmatmul} & \texttt{(5) mvmult} & \texttt{(6) vmmult} \\
\midrule
30    & \texttt{8.15e-07} & \texttt{7.41e-07} & \texttt{6.71e-07} & \texttt{6.66e-07} & \texttt{8.72e-07} & \texttt{1.08e-06} & \texttt{9.20e-07} & \texttt{9.00e-07} \\
60    & \texttt{2.68e-06} & \texttt{2.06e-06} & \texttt{1.95e-06} & \texttt{1.84e-06} & \texttt{3.26e-06} & \texttt{3.33e-06} & \texttt{3.19e-06} & \texttt{3.00e-06} \\
120   & \texttt{9.65e-06} & \texttt{5.18e-06} & \texttt{4.87e-06} & \texttt{3.34e-06} & \texttt{1.27e-05} & \texttt{1.09e-05} & \texttt{9.67e-06} & \texttt{8.09e-06} \\
190   & \texttt{2.28e-05} & \texttt{1.08e-05} & \texttt{1.04e-05} & \texttt{6.36e-06} & \texttt{3.25e-05} & \texttt{1.22e-05} & \texttt{1.09e-05} & \texttt{7.86e-06} \\
250   & \texttt{1.82e-05} & \texttt{1.67e-05} & \texttt{1.64e-05} & \texttt{9.16e-06} & \texttt{3.38e-05} & \texttt{2.88e-05} & \texttt{2.94e-05} & \texttt{1.03e-05} \\
500   & \texttt{4.92e-05} & \texttt{4.21e-05} & \texttt{4.11e-05} & \texttt{2.91e-05} & \texttt{5.64e-05} & \texttt{6.35e-05} & \texttt{6.10e-05} & \texttt{1.69e-05} \\
750   & \texttt{1.63e-04} & \texttt{1.46e-04} & \texttt{1.44e-04} & \texttt{1.34e-04} & \texttt{1.27e-04} & \texttt{1.18e-04} & \texttt{1.13e-04} & \texttt{3.30e-05} \\
1000  & \texttt{5.91e-04} & \texttt{6.09e-04} & \texttt{6.97e-04} & \texttt{5.12e-04} & \texttt{1.29e-04} & \texttt{1.50e-04} & \texttt{1.61e-04} & \texttt{4.63e-05} \\
2000  & \texttt{3.13e-03} & \texttt{2.87e-03} & \texttt{2.90e-03} & \texttt{2.84e-03} & \texttt{9.80e-04} & \texttt{1.25e-03} & \texttt{1.25e-03} & \texttt{4.16e-04} \\
4000  & \texttt{1.19e-02} & \texttt{1.16e-02} & \texttt{1.14e-02} & \texttt{1.14e-02} & \texttt{5.58e-03} & \texttt{5.66e-03} & \texttt{5.60e-03} & \texttt{2.55e-03} \\
8000  & \texttt{4.79e-02} & \texttt{4.60e-02} & \texttt{4.66e-02} & \texttt{4.51e-02} & \texttt{2.27e-02} & \texttt{2.25e-02} & \texttt{2.27e-02} & \texttt{1.11e-02} \\
16000 & \texttt{1.96e-01} & \texttt{1.89e-01} & \texttt{1.88e-01} & \texttt{1.80e-01} & \texttt{8.43e-02} & \texttt{8.33e-02} & \texttt{8.46e-02} & \texttt{4.21e-02} \\
\bottomrule
\end{tabular}
\caption{\label{tab:b5}\label{tab:b7}Time per complex matrix--vector multiplication in seconds on the desktop workstation (left, Table~B.5) and on the ARIS supercomputer (right, Table~B.7). Column definitions are given in Table~\ref{tab:b6}.}
\end{table}

% ===================================================================
% References
% ===================================================================
\begin{thebibliography}{9}

\bibitem{Wimmer_2012}
M.~Wimmer,
\emph{Algorithm 923: Efficient Numerical Computation of the Pfaffian for Dense and Banded Skew-Symmetric Matrices},
ACM Transactions on Mathematical Software (TOMS) \textbf{38} (2012) 30,
\href{https://arxiv.org/abs/1102.3440}{arXiv:1102.3440 [cond-mat.mes-hall]}.

\bibitem{pfapack}
M.~Wimmer,
\emph{PFAPACK: A library for computing the Pfaffian of skew-symmetric matrices},
\url{https://www.michaelwimmer.org/downloads.html}.

\bibitem{Catterall_2003}
S.~Catterall,
\emph{ Lattice Supersymmetry and Topological Field Theory},
JHEP \textbf{05} (2003) 038 / Phys.\ Rev.\ D \textbf{68}, 014503 (2003),
\href{https://doi.org/10.1103/PhysRevD.68.014503}{doi:10.1103/PhysRevD.68.014503}.

\end{thebibliography}


\end{document}
